%%%%%%%%%%%%%%%%%%%%%%%%%%%%%%%%%%%%%%%%%%%%%%%%%%%%%%%%%%%%%%%%%
\chapter{Unit 14: Order of Differential Equations}
\label{chap:unit14}
%%%%%%%%%%%%%%%%%%%%%%%%%%%%%%%%%%%%%%%%%%%%%%%%%%%%%%%%%%%%%%%%%

%%===============================================================
\section{Order of Differential Equations}
\label{sec:order-of-differential-equations}
%%===============================================================

Introduction



% [Image: Visualization of heat transfer in a pump casing, created by solving

the heat equation. Heat is being generated internally in the casing and

being cooled at the boundary, providing a steady state temperature

distribution.]{width=""350""}

In mathematics, a \textbf{differential equation} is an equation that relates

one or more functions and their derivatives. In applications, the

functions generally represent physical quantities, the derivatives

represent their rates of change, and the differential equation defines a

relationship between the two. Such relations are common; therefore,

differential equations play a prominent role in many disciplines

including engineering, physics, economics, and biology.



Mainly the study of differential equations consists of the study of

their solutions (the set of functions that satisfy each equation), and

of the properties of their solutions. Only the simplest differential

equations are solvable by explicit formulas; however, many properties of

solutions of a given differential equation may be determined without

computing them exactly.



Often when a closed-form expression for the solutions is not available,

solutions may be approximated numerically using computers. The theory of

dynamical systems puts emphasis on qualitative analysis of systems

described by differential equations, while many numerical methods have

been developed to determine solutions with a given degree of accuracy.

The order of a differential equation is determined by the highest-order

derivative. The higher the order of the differential equation, the more

arbitrary constants need to be added to the general solution. A

first-order equation will have one, a second-order two, and so on. The

degree of a differential equation, similarly, is determined by the

highest exponent on any variables involved.



Examples:



\begin{itemize}
    \item $y'' + xy' - x^3y = \sin{x}$ is second order (or \""of order 2\"").
    \item $x't + x = t^2$ is first order.
    \item $y''' + 3y'' + 3y' + y = x^2$ is of order 3.
\end{itemize}

\paragraph{Types}



Differential equations can be divided into several types. Apart from

describing the properties of the equation itself, these classes of

differential equations can help inform the choice of approach to a

solution. Commonly used distinctions include whether the equation is

ordinary or partial, linear or non-linear, and homogeneous or

heterogeneous. This list is far from exhaustive; there are many other

properties and subclasses of differential equations which can be very

useful in specific contexts.



\paragraph{Ordinary differential equations}



An ordinary differential equation (\textit{ODE}) is an equation containing an

unknown function of one real or complex variable

`{{mvar|x}}`{=mediawiki}, its derivatives, and some given functions of

`{{mvar|x}}`{=mediawiki}. The unknown function is generally represented

by a variable (often denoted `{{mvar|y}}`{=mediawiki}), which,

therefore, \textit{depends} on `{{mvar|x}}`{=mediawiki}. Thus

`{{mvar|x}}`{=mediawiki} is often called the independent variable of the

equation. The term \""\textit{ordinary}\"" is used in contrast with the term

partial differential equation, which may be with respect to \textit{more than}

one independent variable.



Linear differential equations are the differential equations that are

linear in the unknown function and its derivatives. Their theory is well

developed, and in many cases one may express their solutions in terms of

integrals.



Most ODEs that are encountered in physics are linear. Therefore, most

special functions may be defined as solutions of linear differential

equations (see Holonomic function).



As, in general, the solutions of a differential equation cannot be

expressed by a closed-form expression, numerical methods are commonly

used for solving differential equations on a computer.



\paragraph{Partial differential equations}



A partial differential equation (\textit{PDE}) is a differential equation that

contains unknown multivariable functions and their partial derivatives.

(This is in contrast to ordinary differential equations, which deal with

functions of a single variable and their derivatives.) PDEs are used to

formulate problems involving functions of several variables, and are

either solved in closed form, or used to create a relevant computer

model.



PDEs can be used to describe a wide variety of phenomena in nature such

as sound, heat, electrostatics, electrodynamics, fluid flow, elasticity,

or quantum mechanics. These seemingly distinct physical phenomena can be

formalized similarly in terms of PDEs. Just as ordinary differential

equations often model one-dimensional dynamical systems, partial

differential equations often model multidimensional systems. Stochastic

partial differential equations generalize partial differential equations

for modeling randomness.



\paragraph{Non-linear differential equations}



A \textbf{non-linear differential equation} is a differential equation that

is not a linear equation in the unknown function and its derivatives

(the linearity or non-linearity in the arguments of the function are not

considered here). There are very few methods of solving nonlinear

differential equations exactly; those that are known typically depend on

the equation having particular symmetries. Nonlinear differential

equations can exhibit very complicated behavior over extended time

intervals, characteristic of chaos. Even the fundamental questions of

existence, uniqueness, and extendability of solutions for nonlinear

differential equations, and well-posedness of initial and boundary value

problems for nonlinear PDEs are hard problems and their resolution in

special cases is considered to be a significant advance in the

mathematical theory (cf. Navier--Stokes existence and smoothness).

However, if the differential equation is a correctly formulated

representation of a meaningful physical process, then one expects it to

have a solution.



Linear differential equations frequently appear as approximations to

nonlinear equations. These approximations are only valid under

restricted conditions. For example, the harmonic oscillator equation is

an approximation to the nonlinear pendulum equation that is valid for

small amplitude oscillations (see below).



\paragraph{Equation order}



Differential equations are described by their order, determined by the

term with the highest derivatives. An equation containing only first

derivatives is a \textit{first-order differential equation}, an equation

containing the second derivative is a *second-order differential

equation*, and so on. Differential equations that describe natural

phenomena almost always have only first and second order derivatives in

them, but there are some exceptions, such as the thin film equation,

which is a fourth order partial differential equation.



\paragraph{Examples}



In the first group of examples \textit{u} is an unknown function of \textit{x}, and

\textit{c} and \textit{?} are constants that are supposed to be known. Two broad

classifications of both ordinary and partial differential equations

consist of distinguishing between \textit{linear} and \textit{nonlinear} differential

equations, and between \textit{homogeneous} differential equations and

\textit{heterogeneous} ones.



\begin{itemize}
    \item Heterogeneous first-order linear constant coefficient ordinary
\end{itemize}
    differential equation:



:   



    :   $\frac{du}{dx} = cu+x^2.$



\begin{itemize}
    \item Homogeneous second-order linear ordinary differential equation:
\end{itemize}

:   



    :   $\frac{d^2u}{dx^2} - x\frac{du}{dx} + u = 0.$



\begin{itemize}
    \item Homogeneous second-order linear constant coefficient ordinary
\end{itemize}
    differential equation describing the harmonic oscillator:



:   



    :   $\frac{d^2u}{dx^2} + \omega^2u = 0.$



\begin{itemize}
    \item Heterogeneous first-order nonlinear ordinary differential equation:
\end{itemize}

:   



    :   $\frac{du}{dx} = u^2 + 4.$



\begin{itemize}
    \item Second-order nonlinear (due to sine function) ordinary differential
\end{itemize}
    equation describing the motion of a pendulum of length \textit{L}:



:   



    :   $L\frac{d^2u}{dx^2} + g\sin u = 0.$



In the next group of examples, the unknown function \textit{u} depends on two

variables \textit{x} and \textit{t} or \textit{x} and \textit{y}.



\begin{itemize}
    \item Homogeneous first-order linear partial differential equation:
\end{itemize}

:   



    :   $\frac{\partial u}{\partial t} + t\frac{\partial u}{\partial x} = 0.$



\begin{itemize}
    \item Homogeneous second-order linear constant coefficient partial
\end{itemize}
    differential equation of elliptic type, the Laplace equation:



:   



    :   $\frac{\partial^2 u}{\partial x^2} + \frac{\partial^2 u}{\partial y^2} = 0.$



\begin{itemize}
    \item Homogeneous third-order non-linear partial differential equation :
\end{itemize}

:   



    :   $\frac{\partial u}{\partial t} = 6u\frac{\partial u}{\partial x} - \frac{\partial^3 u}{\partial x^3}.$



\paragraph{Existence of solutions}



Solving differential equations is not like solving algebraic equations.

Not only are their solutions often unclear, but whether solutions are

unique or exist at all are also notable subjects of interest.



For first order initial value problems, the Peano existence theorem

gives one set of circumstances in which a solution exists. Given any

point $(a,b)$ in the xy-plane, define some rectangular region $Z$, such

that $Z = [l,m]\times[n,p]$ and $(a,b)$ is in the interior of $Z$. If we

are given a differential equation $\frac{dy}{dx} = g(x,y)$ and the

condition that $y=b$ when $x=a$, then there is locally a solution to

this problem if $g(x,y)$ and $\frac{\partial g}{\partial x}$ are both

continuous on $Z$. This solution exists on some interval with its center

at $a$. The solution may not be unique. (See Ordinary differential

equation for other results.)



However, this only helps us with first order initial value problems.

Suppose we had a linear initial value problem of the nth order:



$$f_{n}(x)\frac{d^n y}{dx^n} + \cdots + f_{1}(x)\frac{d y}{dx} + f_{0}(x)y = g(x)$$

such that



$$y(x_{0})=y_{0}, y'(x_{0}) = y'_{0}, y''(x_{0}) = y''_{0}, \ldots$$



For any nonzero $f_{n}(x)$, if $\{f_{0},f_{1},\ldots\}$ and $g$ are

continuous on some interval containing $x_{0}$, $y$ is unique and

exists.



\paragraph{Connection to difference equations}



The theory of differential equations is closely related to the theory of

difference equations, in which the coordinates assume only discrete

values, and the relationship involves values of the unknown function or

functions and values at nearby coordinates. Many methods to compute

numerical solutions of differential equations or study the properties of

differential equations involve the approximation of the solution of a

differential equation by the solution of a corresponding difference

equation.



\paragraph{Applications}



The study of differential equations is a wide field in pure and applied

mathematics, physics, and engineering. All of these disciplines are

concerned with the properties of differential equations of various

types. Pure mathematics focuses on the existence and uniqueness of

solutions, while applied mathematics emphasizes the rigorous

justification of the methods for approximating solutions. Differential

equations play an important role in modeling virtually every physical,

technical, or biological process, from celestial motion, to bridge

design, to interactions between neurons. Differential equations such as

those used to solve real-life problems may not necessarily be directly

solvable, i.e. do not have closed form solutions. Instead, solutions can

be approximated using numerical methods.



Many fundamental laws of physics and chemistry can be formulated as

differential equations. In biology and economics, differential equations

are used to model the behavior of complex systems. The mathematical

theory of differential equations first developed together with the

sciences where the equations had originated and where the results found

application. However, diverse problems, sometimes originating in quite

distinct scientific fields, may give rise to identical differential

equations. Whenever this happens, mathematical theory behind the

equations can be viewed as a unifying principle behind diverse

phenomena. As an example, consider the propagation of light and sound in

the atmosphere, and of waves on the surface of a pond. All of them may

be described by the same second-order partial differential equation, the

wave equation, which allows us to think of light and sound as forms of

waves, much like familiar waves in the water. Conduction of heat, the

theory of which was developed by Joseph Fourier, is governed by another

second-order partial differential equation, the heat equation. It turns

out that many diffusion processes, while seemingly different, are

described by the same equation; the Black--Scholes equation in finance

is, for instance, related to the heat equation.



\paragraph{Resources}



\begin{itemize}
    \item \href{https://courses.lumenlearning.com/boundless-calculus/chapter/differential-equations/}{Differential
\end{itemize}
    Equations},

    Lumen Learning

\begin{itemize}
    \item \href{https://www.analyzemath.com/calculus/Differential_Equations/order_linearity.html}{Order and Linearity of Differential
\end{itemize}
    Equations},

    Analyze Math



\paragraph{Licensing}



Content obtained and/or adapted from:



\begin{itemize}
    \item \href{https://en.wikipedia.org/wiki/Differential_equation}{Differential equation,
\end{itemize}
    Wikipedia}

    under a CC BY-SA license"

%%===============================================================
\section{Solutions of Differential Equations}
\label{sec:solutions-of-differential-equations}
%%===============================================================

"A solution of a differential equation is an expression of the dependent

variable that satisfies the relation established in the differential

equation. For example, the solution of $y' + y - 2 = 4x$ will be some

equation y = f(x) such that y and its first derivative, y\', satisfy the

relation $y' + y - 2 = 4x$. The general solution of a differential

equation will have one or more arbitrary constants, depending on the

order of the original differential equation (the solution of a first

order diff. eq. will have one arbitrary constant, a second order one

will have two, etc.).



Examples:



\begin{itemize}
    \item $y' = 2x$. Through simple integration, we can calculate the general
\end{itemize}
    solution of this equation to be $y = x^2 + C$, where C is an

    arbitrary constant.

\begin{itemize}
    \item $y' - y = 0$. The G.S. is $y = Ce^{x}$. $y' = Ce^{x}$, so
\end{itemize}
    $y' - y = 0 \to Ce^{x} - Ce^{x} = 0$, so this solution satisfies the

    relationship for all arbitrary constants C.

\begin{itemize}
    \item $y'' + y' - 2y = 0$. The G.S. is $y = Ce^{x} + De^{-2x}$.
\end{itemize}
    $y' = Ce^{x} - 2De^{-2x}$ and $y'' = Ce^{x} + 4De^{-2x}$, so

    $0 = y'' + y' - 2y$ becomes

    $0 = Ce^{x} + 4De^{-2x} + Ce^{x} - 2De^{-2x} - 2(Ce^{x} + De^{-2x}) = Ce^{x} + Ce^{x} - 2Ce^{x} + 4De^{-2x} - 2De^{-2x} - 2De^{-2x}) = 0$.



The particular solution of a differential equation can be solved if we

have enough points to solve for the arbitrary constants (see initial

value problems for more).



\paragraph{Existence of solutions}



Solving differential equations is not like solving algebraic equations.

Not only are their solutions often unclear, but whether solutions are

unique or exist at all are also notable subjects of interest.



For first order initial value problems, the Peano existence theorem

gives one set of circumstances in which a solution exists. Given any

point $(a,b)$ in the xy-plane, define some rectangular region $Z$, such

that $Z = [l,m]\times[n,p]$ and $(a,b)$ is in the interior of $Z$. If we

are given a differential equation $\frac{dy}{dx} = g(x,y)$ and the

condition that $y=b$ when $x=a$, then there is locally a solution to

this problem if $g(x,y)$ and $\frac{\partial g}{\partial x}$ are both

continuous on $Z$. This solution exists on some interval with its center

at $a$. The solution may not be unique.



However, this only helps us with first order initial value problems.

Suppose we had a linear initial value problem of the nth order:



$$f_{n}(x)\frac{d^n y}{dx^n} + \cdots + f_{1}(x)\frac{d y}{dx} + f_{0}(x)y = g(x)$$

such that



$$y(x_{0})=y_{0}, y'(x_{0}) = y'_{0}, y''(x_{0}) = y''_{0}, \ldots$$



For any nonzero $f_{n}(x)$, if $\{f_{0},f_{1},\ldots\}$ and $g$ are

continuous on some interval containing $x_{0}$, $y$ is unique and

exists.



\paragraph{Resources}



\begin{itemize}
    \item \href{http://www.maths.gla.ac.uk/\textasciitilde{}cc/2x/2005_2xnotes/2x_chap5.pdf}{Differential
\end{itemize}
    Equations},

    University of Glascow

\begin{itemize}
    \item \href{https://www.api.simply.science/index.php/math/calculus/differential-equations/types-of-differential-equations/9869-general-and-particular-solutions-of-a-differential-equation}{General and Particular
\end{itemize}
    Solutions},

    Simply Math



\paragraph{Licensing}



Content obtained and/or adapted from:



\begin{itemize}
    \item \href{https://en.wikipedia.org/wiki/Differential_equation}{Differential equation,
\end{itemize}
    Wikipedia}

    under a CC BY-SA license"

%%===============================================================
\section{Initial Value Problem (IVP)}
\label{sec:initial-value-problem-ivp}
%%===============================================================

Definition



An \textbf{initial value problem} is a differential equation



$$y'(t) = f(t, y(t))$$ with

$f\colon \Omega \subset \mathbb{R} \times \mathbb{R}^n \to \mathbb{R}^n$

where $\Omega$ is an open set of $\mathbb{R} \times \mathbb{R}^n$,

together with a point in the domain of $f$



$$(t_0, y_0) \in \Omega,$$ called the initial condition.



A \textbf{solution} to an initial value problem is a function $y$ that is a

solution to the differential equation and satisfies



$$y(t_0) = y_0.$$



In higher dimensions, the differential equation is replaced with a

family of equations $y_i'(t)=f_i(t, y_1(t), y_2(t), \dotsc)$, and $y(t)$

is viewed as the vector $(y_1(t), \dotsc, y_n(t))$, most commonly

associated with the position in space. More generally, the unknown

function $y$ can take values on infinite dimensional spaces, such as

Banach spaces or spaces of distributions.



Initial value problems are extended to higher orders by treating the

derivatives in the same way as an independent function, e.g.

$y''(t)=f(t,y(t),y'(t))$.



\paragraph{Examples}



With initial value problems, we are given a differential equation, and

one or more points (depending on the order of the equation) to solve the

constants in the general solution. For a first order differential

equation we need 1 point $(x, y(x))$, for a second order equation we

need 2 points (typically $(x_1, y(x_1))$ and either $(x_2, y(x_2))$ or

$(x_2, y'(x_2))$), and so on.



A simple example is to solve $y'(t) = 0.85 y(t)$ and $y(0) = 19$. We are

trying to find a formula for $y(t)$ that satisfies these two equations.



Rearrange the equation so that $y$ is on the left hand side



:   $\frac{y'(t)}{y(t)} = 0.85$



Now integrate both sides with respect to $t$ (this introduces an unknown

constant $B$).



:   $\int \frac{y'(t)}{y(t)}\,dt = \int 0.85\,dt$

:   $\ln |y(t)| = 0.85t + B$



Eliminate the logarithm with exponentiation on both sides



:   $| y(t) | = e^Be^{0.85t}$



Let $C$ be a new unknown constant, $C = \pm e^B$, so



:   $y(t) = Ce^{0.85t}$



Now we need to find a value for $C$. Use $y(0) = 19$ as given at the

start and substitute 0 for $t$ and 19 for $y$



:   $19 = C e^{0.85 \cdot 0}$

:   $C = 19$



this gives the final solution of $y(t) = 19e^{0.85t}$.



Second example



The solution of



:   $y'+3y=6t+5,\qquad y(0)=3$



can be found to be



:   $y(t)=2e^{-3t}+2t+1. \,$



Indeed,



:   $\begin{aligned}

    y'+3y \&= \tfrac{d}{dt} (2e^{-3t}+2t+1)+3(2e^{-3t}+2t+1) \\\          \&= (-6e^{-3t}+2)+(6e^{-3t}+6t+3) \\\          \&= 6t+5.

    \end{aligned}$



Some more simple examples of initial value problems:



\begin{itemize}
    \item $y' = 2x$, $y(2) = 0$. With this point and the general solution
\end{itemize}
    $y = x^2 + C$, we can calculate the constant C to be -4. Thus the

    particular solution is $y = x^2 - 4$.

\begin{itemize}
    \item $y' - y = 0$, $y(0) = 3$. $3 = Ce^{0} \implies C = 3$, so the
\end{itemize}
    particular solution is $y = 3e^{x}$.

\begin{itemize}
    \item $y'' + y' - 2y = 0$, $y(0) = 2$, $y'(0) = -1$. So,
\end{itemize}
    $2 = Ce^{0} + De^{0} = C + D$ and $-1 = Ce^{0} - 2De^{0} = C - 2D$.

    Thus C = 1 and D = 1, and the particular solution is

    $y = e^{x} + e^{-2x}$.



\paragraph{Licensing}



Content obtained and/or adapted from:



\begin{itemize}
    \item \href{https://en.wikipedia.org/wiki/Initial_value_problem}{Initial value problem,
\end{itemize}
    Wikipedia}

    under a CC BY-SA license"

%%===============================================================
\section{Cauchy Problem}
\label{sec:cauchy-problem}
%%===============================================================

"A \textbf{Cauchy problem} in mathematics asks for the solution of a partial

differential equation that satisfies certain conditions that are given

on a hypersurface in the domain. A Cauchy problem can be an initial

value problem or a boundary value problem. It is named after

Augustin-Louis Cauchy.



\paragraph{Formal statement}



For a partial differential equation defined on \textbf{R}^\textit{n+1}^ and a smooth

manifold \textit{S} ? \textbf{R}^\textit{n+1}^ of dimension \textit{n} (\textit{S} is called the Cauchy

surface), the Cauchy problem consists of finding the unknown functions

$u_1,\dots,u_N$ of the differential equation with respect to the

independent variables $t,x_1,\dots,x_n$ that satisfies



$$\begin{aligned}&\frac{\partial^{n_i}u_i}{\partial t^{n_i}} = F_i\left(t,x_1,\dots,x_n,u_1,\dots,u_N,\dots,\frac{\partial^k u_j}{\partial t^{k_0}\partial x_1^{k_1}\dots\partial x_n^{k_n}},\dots\right) \\\&\text{for } i,j = 1,2,\dots,N;\, k_0+k_1+\dots+k_n=k\leq n_j;\, k_0<n_j

\end{aligned}$$ subject to the condition, for some value $t=t_0$,



$$\frac{\partial^k u_i}{\partial t^k}=\phi_i^{(k)}(x_1,\dots,x_n)

\quad \text{for } k=0,1,2,\dots,n_i-1$$



where $\phi_i^{(k)}(x_1,\dots,x_n)$ are given functions defined on the

surface $S$ (collectively known as the \textbf{Cauchy data} of the problem).

The derivative of order zero means that the function itself is

specified.



\paragraph{Cauchy--Kowalevski theorem}



The Cauchy--Kowalevski theorem states that *If all the functions $F_i$

are analytic in some neighborhood of the point

$(t^0,x_1^0,x_2^0,\dots,\phi_{j,k_0,k_1,\dots,k_n}^0,\dots)$, and if all

the functions $\phi_j^{(k)}$ are analytic in some neighborhood of the

point $(x_1^0,x_2^0,\dots,x_n^0)$, then the Cauchy problem has a unique

analytic solution in some neighborhood of the point

$(t^0,x_1^0,x_2^0,\dots,x_n^0)$*.



\paragraph{Licensing}



Content obtained and/or adapted from:



\begin{itemize}
    \item \href{https://en.wikipedia.org/wiki/Cauchy_problem}{Cauchy problem,
\end{itemize}
    Wikipedia} under a CC

    BY-SA license"

%%===============================================================
\section{Separation of Variables (1st Order)}
\label{sec:separation-of-variables-1st-order}
%%===============================================================

"In mathematics, \textbf{separation of variables} (also known as the **Fourier

method**) is any of several methods for solving ordinary and partial

differential equations, in which algebra allows one to rewrite an

equation so that each of two variables occurs on a different side of the

equation.



% [Image: Solve proportional first order differential equation by separation of

variables.]



% [Image: Solve linear first order differential equation by separation of

variables.]



\paragraph{Ordinary differential equations (ODE)}



Suppose a differential equation can be written in the form



$$\frac{d}{dx} f(x) = g(x)h(f(x))$$



which we can write more simply by letting $y = f(x)$:



$$\frac{dy}{dx}=g(x)h(y).$$



As long as \textit{h}(\textit{y}) ? 0, we can rearrange terms to obtain:



$${dy \over h(y)} = g(x) \, dx,$$



so that the two variables \textit{x} and \textit{y} have been separated. \textit{dx} (and

\textit{dy}) can be viewed, at a simple level, as just a convenient notation,

which provides a handy mnemonic aid for assisting with manipulations. A

formal definition of \textit{dx} as a differential (infinitesimal) is somewhat

advanced.



\paragraph{Alternative notation}



Those who dislike Leibniz\'s notation may prefer to write this as



$$\frac{1}{h(y)} \frac{dy}{dx} = g(x),$$



but that fails to make it quite as obvious why this is called

\""separation of variables\"". Integrating both sides of the equation with

respect to $x$, we have



$$\int \frac{1}{h(y)} \frac{dy}{dx} \, dx = \int g(x) \, dx,$$ \textbf{(A1)}



:   or equivalently,



$$\int \frac{1}{h(y)} \, dy = \int g(x) \, dx$$



because of the substitution rule for integrals.



If one can evaluate the two integrals, one can find a solution to the

differential equation. Observe that this process effectively allows us

to treat the derivative $\frac{dy}{dx}$ as a fraction which can be

separated. This allows us to solve separable differential equations more

conveniently, as demonstrated in the example below.



(Note that we do not need to use two constants of integration, in

equation (A1) as in



$$\int \frac{1}{h(y)} \, dy + C_1 = \int g(x) \, dx + C_2,$$ because a

single constant $C = C_2 - C_1$ is equivalent.)



\paragraph{Example}



Population growth is often modeled by the differential equation



:   $\frac{dP}{dt}=kP\left(1-\frac{P}{K}\right)$



where $P$ is the population with respect to time $t$, $k$ is the rate of

growth, and $K$ is the carrying capacity of the environment.



Separation of variables may be used to solve this differential equation.



:   $\begin{aligned}

    \& \frac{dP}{dt}=kP\left(1-\frac{P}{K}\right) \\\\[5pt]

    \& \int\frac{dP}{P\left(1-\frac{P}{K}\right)}=\int k\,dt

    \end{aligned}$



To evaluate the integral on the left side, we simplify the fraction



$$\begin{aligned} \frac{1}{P\left(1-\frac{P}{K}\right)}=\frac{K}{P\left(K-P\right)}

\hspace{45em} \end{aligned}$$



and then, we decompose the fraction into partial fractions



:   $\frac{K}{P(K-P)}=\frac{1}{P}+\frac{1}{K-P}$



Thus we have



:   $\begin{aligned}

    \& \int\left(\frac{1}{P}+\frac{1}{K-P}\right) dP=\int k\,dt \\\\[6pt]

    \& \ln|P|-\ln|K-P|=kt+C \\\\[6pt]

    \& \ln|K-P|-\ln|P|=-kt-C \\\\[6pt]

    \& \ln\left|\cfrac{K-P}{P}\right|=-kt-C \\\\[6pt]

    \& \left|\dfrac{K-P}{P}\right|=e^{-kt-C} \\\\[6pt]

    \& \left|\dfrac{K-P}{P}\right|=e^{-C}e^{-kt} \\\\[6pt]

    \& \frac{K-P}{P}=\pm e^{-C}e^{-kt}

    \end{aligned}$



Let $A=\pm e^{-C}$.



$$\begin{aligned}

& \frac{K-P}{P}=Ae^{-kt} \\\\[6pt]

& \frac{K}{P}-1=Ae^{-kt} \\\\[6pt]

& \frac{K}{P}=1+Ae^{-kt} \\\\[6pt]

& \frac{P}{K}=\frac{1}{1+Ae^{-kt}} \\\\[6pt]

& P=\frac{K}{1+Ae^{-kt}}

\end{aligned}$$



Therefore, the solution to the logistic equation is



:   $P(t)=\frac{K}{1+Ae^{-kt}}$



To find $A$, let $t=0$ and $P\left(0\right)=P_0$. Then we have



:   $P_0=\frac{K}{1+Ae^0}$



Noting that $e^0=1$, and solving for \textit{A} we get



:   $A=\frac{K-P_0}{P_0}.$



\paragraph{Generalization of separable ODEs to the nth order}



Much like one can speak of a separable first-order ODE, one can speak of

a separable second-order, third-order or \textit{n}th-order ODE. Consider the

separable first-order ODE:



$$\frac{dy}{dx}=f(y)g(x)$$ The derivative can alternatively be written

the following way to underscore that it is an operator working on the

unknown function, \textit{y}:



$$\frac{dy}{dx}=\frac{d}{dx}(y)$$ Thus, when one separates variables for

first-order equations, one in fact moves the \textit{dx} denominator of the

operator to the side with the \textit{x} variable, and the \textit{d}(\textit{y}) is left on

the side with the \textit{y} variable. The second-derivative operator, by

analogy, breaks down as follows:



$$\frac{d^2y}{dx^2} = \frac{d}{dx}\left(\frac{dy}{dx}\right) = \frac{d}{dx}\left(\frac{d}{dx}(y)\right)$$

The third-, fourth- and \textit{n}th-derivative operators break down in the

same way. Thus, much like a first-order separable ODE is reducible to

the form



$$\frac{dy}{dx}=f(y)g(x)$$ a separable second-order ODE is reducible to

the form



$$\frac{d^2y}{dx^2}=f\left(y'\right)g(x)$$ and an nth-order separable

ODE is reducible to



$$\frac{d^ny}{dx^n}=f\!\left(y^{(n-1)}\right)g(x)$$



\paragraph{Example}



Consider the simple nonlinear second-order differential

equation$$y''=(y')^2.$$This equation is an equation only of \textit{y\'\'} and

\textit{y\'}, meaning it is reducible to the general form described above and

is, therefore, separable. Since it is a second-order separable equation,

collect all \textit{x} variables on one side and all \textit{y\'} variables on the

other to get$$\frac{d(y')}{(y')^2}=dx.$$Now, integrate the right side

with respect to \textit{x} and the left with respect to

\textit{y}\'$$\int \frac{d(y')}{(y')^2}=\int dx.$$This

gives$-\frac{1}{y'}=x+C_1,$which simplifies

to$$y'=-\frac{1}{x+C_1}~.$$This is now a simple integral problem that

gives the final answer$$y=C_2-\ln|x+C_1|.$$



\paragraph{Partial differential equations}



The method of separation of variables is also used to solve a wide range

of linear partial differential equations with boundary and initial

conditions, such as the heat equation, wave equation, Laplace equation,

Helmholtz equation and biharmonic equation.



The analytical method of separation of variables for solving partial

differential equations has also been generalized into a computational

method of decomposition in invariant structures that can be used to

solve systems of partial differential equations.



\paragraph{Example: homogeneous case}



Consider the one-dimensional heat equation. The equation is



$$\frac{\partial u}{\partial t} - \alpha\frac{\partial^{2}u}{\partial x^{2}} = 0$$

\textbf{(1)}



The variable \textit{u} denotes temperature. The boundary condition is

homogeneous, that is



$$u\big|_{x=0}=u\big|_{x=L}=0$$ \textbf{(2)}



Let us attempt to find a solution which is not identically zero

satisfying the boundary conditions but with the following property: \textit{u}

is a product in which the dependence of \textit{u} on \textit{x}, \textit{t} is separated,

that is:



$$u(x,t) = X(x) T(t).$$ \textbf{(3)}



Substituting \textit{u} back into equation \textbf{(1)} and using the product rule,



$$\frac{T'(t)}{\alpha T(t)} = \frac{X''(x)}{X(x)}.$$ \textbf{(4)}



Since the right hand side depends only on \textit{x} and the left hand side

only on \textit{t}, both sides are equal to some constant value -\textit{?}. Thus:



$$T'(t) = - \lambda \alpha T(t),$$ \textbf{(5)}



and



$$X''(x) = - \lambda X(x).$$ \textbf{(6)}



-\textit{?} here is the eigenvalue for both differential operators, and

\textit{T}(\textit{t}) and \textit{X}(\textit{x}) are corresponding eigenfunctions.



We will now show that solutions for \textit{X}(\textit{x}) for values of \textit{?} = 0

cannot occur:



Suppose that \textit{?} \< 0. Then there exist real numbers \textit{B}, \textit{C} such that



$$X(x) = B e^{\sqrt{-\lambda} \, x} + C e^{-\sqrt{-\lambda} \, x}.$$



From \textbf{(2)} we get



$$X(0) = 0 = X(L),$$ \textbf{(7)}



and therefore \textit{B} = 0 = \textit{C} which implies \textit{u} is identically 0.



Suppose that \textit{?} = 0. Then there exist real numbers \textit{B}, \textit{C} such that



$$X(x) = Bx + C.$$



From \textbf{(7)} we conclude in the same manner as in 1 that \textit{u} is

identically 0.



Therefore, it must be the case that \textit{?} \> 0. Then there exist real

numbers \textit{A}, \textit{B}, \textit{C} such that



$$T(t) = A e^{-\lambda \alpha t},$$ and



$$X(x) = B \sin(\sqrt{\lambda} \, x) + C \cos(\sqrt{\lambda} \, x).$$



From \textbf{(7)} we get \textit{C} = 0 and that for some positive integer \textit{n},



$$\sqrt{\lambda} = n \frac{\pi}{L}.$$



This solves the heat equation in the special case that the dependence of

\textit{u} has the special form of \textbf{(3)}.



In general, the sum of solutions to \textbf{(1)} which satisfy the boundary

conditions \textbf{(2)} also satisfies \textbf{(1)} and \textbf{(3)}. Hence a complete

solution can be given as



$$u(x,t) = \sum_{n = 1}^{\infty} D_n \sin \frac{n\pi x}{L} \exp\left(-\frac{n^2 \pi^2 \alpha t}{L^2}\right),$$



where \textit{D}\textasciitilde{}\textit{n}\textasciitilde{} are coefficients determined by initial condition.



Given the initial condition



$$u\big|_{t=0}=f(x),$$



we can get



$$f(x) = \sum_{n = 1}^{\infty} D_n \sin \frac{n\pi x}{L}.$$



This is the sine series expansion of \textit{f}(\textit{x}). Multiplying both sides

with $\sin \frac{n\pi x}{L}$ and integrating over \[0, \textit{L}\] result in



$$D_n = \frac{2}{L} \int_0^L f(x) \sin \frac{n\pi x}{L} \, dx.$$



This method requires that the eigenfunctions \textit{X}, here

$\left\{\sin \frac{n\pi x}{L}\right\}_{n=1}^{\infty}$, are orthogonal

and complete. In general this is guaranteed by Sturm-Liouville theory.



\paragraph{Example: nonhomogeneous case}



Suppose the equation is nonhomogeneous,



$$\frac{\partial u}{\partial t}-\alpha\frac{\partial^{2}u}{\partial x^{2}}=h(x,t)$$

\textbf{(8)}



with the boundary condition the same as \textbf{(2)}.



Expand \textit{h}(\textit{x,t}), \textit{u}(\textit{x},\textit{t}) and \textit{f}(\textit{x}) into



$$h(x,t)=\sum_{n=1}^{\infty}h_{n}(t)\sin\frac{n\pi x}{L},$$ \textbf{(9)}



$$u(x,t)=\sum_{n=1}^{\infty}u_{n}(t)\sin\frac{n\pi x}{L},$$ \textbf{(10)}



:   $f(x)=\sum_{n=1}^{\infty}b_{n}\sin\frac{n\pi x}{L},$ \textbf{(11)}



where \textit{h}\textasciitilde{}\textit{n}\textasciitilde{}(\textit{t}) and \textit{b}\textasciitilde{}\textit{n}\textasciitilde{} can be calculated by integration, while

\textit{u}\textasciitilde{}\textit{n}\textasciitilde{}(\textit{t}) is to be determined.



Substitute \textbf{(9)} and \textbf{(10)} back to \textbf{(8)} and considering the

orthogonality of sine functions we get



:   $u'_{n}(t)+\alpha\frac{n^{2}\pi^{2}}{L^{2}}u_{n}(t)=h_{n}(t),$



which are a sequence of linear differential equations that can be

readily solved with, for instance, Laplace transform, or Integrating

factor. Finally, we can get



:   $u_{n}(t)=e^{-\alpha\frac{n^{2}\pi^{2}}{L^{2}} t} \left (b_{n}+\int_{0}^{t}h_{n}(s)e^{\alpha\frac{n^{2}\pi^{2}}{L^{2}} s} \, ds \right).$



If the boundary condition is nonhomogeneous, then the expansion of

\textbf{(9)} and \textbf{(10)} is no longer valid. One has to find a function \textit{v}

that satisfies the boundary condition only, and subtract it from \textit{u}.

The function \textit{u-v} then satisfies homogeneous boundary condition, and

can be solved with the above method.



\paragraph{Example: mixed derivatives}



For some equations involving mixed derivatives, the equation does not

separate as easily as the heat equation did in the first example above,

but nonetheless separation of variables may still be applied. Consider

the two-dimensional biharmonic equation



$$\frac{\partial^4 u}{\partial x^4} + 2\frac{\partial^4 u}{\partial x^2\partial y^2} + \frac{\partial^4 u}{\partial y^4} = 0.$$



Proceeding in the usual manner, we look for solutions of the form



$$u(x,y) = X(x)Y(y)$$



and we obtain the equation



$$\frac{X^{(4)}(x)}{X(x)} + 2\frac{X''(x)}{X(x)}\frac{Y''(y)}{Y(y)} + \frac{Y^{(4)}(y)}{Y(y)} = 0.$$



Writing this equation in the form



$$E(x) + F(x)G(y) + H(y) = 0,$$



we see that the derivative with respect to \textit{x} and \textit{y} eliminates the

first and last terms, so that



$$F'(x)G'(y) = 0,$$



i.e. either \textit{F}(\textit{x}) or \textit{G}(\textit{y}) must be a constant, say -?. This

further implies that either $-E(x)=F(x)G(y)+H(y)$ or

$-H(y)=E(x)+F(x)G(y)$ are constant. Returning to the equation for \textit{X}

and \textit{Y}, we have two cases



$$\begin{aligned}

    X''(x) &= -\lambda_1X(x) \\\    X^{(4)}(x) &= \mu_1X(x) \\\    Y^{(4)}(y) - 2\lambda_1Y''(y) &= -\mu_1Y(y)

\end{aligned}$$



and



$$\begin{aligned}

    Y''(y) &= -\lambda_2Y(y) \\\    Y^{(4)}(y) &= \mu_2Y(y) \\\    X^{(4)}(x) - 2\lambda_2X''(x) &= -\mu_2X(x)

\end{aligned}$$



which can each be solved by considering the separate cases for

$\lambda_i0$ and noting that

$\mu_i=\lambda_i^2$.



\paragraph{Curvilinear coordinates}



In orthogonal curvilinear coordinates, separation of variables can still

be used, but in some details different from that in Cartesian

coordinates. For instance, regularity or periodic condition may

determine the eigenvalues in place of boundary conditions. See spherical

harmonics for example.



\paragraph{Applicability}



\textbf{Partial differential equations}



For many PDEs, such as the wave equation, Helmholtz equation and

Schrodinger equation, the applicability of separation of variables is a

result of the spectral theorem. In some cases, separation of variables

may not be possible. Separation of variables may be possible in some

coordinate systems but not others, and which coordinate systems allow

for separation depends on the symmetry properties of the equation. Below

is an outline of an argument demonstrating the applicability of the

method to certain linear equations, although the precise method may

differ in individual cases (for instance in the biharmonic equation

above).



Consider an initial boundary value problem for a function $u(x,t)$ on

$D = \{(x,t): x \in [0,l], t \geq 0 \}$ in two variables:



$$(Tu)(x,t) = (Su)(x,t)$$



where $T$ is a differential operator with respect to $x$ and $S$ is a

differential operator with respect to $t$ with boundary data:



$$(Tu)(0,t) = (Tu)(a) = 0$$ for $t \geq 0$



$$(Su)(x,0)=h(x)$$ for $0 \leq x \leq a$



where $h$ is a known function.



We look for solutions of the form $u(x,t) = f(x) g(t)$. Dividing the PDE

through by $f(x)g(t)$ gives



$$\frac{Tf}{f} = \frac{Sg}{g}$$



The right hand side depends only on $x$ and the left hand side only on

$t$ so both must be equal to a constant $K$, which gives two ordinary

differential equations



$$Tf = Kf, Sg = Kg$$



which we can recognize as eigenvalue problems for the operators for $T$

and $S$. If $T$ is a compact, self-adjoint operator on the space

$L^2[0,l]$ along with the relevant boundary conditions, then by the

Spectral theorem there exists a basis for $L^2[0,l]$ consisting of

eigenfunctions for $T$. Let the spectrum of $T$ be $E$ and let

$f_{\lambda}$ be an eigenfunction with eigenvalue $\lambda \in E$. Then

for any function which at each time $t$ is square-integrable with

respect to $x$, we can write this function as a linear combination of

the $f_{\lambda}$. In particular, we know the solution $u$ can be

written as



$$u(x,t) = \sum_{\lambda \in E} c_{\lambda}(t)f_{\lambda}(x)$$



For some functions $c_{\lambda}(t)$. In the separation of variables,

these functions are given by solutions to $Sg = Kg$



Hence, the spectral theorem ensures that the separation of variables

will (when it is possible) find all the solutions.



For many differential operators, such as $\frac{d^2}{dx^2}$, we can show

that they are self-adjoint by integration by parts. While these

operators may not be compact, their inverses (when they exist) may be,

as in the case of the wave equation, and these inverses have the same

eigenfunctions and eigenvalues as the original operator (with the

possible exception of zero).



\paragraph{Matrices}



The matrix form of the separation of variables is the Kronecker sum.



As an example we consider the 2D discrete Laplacian on a regular grid:



$$L = \mathbf{D_{xx}}\oplus\mathbf{D_{yy}}=\mathbf{D_{xx}}\otimes\mathbf{I}+\mathbf{I}\otimes\mathbf{D_{yy}}, \,$$



where $\mathbf{D_{xx}}$ and $\mathbf{D_{yy}}$ are 1D discrete Laplacians

in the \textit{x}- and \textit{y}-directions, correspondingly, and $\mathbf{I}$ are

the identities of appropriate sizes. See the main article Kronecker sum

of discrete Laplacians for details.



\paragraph{Software}



Some mathematical programs are able to do separation of variables: Xcas

among others.



\paragraph{Resources}



\begin{itemize}
    \item \href{https://tutorial.math.lamar.edu/Classes/DE/Separable.aspx}{Separation of Variables
\end{itemize}
    Notes}.

    Produced by Paul Dawkins, Lamar University



\paragraph{Licensing}



Content obtained and/or adapted from:



\begin{itemize}
    \item \href{https://en.wikipedia.org/wiki/Separation_of_variables}{Separation of variables,
\end{itemize}
    Wikipedia}

    under a CC BY-SA license"

%%===============================================================
\section{Homogeneous Differential Equations (1st Order)}
\label{sec:homogeneous-differential-equations-1st-order}
%%===============================================================

"A differential equation can be \textbf{homogeneous} in either of two

respects.



A first order differential equation is said to be homogeneous if it may

be written



$$f(x,y) \, dy = g(x,y) \, dx,$$ where `{{mvar|f}}`{=mediawiki} and

`{{mvar|g}}`{=mediawiki} are homogeneous functions of the same degree of

`{{mvar|x}}`{=mediawiki} and `{{mvar|y}}`{=mediawiki}. In this case, the

change of variable `{{math|1=''y'' = ''ux''}}`{=mediawiki} leads to an

equation of the form



$$\frac{dx}{x} = h(u) \, du,$$ which is easy to solve by integration of

the two members.



Otherwise, a differential equation is homogeneous if it is a homogeneous

function of the unknown function and its derivatives. In the case of

linear differential equations, this means that there are no constant

terms. The solutions of any linear ordinary differential equation of any

order may be deduced by integration from the solution of the homogeneous

equation obtained by removing the constant term.



\paragraph{Homogeneous first-order differential equations}



A first-order ordinary differential equation in the form:



$$M(x,y)\,dx + N(x,y)\,dy = 0$$



is a homogeneous type if both functions

`{{math|''M''(''x'', ''y'')}}`{=mediawiki} and

`{{math|''N''(''x'', ''y'')}}`{=mediawiki} are homogeneous functions of

the same degree `{{mvar|n}}`{=mediawiki}. That is, multiplying each

variable by a parameter `{{math|?}}`{=mediawiki}, we find



$$M(\lambda x, \lambda y) = \lambda^n M(x,y) \quad \text{and} \quad N(\lambda x, \lambda y) = \lambda^n N(x,y)\,.$$



Thus,



$$\frac{M(\lambda x, \lambda y)}{N(\lambda x, \lambda y)} = \frac{M(x,y)}{N(x,y)}\,.$$



\paragraph{Solution method}



In the quotient $\frac{M(tx,ty)}{N(tx,ty)} = \frac{M(x,y)}{N(x,y)}$, we

can let $t = \frac{1}{x}$ to simplify this quotient to a function

`{{mvar|f}}`{=mediawiki} of the single variable $\frac{y}{x}$:



$$\frac{M(x,y)}{N(x,y)} = \frac{M(tx,ty)}{N(tx,ty)} = \frac{M(1,y/x)}{N(1,y/x)}=f(y/x)\,.$$

That is



$$\frac{dy}{dx} = -f(y/x).$$



Introduce the change of variables

`{{math|1=''y'' = ''ux''}}`{=mediawiki}; differentiate using the product

rule:



$$\frac{dy}{dx}=\frac{d(ux)}{dx} = x\frac{du}{dx} + u\frac{dx}{dx} = x\frac{du}{dx} + u.$$



This transforms the original differential equation into the separable

form



:   $x\frac{du}{dx} = -f(u) - u,$



or



:   $\frac 1x\frac{dx}{du} = \frac {-1}{f(u) + u},$



which can now be integrated directly: `{{math|ln ''x''}}`{=mediawiki}

equals the antiderivative of the right-hand side (see ordinary

differential equation).



\paragraph{Special case}



A first order differential equation of the form

(`{{mvar|a}}`{=mediawiki}, `{{mvar|b}}`{=mediawiki},

`{{mvar|c}}`{=mediawiki}, `{{mvar|e}}`{=mediawiki},

`{{mvar|f}}`{=mediawiki}, `{{mvar|g}}`{=mediawiki} are all constants)



$$\left(ax + by + c\right) dx + \left(ex + fy + g\right) dy = 0$$ where

`{{math|''af'' ? ''be''}}`{=mediawiki} can be transformed into a

homogeneous type by a linear transformation of both variables

(`{{mvar|a}}`{=mediawiki} and `{{mvar|ß}}`{=mediawiki} are constants):



$$t = x + \alpha; \; \; z = y + \beta \,.$$



\paragraph{Homogeneous linear differential equations}



A linear differential equation is \textbf{homogeneous} if it is a homogeneous

linear equation in the unknown function and its derivatives. It follows

that, if `{{math|''f''(''x'')}}`{=mediawiki} is a solution, so is

`{{math|''cf''(''x'')}}`{=mediawiki}, for any (non-zero) constant

`{{mvar|c}}`{=mediawiki}. In order for this condition to hold, each

nonzero term of the linear differential equation must depend on the

unknown function or any derivative of it. A linear differential equation

that fails this condition is called \textbf{inhomogeneous.}



A linear differential equation can be represented as a linear operator

acting on `{{math|''y''(''x'')}}`{=mediawiki} where

`{{mvar|x}}`{=mediawiki} is usually the independent variable and

`{{mvar|y}}`{=mediawiki} is the dependent variable. Therefore, the

general form of a linear homogeneous differential equation is



$$L(y) = 0$$



where `{{mvar|L}}`{=mediawiki} is a differential operator, a sum of

derivatives (defining the \""0th derivative\"" as the original,

non-differentiated function), each multiplied by a function

`{{math|''f''''i''}}`{=mediawiki} of

`{{mvar|x}}`{=mediawiki}:



$$L = \sum_{i=0}^n f_i(x)\frac{d^i}{dx^i} \, ,$$ where

`{{math|''f''''i''}}`{=mediawiki} may be constants, but not

all `{{math|''f''''i''}}`{=mediawiki} may be zero.



For example, the following linear differential equation is homogeneous:



$$\sin(x) \frac{d^2y}{dx^2} + 4 \frac{dy}{dx} + y = 0 \,,$$



whereas the following two are inhomogeneous:



$$2 x^2 \frac{d^2y}{dx^2} + 4 x \frac{dy}{dx} + y = \cos(x) \,;$$



$$2 x^2 \frac{d^2y}{dx^2} - 3 x \frac{dy}{dx} + y = 2 \,.$$ The

existence of a constant term is a sufficient condition for an equation

to be inhomogeneous, as in the above example.



\paragraph{Second order equations}



In order to show how these equations are solved, lets start with the

most basic case - a second order equation



$$A\frac{d^2y}{dx^2}+B\frac{dy}{dx}+Cy=0$$



where A, B, and C are constants.



\paragraph{The Auxiliary Quadratic}



For a DE



$$a\frac{d^2y}{dx^2}+b\frac{dy}{dx}+cy=0,$$



Make the following substitution:



$$y=e^{mx}\,$$



This gives also



$$\frac{dy}{dx}=me^{mx}\,$$



$$\frac{d^2y}{dx^2}=m^2e^{mx}\,$$



The DE is now



$$am^2e^{mx}+bme^{mx}+ce^{mx}=0 \,$$



Dividing by $e^{mx}$ gives (note$$e^{mx}$$ can never equal zero)



$$am^2+bm+c=0 \,$$



This is the \textbf{auxiliary quadratic} (AQ) of the DE. There are four

classes of outcomes to the auxiliary quadratic:



1.  $b^2-4ac > 0, \,$, giving two distinct, real, roots.

2.  $b^2-4ac = 0, \,$, giving two coincident real, roots.

3.  $b^2-4ac < 0, \,$, giving complex roots.



:   \textit{a}: Purely imaginary roots.

:   \textit{b}: Complex-conjugate pair.



The method of solution of the DE depends on the class of AQ.



\paragraph{Class 1: Distinct, Real Roots}



Consider the DE



$$\frac{d^2y}{dx^2}+\frac{dy}{dx}-6y=0$$



The AQ is



$$m^2+m-6=0 \,$$



$$(m+3)(m-2)=0 \,$$



This gives us the following roots:



$$m=2, \ m=-3$$



Going back to the substitution we made to obtain the AQ, we have



$$y=e^{2x},\ y=e^{-3x}$$



as two distinct solutions to the DE. According to the superposition

principle, the general solution is therefore



$$y=Ae^{2x}+Be^{-3x} \,$$



\paragraph{General Solution to Class 1 DEs}



:   Generalizing, for the Second Order DE



$$a\frac{d^2y}{dx^2}+b\frac{dy}{dx}+cy=0$$



:   with the auxiliary quadratic given by



$$am^2+bm+c=0 \,$$



:   with roots \textit{a} and \textit{ß}, the general solution is



$$y=Ae^{\alpha x}+Be^{\beta x} \,$$



\paragraph{Class 2: Coincident, Real Roots}



Consider the DE



$$\frac{d^2y}{dx^2}-4\frac{dy}{dx}+4y=0$$



The AQ is



$$m^2-4m+4=0 \,$$



$$(m-2)(m-2)=0 \,$$



so, $e^{2x}\,$ is a solution. However, we cannot have it as both

solutions as the factor of two produced will be absorbed into the

constant, leaving us with only one constant, and therefore a DE without

a full solution.



For the other solution we will use the Method of Reduction of Order. To

do so we assume that it is in the form of:



$$y_2=u(x)y_1=u(x)e^{2x}\,$$ At the end we will check if our assumption

is correct. We will now substitute this equation and solve for u(x)



$$y_2''-4y_2'+4y_2=0\,$$



$$u''(x)e^{2x}+2u'(x)e^{2x}+2u'(x)e^{2x}+4u(x)e^{2x}-4(u'(x)e^{2x}+2u(x)e^{2x})+4u(x)e^{2x}=0\,$$



$$u''(x)e^{2x}=0\,$$



$$e^{2x}\,$$ is always non-zero so the only way the product can equal

zero is if:



$$u''(x)=0\,$$ Integrating twice offers



$$u(x)=a_1x + a_2\,$$ Therefore



$$y_2=(a_1x+a_2)y_1\,$$



$$y_2=(a_1x+a_2)e^{2x}\,$$ Our general solution is:



$$y=C_1y_1+C_2y_2\,$$



$$y=C_1e^{2x}+C_2(a_1x+a_2)e^{2x}\,$$



$$y=(C_1+a_2)e^{2x}+C_2a_1xe^{2x}\,$$ Because each constant is arbitrary

we can simply write



$$y=C_1e^{2x}+C_2xe^{2x}\,$$



The Method of Reduction of Order can be used on different equations and

u(x) does not always equal x. You can see below that $y=xe^{x}$ is a

valid solution.



$$y=xe^{2x}\,$$



$$\frac{dy}{dx}=2xe^{2x}+e^{2x}=e^{2x}(2x+1)\,$$



$$\frac{d^2y}{dx^2}=2e^{2x}+2e^{2x}(2x+1)=e^{2x}(4x+4)\,$$



To check substitute these into the original DE:



$$e^{2x}(4x+4)-4e^{2x}(2x+1)+4xe^{2x}=0 \,$$



$$4xe^{2x}+4e^{2x}-8xe^{2x}-4e^{2x}+4xe^{2x}=0 \,$$



$$0=0\,$$



Therefore, $y=xe^{2x}\,$ is a solution as well.



\paragraph{General Solution to Class 2 DEs}



:   Generalizing, for the Second Order DE



$$a\frac{d^2y}{dx^2}+b\frac{dy}{dx}+cy=0$$



:   with the coincident root \textit{a} of the AQ, the general solution is



$$y=(A+Bx)e^{\alpha x} \,$$



\paragraph{Class 3a: Purely Imaginary Roots}



To have complex roots, the AQ must have a discriminant less than zero,

so



$$b^2-4ac < 0 \,$$



Also, for the solution to be purely imaginary, the value of \textit{b} must be

exactly zero.



Therefore,



$$4ac>0\,$$



This means that \textit{a} and \textit{c} have to have the same signs: either \textit{a} and

\textit{c} are both positive or they are both negative. If we consider our

general second-order DE:



$$a\frac{d^2y}{dx^2}+b\frac{dy}{dx}+cy=0$$



Setting \textit{b} to zero gives



$$a\frac{d^2y}{dx^2}+cy=0$$



Dividing through by \textit{a} gives



$$\frac{d^2y}{dx^2}+\frac{c}{a}y=0$$.



Therefore, the \textit{y} term is always positive, and this can be represented

by



$$\frac{d^2y}{dx^2}+\omega^2y=0$$.



(I\'m using \textit{?} here as it is used for simple harmonic motion, which is

the primary use of this DE). There are now two paths to the solution of

the DE. The first relies on us spotting that we can use the cyclical

nature of trig. functions when derived. Substitute the following



$$y=\cos \omega x \,$$



$$\frac{dy}{dx}=-\omega \sin \omega x\,$$



$$\frac{d^2y}{dx^2}=-\omega^2 \cos \omega x\,$$



And check in our DE:



$$\frac{d^2y}{dx^2}+\omega^2 y=-\omega^2 \cos \omega x+\omega^2 \cos \omega x=0$$



This checks out, so $\cos \omega x\,$ is a solution. A similar result

holds true for the substitution using $y=\sin \omega x \,$.



Our solutions are therefore



$$y=\cos \omega x \,$$



$$y=\sin \omega x \,$$



So the general solution is



$$y=A\cos \omega x+B\sin \omega x. \,$$



The other method of solving this equation is to use Euler\'s Formula:



$$e^{\omega ix}=\cos \omega x +i \sin x \,$$



:   and



$$e^{-\omega ix}=\cos \omega x -i \sin x \,$$



From our original DE, we have an AQ of



$$m^2+\omega^2=0 \,$$



giving us roots of



$$m=\pm i \omega \,$$



so the general solution, similar to the Class 1 DEs, is



$$y=Ae^{\omega ix}+Be^{-\omega ix} \,$$



$$y=A(\cos \omega x +i \sin \omega x)+B(\cos \omega x -i \sin \omega x)\,$$



$$y=(A+B)\cos \omega x+i(A-B)\sin \omega x \,$$



Since A and B are arbitrary, we can set new constants for convenience,

letting our new A equal A+B and our new B equal \textit{i}(A-B).



Thus we have as our general solution



$$y=A\cos \omega x+B\sin \omega x \,$$



\paragraph{General Solution to Class 3a DEs}



:   Generalizing, for the Second Order DE



$$\frac{d^2y}{dx^2}+\omega^2 y=0$$



:   the general solution is



$$y=A\cos \omega x+B\sin \omega x \,$$



\paragraph{Class 3b: Complex Conjugate Roots}



Since it is a proven theorem that complex roots of polynomials always

occur in conjugate pairs, the only remaining class of AQ is the one with

complex conjugates for solutions.



Given that the solutions are complex, we know that in the AQ



$$am^2+bm+c=0 \,$$



$$b^2-4ac,<0,\ b \ne 0 \,$$ (see Class 3a).



The roots of this are in the form



$$(p \pm iq)\,$$



The general solution is then



$$y=Ae^{(p+iq)x}+Be^{(p-iq)x}\,$$



$$y=e^{px} \left( Ae^{iqx}+Be^{-iqx}\right)\,$$



From Euler\'s Formulas, we can now get



$$y=e^{px}\left( A\left(\cos q x + i \sin q x \right)+B\left(\cos q x - i \sin q x \right) \right)\,$$



$$y=e^{px}\left( (A+B)\cos q x+i(A-B)\sin q x \right)\,$$



As A and B are arbitrary, we can collapse them as in Class 3a, so that

we have the general solution



$$y=e^{px}\left( A\cos q x+B\sin q x \right)\,$$



\paragraph{General Solution to Class 3b DEs}



:   Generalizing, for the Second Order DE



$$a\frac{d^2y}{dx^2}+b\frac{dy}{dx}+cy=0$$



:   with an AQ with roots



$$m=p \pm iq,$$



:   The general solution is



$$y=e^{px}\left( A\cos q x+B\sin q x \right)\,$$



We have now covered all possible types of homogeneous second-order

differential equation, and we didn\'t even have to integrate anything!

We will now have a look at higher order equivalents.



\paragraph{Licensing}



Content obtained and/or adapted from:



\begin{itemize}
    \item \href{https://en.wikipedia.org/wiki/Homogeneous_differential_equation}{Homogeneous differential equation,
\end{itemize}
    Wikipedia}

    under a CC BY-SA license

\begin{itemize}
    \item \href{https://en.wikibooks.org/wiki/Ordinary_Differential_Equations/Homogenous_1}{Homogeneous 1, Wikibooks: Ordinary Differential
\end{itemize}
    Equations}

    under a CC BY-SA license"

%%===============================================================
\section{Linear Differential Equations (1st Order)}
\label{sec:linear-differential-equations-1st-order}
%%===============================================================

"In mathematics, a \textbf{linear differential equation} is a differential

equation that is defined by a linear polynomial in the unknown function

and its derivatives, that is an equation of the form



$$a_0(x)y + a_1(x)y' + a_2(x)y'' \cdots + a_n(x)y^{(n)} = b(x)$$ where

\textit{a}\textasciitilde{}0\textasciitilde{}(\textit{x}), ..., \textit{a}\textasciitilde{}\textit{n}\textasciitilde{}(\textit{x}) and `{{math|''b''(''x'')}}`{=mediawiki}

are arbitrary differentiable functions that do not need to be linear,

and `{{math|''y''', …, ''y''(''n'') }}`{=mediawiki} are the

successive derivatives of an unknown function `{{mvar|y}}`{=mediawiki}

of the variable `{{mvar|x}}`{=mediawiki}.



Such an equation is an ordinary differential equation (ODE). A *linear

differential equation* may also be a linear partial differential

equation (PDE), if the unknown function depends on several variables,

and the derivatives that appear in the equation are partial derivatives.



A linear differential equation or a system of linear equations such that

the associated homogeneous equations have constant coefficients may be

solved by quadrature, which means that the solutions may be expressed in

terms of integrals. This is also true for a linear equation of order

one, with non-constant coefficients. An equation of order two or higher

with non-constant coefficients cannot, in general, be solved by

quadrature. For order two, Kovacic\'s algorithm allows deciding whether

there are solutions in terms of integrals, and computing them if any.



The solutions of linear differential equations with polynomial

coefficients are called holonomic functions. This class of functions is

stable under sums, products, differentiation, integration, and contains

many usual functions and special functions such as exponential function,

logarithm, sine, cosine, inverse trigonometric functions, error

function, Bessel functions and hypergeometric functions. Their

representation by the defining differential equation and initial

conditions allows making algorithmic (on these functions) most

operations of calculus, such as computation of antiderivatives, limits,

asymptotic expansion, and numerical evaluation to any precision, with a

certified error bound.



\paragraph{Basic terminology}



The highest order of derivation that appears in a (linear) differential

equation is the \textit{order} of the equation. The term

`{{math|''b''(''x'')}}`{=mediawiki}, which does not depend on the

unknown function and its derivatives, is sometimes called the *constant

term* of the equation (by analogy with algebraic equations), even when

this term is a non-constant function. If the constant term is the zero

function, then the differential equation is said to be \textit{homogeneous}, as

it is a homogeneous polynomial in the unknown function and its

derivatives. The equation obtained by replacing, in a linear

differential equation, the constant term by the zero function is the

\textit{associated homogeneous equation}. A differential equation has *constant

coefficients* if only constant functions appear as coefficients in the

associated homogeneous equation.



A \textit{solution} of a differential equation is a function that satisfies the

equation. The solutions of a homogeneous linear differential equation

form a vector space. In the ordinary case, this vector space has a

finite dimension, equal to the order of the equation. All solutions of a

linear differential equation are found by adding to a particular

solution any solution of the associated homogeneous equation.



\paragraph{Linear differential operator}



A \textit{basic differential operator} of order `{{mvar|i}}`{=mediawiki} is a

mapping that maps any differentiable function to its

`{{mvar|i}}`{=mediawiki}th derivative, or, in the case of several

variables, to one of its partial derivatives of order

`{{mvar|i}}`{=mediawiki}. It is commonly denoted



$$\frac{d^i}{dx^i}$$ in the case of univariate functions, and



$$\frac{\partial^{i_1+\cdots +i_n}}{\partial x_1^{i_1}\cdots \partial x_n^{i_n}}$$

in the case of functions of `{{mvar|n}}`{=mediawiki} variables. The

basic differential operators include the derivative of order 0, which is

the identity mapping.



A \textbf{linear differential operator} (abbreviated, in this article, as

\textit{linear operator} or, simply, \textit{operator}) is a linear combination of

basic differential operators, with differentiable functions as

coefficients. In the univariate case, a linear operator has thus the

form



$$a_0(x)+a_1(x)\frac{d}{dx} + \cdots +a_n(x)\frac{d^n}{dx^n},$$ where

`{{math|''a''0(''x''), …, ''a''''n''(''x'')}}`{=mediawiki}

are differentiable functions, and the nonnegative integer

`{{mvar|n}}`{=mediawiki} is the \textit{order} of the operator (if

`{{math|''a''''n''(''x'')}}`{=mediawiki} is not the zero

function).



Let `{{mvar|L}}`{=mediawiki} be a linear differential operator. The

application of `{{mvar|L}}`{=mediawiki} to a function

`{{mvar|f}}`{=mediawiki} is usually denoted

`{{math|''Lf''}}`{=mediawiki} or `{{math|''Lf''(''X'')}}`{=mediawiki},

if one needs to specify the variable (this must not be confused with a

multiplication). A linear differential operator is a linear operator,

since it maps sums to sums and the product by a scalar to the product by

the same scalar.



As the sum of two linear operators is a linear operator, as well as the

product (on the left) of a linear operator by a differentiable function,

the linear differential operators form a vector space over the real

numbers or the complex numbers (depending on the nature of the functions

that are considered). They form also a free module over the ring of

differentiable functions.



The language of operators allows a compact writing for differentiable

equations: if



$$L=a_0(x)+a_1(x)\frac{d}{dx} + \cdots +a_n(x)\frac{d^n}{dx^n},$$ is a

linear differential operator, then the equation



$$a_0(x)y +a_1(x)y' + a_2(x)y'' +\cdots +a_n(x)y^{(n)}=b(x)$$ may be

rewritten



$$Ly=b(x).$$



There may be several variants to this notation; in particular the

variable of differentiation may appear explicitly or not in

`{{mvar|y}}`{=mediawiki} and the right-hand and of the equation, such as

`{{math|1=''Ly''(''x'') = ''b''(''x'')}}`{=mediawiki} or

`{{math|1=''Ly'' = ''b''}}`{=mediawiki}.



The \textit{kernel} of a linear differential operator is its kernel as a linear

mapping, that is the vector space of the solutions of the (homogeneous)

differential equation `{{math|1=''Ly'' = 0}}`{=mediawiki}.



In the case of an ordinary differential operator of order

`{{mvar|n}}`{=mediawiki}, Carathéodory\'s existence theorem implies

that, under very mild conditions, the kernel of `{{mvar|L}}`{=mediawiki}

is a vector space of dimension `{{mvar|n}}`{=mediawiki}, and that the

solutions of the equation

`{{math|1=''Ly''(''x'') = ''b''(''x'')}}`{=mediawiki} have the form



$$S_0(x) + c_1S_1(x) + \cdots +c_nS_n(x),$$ where

`{{math|''c''1, …, ''c''''n''}}`{=mediawiki} are

arbitrary numbers. Typically, the hypotheses of Carathéodory\'s theorem

are satisfied in an interval `{{mvar|I}}`{=mediawiki}, if the functions

`{{math|''b'', ''a''0, …, ''a''''n''}}`{=mediawiki}

are continuous in `{{mvar|I}}`{=mediawiki}, and there is a positive real

number `{{mvar|k}}`{=mediawiki} such that \|\textit{a}\textasciitilde{}\textit{n}\textasciitilde{}(\textit{x})\| \> \textit{k} for

every `{{mvar|x}}`{=mediawiki} in `{{mvar|I}}`{=mediawiki}.



\paragraph{Homogeneous equation with constant coefficients}



A homogeneous linear differential equation has \textit{constant coefficients}

if it has the form



$$a_0y + a_1y' + a_2y'' + \cdots + a_n y^{(n)} = 0$$ where

`{{math|''a''1, …, ''a''''n''}}`{=mediawiki} are

(real or complex) numbers. In other words, it has constant coefficients

if it is defined by a linear operator with constant coefficients.



The study of these differential equations with constant coefficients

dates back to Leonhard Euler, who introduced the exponential function

`{{math|''e''''x''}}`{=mediawiki}, which is the unique

solution of the equation `{{math|1=''f''' = ''f''}}`{=mediawiki} such

that `{{math|1=''f''(0) = 1}}`{=mediawiki}. It follows that the

`{{mvar|n}}`{=mediawiki}th derivative of

`{{math|''e''''cx'' }}`{=mediawiki} is

`{{math|''c''''n''''e''''cx''}}`{=mediawiki}, and

this allows solving homogeneous linear differential equations rather

easily.



Let



$$a_0y + a_1y' + a_2y'' + \cdots + a_ny^{(n)} = 0$$ be a homogeneous

linear differential equation with constant coefficients (that is

`{{math|''a''0, …, ''a''''n''}}`{=mediawiki} are

real or complex numbers).



Searching solutions of this equation that have the form

`{{math|''e''''ax''}}`{=mediawiki} is equivalent to searching

the constants `{{mvar|a}}`{=mediawiki} such that



$$a_0e^{\alpha x} + a_1\alpha e^{\alpha x} + a_2\alpha^2 e^{\alpha x}+\cdots + a_n\alpha^n e^{\alpha x} = 0.$$

Factoring out `{{math|''e''''ax''}}`{=mediawiki} (which is

never zero), shows that `{{mvar|a}}`{=mediawiki} must be a root of the

\textit{characteristic polynomial}



$$a_0 + a_1t + a_2t^2 + \cdots + a_nt^n$$ of the differential equation,

which is the left-hand side of the characteristic equation



$$a_0 + a_1t + a_2t^2 + \cdots + a_nt^n = 0.$$



When these roots are all distinct, one has `{{mvar|n}}`{=mediawiki}

distinct solutions that are not necessarily real, even if the

coefficients of the equation are real. These solutions can be shown to

be linearly independent, by considering the Vandermonde determinant of

the values of these solutions at

`{{math|1=''x'' = 0, …, ''n'' – 1}}`{=mediawiki}. Together they form a

basis of the vector space of solutions of the differential equation

(that is, the kernel of the differential operator).



+----------------------------------------------------------------------+

| Example                                                              |

+======================================================================+

| $$y''''-2y'''+2y''-2y'+y=0$$ has the characteristic equation         |

|                                                                      |

| :   $z^4-2z^3+2z^2-2z+1=0.$                                          |

|                                                                      |

| This has zeros, `{{mvar|i}}`{=mediawiki},                            |

| `{{math|-''i''}}`{=mediawiki}, and `{{math|1}}`{=mediawiki}          |

| (multiplicity 2). The solution basis is thus                         |

|                                                                      |

| :   $e^{ix}, \; e^{-ix}, \; e^x, \; xe^x.$                              |

|                                                                      |

| A real basis of solution is thus                                     |

|                                                                      |

| :   $\cos x, \; \sin x, \; e^x, \; xe^x.$                               |

+----------------------------------------------------------------------+



In the case where the characteristic polynomial has only simple roots,

the preceding provides a complete basis of the solutions vector space.

In the case of multiple roots, more linearly independent solutions are

needed for having a basis. These have the form



$$x^ke^{\alpha x},$$ where `{{mvar|k}}`{=mediawiki} is a nonnegative

integer, `{{mvar|a}}`{=mediawiki} is a root of the characteristic

polynomial of multiplicity `{{mvar|m}}`{=mediawiki}, and

`{{math|''k'' < ''m''}}`{=mediawiki}. For proving that these functions

are solutions, one may remark that if `{{mvar|a}}`{=mediawiki} is a root

of the characteristic polynomial of multiplicity

`{{mvar|m}}`{=mediawiki}, the characteristic polynomial may be factored

as

`{{math|''P''(''t'') (''t'' - ''a'')''m''}}`{=mediawiki}.

Thus, applying the differential operator of the equation is equivalent

with applying first `{{mvar|m}}`{=mediawiki} times the operator

$\frac{d}{dx} - \alpha$, and then the operator that has

`{{mvar|P}}`{=mediawiki} as characteristic polynomial. By the

exponential shift theorem,



$$\left(\frac{d}{dx}-\alpha\right)\left(x^ke^{\alpha x}\right)= kx^{k-1}e^{\alpha x},$$



and thus one gets zero after `{{math|''k'' + 1}}`{=mediawiki}

application of $\frac{d}{dx} - \alpha$.



As, by the fundamental theorem of algebra, the sum of the multiplicities

of the roots of a polynomial equals the degree of the polynomial, the

number of above solutions equals the order of the differential equation,

and these solutions form a base of the vector space of the solutions.



In the common case where the coefficients of the equation are real, it

is generally more convenient to have a basis of the solutions consisting

of real-valued functions. Such a basis may be obtained from the

preceding basis by remarking that, if

`{{math|''a'' + ''ib''}}`{=mediawiki} is a root of the characteristic

polynomial, then `{{math|''a'' – ''ib''}}`{=mediawiki} is also a root,

of the same multiplicity. Thus a real basis is obtained by using

Euler\'s formula, and replacing

`{{math|''x''''k''''e''(''a''+''ib'')''x''}}`{=mediawiki}

and

`{{math|''x''''k''''e''(''a''-''ib'')''x''}}`{=mediawiki}

by

`{{math|''x''''k''''e''''ax'' cos(''bx'')}}`{=mediawiki}

and

`{{math|''x''''k''''e''''ax'' sin(''bx'')}}`{=mediawiki}.



\paragraph{Second-order case}



A homogeneous linear differential equation of the second order may be

written



$$y'' + ay' + by = 0,$$ and its characteristic polynomial is



$$r^2 + ar + b.$$



If `{{mvar|a}}`{=mediawiki} and `{{mvar|b}}`{=mediawiki} are real, there

are three cases for the solutions, depending on the discriminant

`{{math|1=''D'' = ''a''2 - 4''b''}}`{=mediawiki}. In all

three cases, the general solution depends on two arbitrary constants

`{{math|''c''1}}`{=mediawiki} and

`{{math|''c''2}}`{=mediawiki}.



\begin{itemize}
    \item If `{{math|''D'' > 0}}`{=mediawiki}, the characteristic polynomial
\end{itemize}
    has two distinct real roots `{{mvar|a}}`{=mediawiki}, and

    `{{mvar|ß}}`{=mediawiki}. In this case, the general solution is



$$c_1 e^{\alpha x} + c_2 e^{\beta x}.$$



\begin{itemize}
    \item If `{{math|1=''D'' = 0}}`{=mediawiki}, the characteristic polynomial
\end{itemize}
    has a double root `{{math|-''a''/2}}`{=mediawiki}, and the general

    solution is



$$(c_1 + c_2 x) e^{-ax/2}.$$



\begin{itemize}
    \item If `{{math|''D'' < 0}}`{=mediawiki}, the characteristic polynomial
\end{itemize}
    has two complex conjugate roots

    `{{math|''a'' ± ''ßi''}}`{=mediawiki}, and the general solution is



$$c_1 e^{(\alpha + \beta i)x} + c_2 e^{(\alpha - \beta i)x},$$



:   which may be rewritten in real terms, using Euler\'s formula as

    $$e^{\alpha x}  (c_1\cos(\beta x) + c_2 \sin(\beta x)).$$



Finding the solution `{{math|''y''(''x'')}}`{=mediawiki} satisfying

`{{math|1=''y''(0) = ''d''1}}`{=mediawiki} and

`{{math|1=''y'''(0) = ''d''2}}`{=mediawiki}, one equates the

values of the above general solution at `{{math|0}}`{=mediawiki} and its

derivative there to `{{math|''d''1}}`{=mediawiki} and

`{{math|''d''2}}`{=mediawiki}, respectively. This results in

a linear system of two linear equations in the two unknowns

`{{math|''c''1}}`{=mediawiki} and

`{{math|''c''2}}`{=mediawiki}. Solving this system gives the

solution for a so-called Cauchy problem, in which the values at

`{{math|0}}`{=mediawiki} for the solution of the DEQ and its derivative

are specified.



\paragraph{Non-homogeneous equation with constant coefficients}



A non-homogeneous equation of order `{{mvar|n}}`{=mediawiki} with

constant coefficients may be written



$$y^{(n)}(x) + a_1 y^{(n-1)}(x) + \cdots + a_{n-1} y'(x)+ a_ny(x) = f(x),$$

where `{{math|''a''1, …, ''a''''n''}}`{=mediawiki}

are real or complex numbers, `{{mvar|f}}`{=mediawiki} is a given

function of `{{mvar|x}}`{=mediawiki}, and `{{mvar|y}}`{=mediawiki} is

the unknown function (for sake of simplicity,

\""`{{math|(''x'')}}`{=mediawiki}\"" will be omitted in the following).



There are several methods for solving such an equation. The best method

depends on the nature of the function `{{mvar|f}}`{=mediawiki} that

makes the equation non-homogeneous. If `{{mvar|f}}`{=mediawiki} is a

linear combination of exponential and sinusoidal functions, then the

exponential response formula may be used. If, more generally,

`{{mvar|f}}`{=mediawiki} is a linear combination of functions of the

form `{{math|''x''''n''''e''''ax''}}`{=mediawiki},

`{{math|''x''''n'' cos(''ax'')}}`{=mediawiki}, and

`{{math|''x''''n'' sin(''ax'')}}`{=mediawiki}, where

`{{mvar|n}}`{=mediawiki} is a nonnegative integer, and

`{{mvar|a}}`{=mediawiki} a constant (which need not be the same in each

term), then the method of undetermined coefficients may be used. Still

more general, the annihilator method applies when

`{{mvar|f}}`{=mediawiki} satisfies a homogeneous linear differential

equation, typically, a holonomic function.



The most general method is the variation of constants, which is

presented here.



The general solution of the associated homogeneous equation



$$y^{(n)} + a_1 y^{(n-1)} + \cdots + a_{n-1} y'+ a_ny = 0$$ is



$$y=u_1y_1+\cdots+ u_ny_n,$$ where

`{{math|(''y''1, …, ''y''''n'')}}`{=mediawiki} is

a basis of the vector space of the solutions and

`{{math|''u''1, …, ''u''''n''}}`{=mediawiki} are

arbitrary constants. The method of variation of constants takes its name

from the following idea. Instead of considering

`{{math|''u''1, …, ''u''''n''}}`{=mediawiki} as

constants, they can considered as unknown functions that have to be

determined for making `{{mvar|y}}`{=mediawiki} a solution of the

non-homogeneous equation. For this purpose, one adds the constraints



$$\begin{aligned}

0 &= u'_1y_1 + u'_2y_2 + \cdots+u'_ny_n \\\0 &= u'_1y'_1 + u'_2y'_2 + \cdots + u'_n y'_n \\\  & \; \;\vdots \\\0 &= u'_1y^{(n-2)}_1+u'_2y^{(n-2)}_2 + \cdots + u'_n y^{(n-2)}_n,

\end{aligned}$$ which imply (by product rule and induction)



$$y^{(i)} = u_1 y_1^{(i)} + \cdots + u_n y_n^{(i)}$$ for

`{{math|1=''i'' = 1, …, ''n'' – 1}}`{=mediawiki}, and



$$y^{(n)} = u_1 y_1^{(n)} + \cdots + u_n y_n^{(n)} +u'_1y_1^{(n-1)}+u'_2y_2^{(n-1)}+\cdots+u'_ny_n^{(n-1)}.$$



Replacing in the original equation `{{mvar|y}}`{=mediawiki} and its

derivatives by these expressions, and using the fact that

`{{math|''y''1, …, ''y''''n''}}`{=mediawiki} are

solutions of the original homogeneous equation, one gets



$$f=u'_1y_1^{(n-1)} + \cdots + u'_ny_n^{(n-1)}.$$



This equation and the above ones with `{{math|0}}`{=mediawiki} as

left-hand side form a system of `{{mvar|n}}`{=mediawiki} linear

equations in

`{{math|''u'''1, …, ''u'''''n''}}`{=mediawiki}

whose coefficients are known functions (`{{mvar|f}}`{=mediawiki}, the

$y_i$, and their derivatives). This system can be solved by any method

of linear algebra. The computation of antiderivatives gives

`{{math|''u''1, …, ''u''''n''}}`{=mediawiki}, and

then

`{{math|1=''y'' = ''u''1''y''1 + ? + ''u''''n''''y''''n''}}`{=mediawiki}.



As antiderivatives are defined up to the addition of a constant, one

finds again that the general solution of the non-homogeneous equation is

the sum of an arbitrary solution and the general solution of the

associated homogeneous equation.



\paragraph{First-order linear equation}



The general form of a linear ordinary differential equation of order 1,

after dividing out the coefficient of

`{{math|''y'''(''x'')}}`{=mediawiki}, is:



$$y'(x) = f(x) y(x) + g(x).$$



If the equation is homogeneous, i.e.

`{{math|1=''g''(''x'') = 0}}`{=mediawiki}, one may rewrite and

integrate:



$$\frac{y'}{y}= f, \qquad \log y = k +F,$$ where

`{{mvar|k}}`{=mediawiki} is an arbitrary constant of integration and

$F=\textstyle\int f\,dx$ is any antiderivative of

`{{mvar|f}}`{=mediawiki}. Thus, the general solution of the homogeneous

equation is



$$y=ce^F,$$ where `{{math|1=''c'' = ''e''''k''}}`{=mediawiki}

is an arbitrary constant.



For the general non-homogeneous equation, one may multiply it by the

reciprocal `{{math|''e''-''F''}}`{=mediawiki} of a solution

of the homogeneous equation. This gives



$$y'e^{-F}-yfe^{-F}= ge^{-F}.$$ As

$-fe^{-F} = \frac{d}{dx} \left(e^{-F}\right),$ the product rule allows

rewriting the equation as



$$\frac{d}{dx}\left(ye^{-F}\right)= ge^{-F}.$$ Thus, the general

solution is



$$y=ce^F + e^F\int ge^{-F}dx,$$ where `{{mvar|c}}`{=mediawiki} is a

constant of integration, and `{{mvar|F}}`{=mediawiki} is any

antiderivative of `{{mvar|f}}`{=mediawiki} (changing of antiderivative

amounts to change the constant of integration).



\paragraph{Example}



Solving the equation



:   $y'(x) + \frac{y(x)}{x} = 3x.$



The associated homogeneous equation $y'(x) + \frac{y(x)}{x} = 0$ gives



$$\frac{y'}{y}=-\frac{1}{x},$$ that is



$$y=\frac{c}{x}.$$



Dividing the original equation by one of these solutions gives



:   $xy'+y=3x^2.$



That is



$$(xy)'=3x^2,$$



$$xy=x^3 +c,$$ and



$$y(x)=x^2+c/x.$$ For the initial condition



:   $y(1)=\alpha,$



one gets the particular solution



$$y(x)=x^2+\frac{\alpha-1}{x}.$$



\paragraph{Making Linear Equations from Non-Linear Equations}



Sometimes a non-linear equation, which is not solvable like this, can be

made linear, and more easily solvable, by applying a substitution.



\paragraph{Example}



$$y\frac{dy}{dx}+xy^2=5x$$



:   Let\'s make the following substitution:



$$v=y^2 \,$$



$$\frac{dv}{dx}=2y\frac{dy}{dx}$$



:   Plugging in, we get



$$\frac{1}{2}\frac{dv}{dx}+xv=5x$$



$$\frac{dv}{dx}+2xv=10x$$



:   We can then solve as a linear equation in \textit{v}, using the

    step-by-step method above:



<!-- -->



:   \textbf{Step 1}: Find the integrating factor:



$$e^{\int P(x)dx}$$



$$e^{\int 2x dx}=e^{x^2+C}$$



$$e^{\int P(x)dx}=Ce^{x^2}$$



:   Letting C=1 for convenience, we get $e^{x^2}$ as our integrating

    factor.



<!-- -->



:   \textbf{Step 2}: Multiply through



$$e^{x^2}\frac{dv}{dx}+e^{x^2}xv=10e^{x^2}x$$



:   \textbf{Step 3}: Recognize that the left hand is

    $\frac{d}{dx} e^{\int P(x)dx}v$



$$\frac{d}{dx} e^{x^2}v=10e^{x^2}x$$



:   \textbf{Step 4}: Integrate both sides w.r.t.\textit{x}.



$$\int (\frac{d}{dx} e^{x^2}v)dx=\int 10e^{x^2}xdx$$



$$e^{x^2}v=5e^{x^2}+C$$



:   \textbf{Step 5}: Solve for \textit{v}.



$$v=5+\frac{C}{e^{x^2}}$$



:   Now that we have \textit{v}, solve for \textit{y}.



$$v=y^2 \,$$



$$y^2=5+\frac{C}{e^{x^2}} \,$$



\paragraph{Higher order linear equations}



A linear ordinary equation of order one with variable coefficients may

be solved by quadrature, which means that the solutions may be expressed

in terms of integrals. This is not the case for order at least two. This

is the main result of Picard--Vessiot theory which was initiated by

Émile Picard and Ernest Vessiot, and whose recent developments are

called differential Galois theory.



The impossibility of solving by quadrature can be compared with the

Abel--Ruffini theorem, which states that an algebraic equation of degree

at least five cannot, in general, be solved by radicals. This analogy

extends to the proof methods and motivates the denomination of

differential Galois theory.



Similarly to the algebraic case, the theory allows deciding which

equations may be solved by quadrature, and if possible solving them.

However, for both theories, the necessary computations are extremely

difficult, even with the most powerful computers.



Nevertheless, the case of order two with rational coefficients has been

completely solved by Kovacic\'s algorithm.



\paragraph{Cauchy--Euler equation}



Cauchy--Euler equations are examples of equations of any order, with

variable coefficients, that can be solved explicitly. These are the

equations of the form



$$x^n y^{(n)}(x) + a_{n-1} x^{n-1} y^{(n-1)}(x) + \cdots + a_0 y(x) = 0,$$

where $a_0, \ldots, a_{n-1}$ are constant coefficients.



\paragraph{Licensing}



Content obtained and/or adapted from:



\begin{itemize}
    \item \href{https://en.wikipedia.org/wiki/Linear_differential_equation}{Linear differential equation,
\end{itemize}
    Wikipedia}

    under a CC BY-SA license

\begin{itemize}
    \item \href{https://en.wikibooks.org/wiki/Ordinary_Differential_Equations/First_Order_Linear_1}{First Order Linear 1, Wikibooks: Differential
\end{itemize}
    Equations}

    under a CC BY-SA license"

%%===============================================================
\section{Integrating Factor}
\label{sec:integrating-factor}
%%===============================================================

The Method of Integrating Factors



Let ``{=html}$p$``{=html} and

``{=html}$g$``{=html} be functions of

``{=html}$t$``{=html} and consider the

following first order differential equation:







$\begin{aligned} \quad \frac{dy}{dt} + p(t) y = g(t) \end{aligned}$







If we multiply the both sides of the equation above by the function

``{=html}$\mu (t)$``{=html} we get

that:







$\begin{aligned} \quad \mu (t) \frac{dy}{dt} + \mu (t) p(t) y = \mu (t) g(t) \end{aligned}$







If we can guarantee that

``{=html}$\mu ' (t) = \mu (t) p(t)$``{=html},

then notice that by applying the product rule for differentiation that

we get:

``{=html}$\frac{d}{dt} \left ( \mu (t) y \right ) = \mu (t) \frac{dy}{dt} + \mu ' (t) y$``{=html}

and substituting

``{=html}$\mu ' (t) = \mu (t) y$``{=html}

and we get that

``{=html}$\frac{d}{dt} \left ( \mu (t) y \right ) = \mu (t) \frac{dy}{dt} + \mu (t) p(t) y$``{=html}

which is exactly the lefthand side of the equation above. Thus we get

that:







$\begin{aligned} \frac{d}{dt} \left ( \mu (t) y \right ) = \mu (t) g(t) \\\\ \end{aligned}$







The above differential equation can be solved by integrating both sides

of the equation with respect to

``{=html}$t$``{=html} and isolating

``{=html}$y$``{=html}. The question now

arises on how we can find such a function

``{=html}$\mu (t)$``{=html}.



> 

>

> ``{=html}Definition:``{=html} If

> ``{=html}$\frac{dy}{dt} + p(t) y = g(t)$``{=html}

> is a first order differential equation, then

> ``{=html}$\mu (t)$``{=html} is called

> an ``{=html}Integrating Factor``{=html} if for

> ``{=html}$\mu (t) \frac{dy}{dt} + \mu (t) p(t) y = \mu (t) g(t)$``{=html}

> we have that

> ``{=html}$\mu ' (t) = \mu (t) p(t)$``{=html}.

>

> 



The following proposition will give us a formula for obtaining the

integrating factor for differential equations in the form

``{=html}$\frac{dy}{dt} + p(t) y = g(t)$``{=html}.



> 

>

> ``{=html}Proposition 1:``{=html} If

> ``{=html}$\frac{dy}{dt} + p(t) y = g(t)$``{=html}

> is a differential equation, then an integrating factor

> ``{=html}$\mu (t)$``{=html} of this

> equation is given by the formula

> ``{=html}$\mu (t) = e^{\int p(t) \; dt}$``{=html}.

>

> 



\begin{itemize}
    \item ``{=html}Proof:``{=html} We want to find
\end{itemize}
    ``{=html}$\mu (t)$``{=html} such

    that

    ``{=html}$\mu ' (t) = \mu (t) p(t)$``{=html}.

    We can rewrite this equation as as

    ``{=html}$\frac{\mu '(t)}{\mu (t)} = p(t)$``{=html}

    and then:







$\begin{aligned} \quad \frac{ \mu '(t)}{ \mu (t)} = p(t) \\\\ \quad \frac{d}{dt} \ln \mid \mu (t) \mid = p(t) \\\\ \quad \int \frac{d}{dt} \ln \mid \mu (t) \mid \; dt = \int p(t) \; dt \\\\ \quad \ln \mid \mu (t) \mid = \int p(t) \; dt \\\\ \quad \mu (t) = \pm e^{\int p(t) \; dt} \end{aligned}$







\begin{itemize}
    \item Since we only need one integrating factor to solve differential
\end{itemize}
    equations in the form

    ``{=html}$\frac{dy}{dt} + p(t) y = g(t)$``{=html},

    we can more generally note that

    ``{=html}$\mu (t) = e^{\int p(t) \; dt}$``{=html}

    is an integrating factor of this differential equation.

    ``{=html}$\blacksquare$``{=html}



``{=html}Notice that from proposition 1 that integrating factors

``{=html}$\mu (t)$``{=html} are not

unique. In fact, there are infinitely many integrating factors. This can

be see when evaluating the indefinite integral in

``{=html}$\mu (t) = e^{\int p(t) \; dt}$``{=html}

which will result in getting

``{=html}$\mu (t) = e^{P(t) + C}$``{=html}

where ``{=html}$P$``{=html} is any

antiderivative if ``{=html}$p$``{=html}

and where ``{=html}$C$``{=html} is a

constant. We will always use the simplest integrating factor in solving

differential equations of this type.``{=html}



Let\'s now look at some examples of applying the method of integrating

factors.



\paragraph{Example 1}



``{=html}Find all solutions to the differential equation

``{=html}$\frac{dy}{dt} + \frac{2y}{t} = \frac{\sin t}{t^2}$``{=html}.``{=html}



We first notice that our differential equation is in the appropriate

form

``{=html}$\frac{dy}{dt} + p(t) y = g(t)$``{=html}

where

``{=html}$p(t) = \frac{2}{t}$``{=html}

and

``{=html}$g(t) = \frac{\sin t}{t^2}$``{=html}.

We compute our integrating factor as:







$\begin{aligned} \quad \mu (t) = e^{\int p(t) \; dt} = e^{\int \frac{2}{t} \; dt} = e^{2 \ln (t)} = e^{\ln (t^2)} = t^2 \end{aligned}$







Thus we have that for

``{=html}$C$``{=html} as a constant:







$\begin{aligned} \quad \mu (t) \frac{dy}{dt} + \mu (t) p(t) y = \mu (t) g(t) \\\\ \quad t^2 \frac{dy}{dt} + 2t y = \sin t \\\\ \quad \frac{d}{dt} \left ( t^2 y \right ) = \sin t \\\\ \quad \int \frac{d}{dt} \left ( t^2 y \right ) \; dt = \int \sin t \; dt \\\\ \quad t^2 y = -\cos t + C \\\\ \quad y = \frac{-\cos t}{t^2} + \frac{C}{t^2} \end{aligned}$







\paragraph{Example 2}



``{=html}Find all solutions to the differential equation

``{=html}$\frac{dy}{dt} - \frac{y}{t} + te^{-t} = 0$``{=html}.``{=html}



We first rewrite our differential equation as

``{=html}$\frac{dy}{dt} - \frac{y}{t} = - te^{-t}$``{=html}.

We note that in this form we have

``{=html}$p(t) = - \frac{1}{t}$``{=html}

and

``{=html}$g(t) = -t e^{-t}$``{=html}.

We now find an integrating factor:







$\begin{aligned} \quad \mu (t) = e^{\int p(t) \; dt} = e^{\int -\frac{1}{t} \; dt} = e^{-\ln t} = e^{\ln \left ( \frac{1}{t} \right)} = \frac{1}{t} \end{aligned}$







Thus we have that for

``{=html}$C$``{=html} as a constant:







$\begin{aligned} \quad \mu (t) \frac{dy}{dt} + \mu (t) p(t) y = \mu (t) g(t) \\\\ \quad \frac{1}{t} \frac{dy}{dt} - \frac{1}{t^2} y = -e^{-t} \\\\ \quad \frac{d}{dt} \left ( \frac{y}{t} \right ) = -e^{-t} \\\\ \quad \int \frac{d}{dt} \left ( \frac{y}{t} \right ) \; dt = -\int e^{-t} \; dt \\\\ \quad \frac{y}{t} = e^{-t} + C \\\\ \quad y = te^{-t} + tC \end{aligned}$







\paragraph{Licensing}



Content obtained and/or adapted from:



\begin{itemize}
    \item \href{http://mathonline.wikidot.com/the-method-of-integrating-factors}{The Method of Integrating Factors,
\end{itemize}
    mathonline.wikidot.com}

    under a CC BY-SA license"

%%===============================================================
\section{Bernoulli Equations (1st Order)}
\label{sec:bernoulli-equations-1st-order}
%%===============================================================

Bernoulli Equations (1st Order)

%%===============================================================
\section{Exact Differential Equations (1st Order)}
\label{sec:exact-differential-equations-1st-order}
%%===============================================================

Definition



Given a simply connected and open subset \textit{D} of \textbf{R}^2^ and two

functions \textit{I} and \textit{J} which are continuous on \textit{D}, an implicit

first-order ordinary differential equation of the form



:   $I(x, y)\, dx + J(x, y)\, dy = 0,$



is called an \textbf{exact differential equation} if there exists a

continuously differentiable function \textit{F}, called the **potential

function**, so that



$$\frac{\partial F}{\partial x} = I$$ and



$$\frac{\partial F}{\partial y} = J.$$



An exact equation may also be presented in the following form:



$$I(x, y) + J(x, y) \, y'(x) = 0$$ where the same constraints on \textit{I} and

\textit{J} apply for the differential equation to be exact.



The nomenclature of \""exact differential equation\"" refers to the exact

differential of a function. For a function

$F(x_0, x_1,...,x_{n-1},x_n)$, the exact or total derivative with

respect to $x_0$ is given by



$$\frac{dF}{dx_0}=\frac{\partial F}{\partial x_0}+\sum_{i=1}^{n}\frac{\partial F}{\partial x_i}\frac{dx_i}{dx_0}.$$



\paragraph{Example}



The function $F:\mathbb{R}^{2}\to\mathbb{R}$ given by



$$F(x,y) = \frac{1}{2}(x^2 + y^2)+c$$



is a potential function for the differential equation



$$x\,dx + y\,dy = 0.\,$$



\paragraph{Existence of potential functions}



In physical applications the functions \textit{I} and \textit{J} are usually not only

continuous but even continuously differentiable. Schwarz\'s Theorem then

provides us with a necessary criterion for the existence of a potential

function. For differential equations defined on simply connected sets

the criterion is even sufficient and we get the following theorem:



Given a differential equation of the form (for example, when \textit{F} has

zero slope in the x and y direction at \textit{F}(\textit{x},\textit{y})):



:   $I(x, y)\, dx + J(x, y)\, dy = 0,$



with \textit{I} and \textit{J} continuously differentiable on a simply connected and

open subset \textit{D} of \textbf{R}^2^ then a potential function \textit{F} exists if and

only if



$$\frac{\partial I}{\partial y}(x, y) = \frac{\partial J}{\partial x}(x, y).$$



\paragraph{Solutions to exact differential equations}



Given an exact differential equation defined on some simply connected

and open subset \textit{D} of \textbf{R}^2^ with potential function \textit{F}, a

differentiable function \textit{f} with (\textit{x}, \textit{f}(\textit{x})) in \textit{D} is a solution if

and only if there exists real number \textit{c} so that



$$F(x, f(x)) = c.$$



For an initial value problem



$$y(x_0) = y_0$$ we can locally find a potential function by



$$\begin{aligned}

F(x,y) &= \int_{x_0}^x I(t,y_0) dt + \int_{y_0}^y J(x,t) dt \\\&= \int_{x_0}^x I(t,y_0) dt + \int_{y_0}^y \left[ J(x_0,t) +  \int_{x_0}^{x} \frac{\partial I}{\partial y}(u, t)\, du \right]dt.

\end{aligned}$$



Solving



$$F(x,y) = c$$ for \textit{y}, where \textit{c} is a real number, we can then

construct all solutions.



\paragraph{Second order exact differential equations}



The concept of exact differential equations can be extended to second

order equations. Consider starting with the first-order exact equation:



$$I\left(x,y\right)+J\left(x,y\right){dy \over dx}=0$$



Since both functions $I\left(x,y\right)$, $J\left(x,y\right)$ are

functions of two variables, implicitly differentiating the multivariate

function yields



$${dI \over dx} +\left({ dJ\over dx}\right){dy \over dx}+{d^2y \over dx^2}\left(J\left(x,y\right)\right)=0$$



Expanding the total derivatives gives that



$${dI \over dx}={\partial I\over\partial x}+{\partial I\over\partial y}{dy \over dx}$$



and that



$${dJ \over dx}={\partial J\over\partial x}+{\partial J\over\partial y}{dy \over dx}$$



Combining the ${dy \over dx}$ terms gives



$${\partial I\over\partial x}+{dy \over dx}\left({\partial I\over\partial y}+{\partial J\over\partial x}+{\partial J\over\partial y}{dy \over dx}\right)+{d^2y \over dx^2}\left(J\left(x,y\right)\right)=0$$



If the equation is exact, then

${\partial J\over\partial x}={\partial I\over\partial y}$. Additionally,

the total derivative of $J\left(x,y\right)$ is equal to its implicit

ordinary derivative ${dJ \over dx}$. This leads to the rewritten

equation



$${\partial I\over\partial x}+{dy \over dx}\left({\partial J\over\partial x}+{dJ \over dx}\right)+{d^2y \over dx^2}\left(J\left(x,y\right)\right)=0$$



Now, let there be some second-order differential equation



$$f\left(x,y\right)+g\left(x,y,{dy \over dx}\right){dy \over dx}+{d^2y \over dx^2}\left(J\left(x,y\right)\right)=0$$



If ${\partial J\over\partial x}={\partial I\over\partial y}$ for exact

differential equations, then



$$\int \left({\partial I\over\partial y}\right)dy=\int \left({\partial J\over\partial x}\right)dy$$



and



$$\int \left({\partial I\over\partial y}\right)dy=\int \left({\partial J\over\partial x}\right)dy=I\left(x,y\right)-h\left(x\right)$$



where $h\left(x\right)$ is some arbitrary function only of $x$ that was

differentiated away to zero upon taking the partial derivative of

$I\left(x,y\right)$ with respect to $y$. Although the sign on

$h\left(x\right)$ could be positive, it is more intuitive to think of

the integral\'s result as $I\left(x,y\right)$ that is missing some

original extra function $h\left(x\right)$ that was partially

differentiated to zero.



Next, if



$${dI\over dx}={\partial I\over\partial x}+{\partial I\over\partial y}{dy \over dx}$$



then the term ${\partial I\over\partial x}$ should be a function only of

$x$ and $y$, since partial differentiation with respect to $x$ will hold

$y$ constant and not produce any derivatives of $y$. In the second order

equation



$$f\left(x,y\right)+g\left(x,y,{dy \over dx}\right){dy \over dx}+{d^2y \over dx^2}\left(J\left(x,y\right)\right)=0$$



only the term $f\left(x,y\right)$ is a term purely of $x$ and $y$. Let

${\partial I\over\partial x}=f\left(x,y\right)$. If

${\partial I\over\partial x}=f\left(x,y\right)$, then



$$f\left(x,y\right)={ dI\over dx}-{\partial I\over\partial y}{dy \over dx}$$



Since the total derivative of $I\left(x,y\right)$ with respect to $x$ is

equivalent to the implicit ordinary derivative ${dI \over dx}$ , then



$$f\left(x,y\right)+{\partial I\over\partial y}{dy \over dx}={dI \over dx}={d \over dx}\left(I\left(x,y\right)-h\left(x\right)\right)+{dh\left(x\right) \over dx}$$



So,



$${dh\left(x\right) \over dx}=f\left(x,y\right)+{\partial I\over\partial y}{dy \over dx}-{d \over dx}\left(I\left(x,y\right)-h\left(x\right)\right)$$



and



$$h\left(x\right)=\int\left(f\left(x,y\right)+{\partial I\over\partial y}{dy \over dx}-{d \over dx}\left(I\left(x,y\right)-h\left(x\right)\right)\right)dx$$



Thus, the second order differential equation



$$f\left(x,y\right)+g\left(x,y,{dy \over dx}\right){dy \over dx}+{d^2y \over dx^2}\left(J\left(x,y\right)\right)=0$$



is exact only if

$g\left(x,y,{dy \over dx}\right)={ dJ\over dx}+{\partial J\over\partial x}={dJ \over dx}+{\partial J\over\partial x}$

and only if the below expression



$$\int\left(f\left(x,y\right)+{\partial I\over\partial y}{dy \over dx}-{d \over dx}\left(I\left(x,y\right)-h\left(x\right)\right)\right)dx=\int \left(f\left(x,y\right)-{\partial \left(I\left(x,y\right)-h\left(x\right)\right)\over\partial x}\right)dx$$



is a function solely of $x$. Once $h\left(x\right)$ is calculated with

its arbitrary constant, it is added to

$I\left(x,y\right)-h\left(x\right)$ to make $I\left(x,y\right)$. If the

equation is exact, then we can reduce to the first order exact form

which is solvable by the usual method for first-order exact equations.



$$I\left(x,y\right)+J\left(x,y\right){dy \over dx}=0$$



Now, however, in the final implicit solution there will be a $C_1x$ term

from integration of $h\left(x\right)$ with respect to $x$ twice as well

as a $C_2$, two arbitrary constants as expected from a second-order

equation.



\paragraph{Example}



Given the differential equation



$$\left(1-x^2\right)y''-4xy'-2y=0$$



one can always easily check for exactness by examining the $y''$ term.

In this case, both the partial and total derivative of $1-x^2$ with

respect to $x$ are $-2x$, so their sum is $-4x$, which is exactly the

term in front of $y'$. With one of the conditions for exactness met, one

can calculate that



$$\int \left(-2x\right)dy=I\left(x,y\right)-h\left(x\right)=-2xy$$



Letting $f\left(x,y\right)=-2y$, then



$$\int \left(-2y-2xy'-{d \over dx}\left(-2xy \right)\right)dx=\int \left(-2y-2xy'+2xy'+2y\right)dx=\int \left(0\right)dx=h\left(x\right)$$



So, $h\left(x\right)$ is indeed a function only of $x$ and the second

order differential equation is exact. Therefore, $h\left(x\right)=C_1$

and $I\left(x,y\right)=-2xy+C_1$. Reduction to a first-order exact

equation yields



$$-2xy+C_1+\left(1-x^2\right)y'=0$$



Integrating $I\left(x,y\right)$ with respect to $x$ yields



$$-x^2y+C_1x+i\left(y\right)=0$$



where $i\left(y\right)$ is some arbitrary function of $y$.

Differentiating with respect to $y$ gives an equation correlating the

derivative and the $y'$ term.



$$-x^2+i'\left(y\right)=1-x^2$$



So, $i\left(y\right)=y+C_2$ and the full implicit solution becomes



$$C_1x+C_2+y-x^2y=0$$



Solving explicitly for $y$ yields



$$y= \frac{C_1x+C_2}{1-x^2}$$



\paragraph{Higher order exact differential equations}



The concepts of exact differential equations can be extended to any

order. Starting with the exact second order equation



$${d^2y \over dx^2}\left(J\left(x,y\right)\right)+{dy \over dx}\left({dJ \over dx}+{\partial J\over\partial x}\right)+f\left(x,y\right)=0$$



it was previously shown that equation is defined such that



$$f\left(x,y\right)={dh\left(x\right) \over dx}+{d \over dx}\left(I\left(x,y\right)-h\left(x\right)\right)-{\partial J\over\partial x}{dy \over dx}$$



Implicit differentiation of the exact second-order equation $n$ times

will yield an $\left(n+2\right)$th order differential equation with new

conditions for exactness that can be readily deduced from the form of

the equation produced. For example, differentiating the above

second-order differential equation once to yield a third-order exact

equation gives the following form



$${d^3y \over dx^3}\left(J\left(x,y\right)\right)+{d^2y \over dx^2}{dJ \over dx}+{d^2y \over dx^2}\left({dJ \over dx}+{\partial J\over\partial x}\right)+{dy \over dx}\left({d^2J \over dx^2}+{d \over dx}\left({\partial J\over\partial x}\right)\right)+{df\left(x,y\right) \over dx}=0$$



where



$${df\left(x,y\right) \over dx}={d^2h\left(x\right) \over dx^2}+{d^2 \over dx^2}\left(I\left(x,y\right)-h\left(x\right)\right)-{d^2y \over dx^2}{\partial J\over\partial x}-{dy \over dx}{d \over dx}\left({\partial J\over\partial x}\right)=F\left(x,y,{dy \over dx}\right)$$

and where $F\left(x,y,{dy \over dx}\right)$ is a function only of $x,y$

and ${dy \over dx}$. Combining all ${dy \over dx}$ and

${d^2y \over dx^2}$ terms not coming from

$F\left(x,y,{dy \over dx}\right)$ gives



$${d^3y \over dx^3}\left(J\left(x,y\right)\right)+{d^2y \over dx^2}\left(2{dJ \over dx}+{\partial J\over\partial x}\right)+{dy \over dx}\left({d^2J \over dx^2}+{d \over dx}\left({\partial J\over\partial x}\right)\right)+F\left(x,y,{dy \over dx}\right)=0$$



Thus, the three conditions for exactness for a third-order differential

equation are: the ${d^2y \over dx^2}$ term must be

$2{dJ \over dx}+{\partial J\over\partial x}$, the ${dy \over dx}$ term

must be

${d^2J \over dx^2}+{d \over dx}\left({\partial J\over\partial x}\right)$

and



$$F\left(x,y,{dy \over dx}\right)-{d^2 \over dx^2}\left(I\left(x,y\right)-h\left(x\right)\right)+{d^2y \over dx^2}{\partial J\over\partial x}+{dy \over dx}{d \over dx}\left({\partial J\over\partial x}\right)$$



must be a function solely of $x$.



\paragraph{Example}



Consider the nonlinear third-order differential equation



$$yy'''+3y'y''+12x^2=0$$



If $J\left(x,y\right)=y$, then

$y''\left(2{dJ \over dx}+{\partial J\over\partial x}\right)$ is $2y'y''$

and

$y'\left({d^2J \over dx^2}+{d \over dx}\left({\partial J\over\partial x}\right)\right)=y'y''$which

together sum to $3y'y''$. Fortunately, this appears in our equation. For

the last condition of exactness,



$$F\left(x,y,{dy \over dx}\right)-{d^2 \over dx^2}\left(I\left(x,y\right)-h\left(x\right)\right)+{d^2y \over dx^2}{\partial J\over\partial x}+{dy \over dx}{d \over dx}\left({\partial J\over\partial x}\right)=12x^2-0+0+0=12x^2$$



which is indeed a function only of $x$. So, the differential equation is

exact. Integrating twice yields that

$h\left(x\right)=x^4+C_1x+C_2=I\left(x,y\right)$. Rewriting the equation

as a first-order exact differential equation yields



$$x^4+C_1x+C_2+yy'=0$$



Integrating $I\left(x,y\right)$ with respect to $x$ gives that

${x^5\over 5}+C_1x^2+C_2x+i\left(y\right)=0$. Differentiating with

respect to $y$ and equating that to the term in front of $y'$ in the

first-order equation gives that $i'\left(y\right)=y$ and that

$i\left(y\right)={y^2\over 2}+C_3$. The full implicit solution becomes



$${x^5\over 5}+C_1x^2+C_2x+C_3+{y^2\over 2}=0$$



The explicit solution, then, is



$$y=\pm\sqrt{C_1x^2+C_2x+C_3-\frac{2x^5}{5}}$$



\paragraph{Licensing}



Content obtained and/or adapted from:



\begin{itemize}
    \item \href{https://en.wikipedia.org/wiki/Exact_differential_equation}{Exact differential equation,
\end{itemize}
    Wikipedia}

    under a CC BY-SA license"

%%===============================================================
\section{Linear Independence of Functions}
\label{sec:linear-independence-of-functions}
%%===============================================================

"If we have an $n^{\mathrm{th}}$ order linear homogenous differential

equation

$\frac{d^ny}{dt^n} + p_1(t) \frac{d^{n-1}y}{dt^{n-1}} + ... + p_{n-1}(t) \frac{dy}{dt} + p_n(t)y = 0$

where $p_1$, $p_2$, ..., $p_n$ are continuous on an open interval $I$

and if $y = y_1(t)$, $y = y_2(t)$, ..., $y = y_n(t)$ are solutions to

this differential equation, then provided that

$W(y_1, y_2, ..., y_n) \neq 0$ for at least one point $t \in I$, then

$y_1$, $y_2$, ..., $y_n$ form a fundamental set of solutions to this

differential equation - that is, for constants $C_1$, $C_2$, ..., $C_n$,

then every solution to this differential equation can be written in the

form:



$\begin{aligned} \quad y = C_1y_1(t) + C_2y_2(t) + ... + C_ny_n(t) \end{aligned}$



We will now look at the connection between the solutions $y_1$, $y_2$,

..., $y_n$ forming a fundamental set of solutions and the linear

independence/dependence of such solutions. We first define linear

independence and linear dependence below.







``{=html}Definition:``{=html} The functions $f_1$,

$f_2$, ..., $f_n$ are said to be ``{=html}Linearly

Independent``{=html} on an interval $I$ if for constants $k_1$,

$k_2$, ..., $k_n$ we have that

$k_1f_1(t) + k_2f_2(t) + ... + k_nf_n(t) = 0$ implies that

$k_1 = k_2 = ... = k_n = 0$ for all $t \in I$. This set of functions is

said to be ``{=html}Linearly Dependent``{=html} if

$k_1f_1(t) + k_2f_2(t) + ... + k_nf_n(t) = 0$ where $k_1$, $k_2$, ...,

$k_n$ are not all zero for all $t \in I$.







Perhaps the simplest linearly independent sets of functions is that set

that contains $f_1(t) = 1$, $f_2(t) = t$, and $f_3(t) = t^2$. Let $k_1$,

$k_2$, and $k_3$ be constants and consider the following equation:



$\begin{aligned} \quad k_1f_1(t) + k_2f_2(t) + k_3f_3(t) = 0 \\\\ \quad k_1 + k_2t + k_3t^2 = 0 \end{aligned}$



It\'s not hard to see that equation above is satisfied if and only if

the constants $k_1 = k_2 = k_3 = 0$.



For another example, consider the functions $f_1(t) = \sin t$ and

$f_2(t) = \sin (t + \pi)$ defined on all of $\mathbb{R}$. This set of

functions is not linearly independent. To show this, let $k_1$ and $k_2$

be constants and consider the following equation:



$\begin{aligned} \quad k_1f_1(t) + k_2f_2(t) = 0 \\\\ \quad k_1 \sin t + k_2 \sin (t + \pi) = 0 \end{aligned}$



Now choose $t = \pi$. Then we have that:



$\begin{aligned} \quad k_1 \sin \pi + k_2 \sin 2\pi = 0 \\\\ \end{aligned}$



But the above equation is true for any choice of constants $k_1$ and

$k_2$ since $\sin \pi = \sin 2\pi = 0$, and thus $f_1$ and $f_2$ do not

form a linearly independent set on all of $\mathbb{R}$.



From the concept of linear independence/dependence, we obtain the

following theorem on fundamental sets of solutions for $n^{\mathrm{th}}$

order linear homogenous differential equations.







``{=html}Theorem 1:``{=html} Let

$\frac{d^ny}{dt^n} + p_1(t) \frac{d^{n-1}y}{dt^{n-1}} + ... + p_{n-1}(t) \frac{dy}{dt} + p_n(t)y = 0$

be an $n^{\mathrm{th}}$ order linear homogenous differential equation.

If $y = y_1(t)$, $y = y_2(t)$, ..., $y = y_n(t)$ are solutions to this

differential equation then $y_1$, $y_2$, ..., $y_n$ form a fundamental

set of solutions to this differential equation on the open interval $I$

if and only if $y_1$, $y_2$, ..., $y_n$ are linearly dependent on $I$.







\begin{itemize}
    \item ``{=html}Proof:``{=html} Consider the following
\end{itemize}
    $n^{\mathrm{th}}$ order linear homogenous differential equation:



$\begin{aligned} \quad \frac{d^ny}{dt^n} + p_1(t) \frac{d^{n-1}y}{dt^{n-1}} + ... + p_{n-1}(t) \frac{dy}{dt} + p_n(t)y = 0 \end{aligned}$



\begin{itemize}
    \item $\Rightarrow$ Suppose that $y = y_1(t)$, $y = y_2(t)$, ...,
\end{itemize}
    $y = y_n(t)$ form a fundamental set of solutions to this

    differential equation on the open interval $I$. Then this implies

    that for all $t_0 \in I$ we have that :



$\begin{aligned} \quad W(y_1, y_2, ..., y_n)  \bigg|_{t_0} \neq 0 \end{aligned}$



\begin{itemize}
    \item Thus this implies that following system of equations have only
\end{itemize}
    trivial solution $k_1 = k_2 = ... = k_n = 0$:



$\begin{aligned} \quad k_1y_1(t) + k_2y_2(t) + ... + k_ny_n(t) = 0 \\\\ \quad k_1y_1'(t) + k_2y_2'(t) + ... + k_ny_n'(t) = 0 \\\\ \quad \quad \quad \quad \quad \quad \vdots \quad \quad \quad \quad \quad \quad \\\\ \quad k_1y_1^{(n-2)}(t) + k_2y_2^{(n-2)}(t) + ... + k_ny_n^{(n-2)}(t) = 0 \\\\ \quad k_1y_1^{(n-1)}(t) + k_2y_2^{(n-1)}(t) + ... + k_ny_n^{(n-1)}(t) = 0 \end{aligned}$



\begin{itemize}
    \item Thus the equation $k_1y_1(t) + k_2y_2(t) + ... + k_ny_n(t) = 0$
\end{itemize}
    implies that $k_1 = k_2 = ... = k_n = 0$. Thus $y_1$, $y_2$, ...,

    $y_n$ are linearly independent on $I$.



<!-- -->



\begin{itemize}
    \item $\Leftarrow$ We will prove the converse of Theorem 1 by
\end{itemize}
    contradiction. Suppose that $y_1$, $y_2$, ..., $y_n$ are linearly

    independent on $I$, and assume that instead $y_1$, $y_2$, ..., $y_n$

    do NOT form a fundamental set of solutions on $I$. Then for some

    $t_0 \in I$, the Wronskian

    $W(y_1, y_2, ..., y_n)  \bigg|_{t_0} = 0$. Thus the system of

    equations above does not have only the trivial solution. Let the

    constants $k_1^*$, $k_2^*$, ..., $k_n^*$ be a nontrivial solution to

    this system. Define $\phi(t)$ as:



$\begin{aligned} \quad \phi(t) = k_1^* y_1(t) + k_2^* y_2(t) + ... + k_n^* y_n(t) \end{aligned}$



\begin{itemize}
    \item Note that $y = \phi(t)$ satisfies the initial conditions
\end{itemize}
    $y(t_0) = 0$, $y'(t_0) = 0$, ..., $y^{(n-1)} (t_0) = 0$, and

    $\phi(t)$ satisfies our $n^{\mathrm{th}}$ order linear homogenous

    differential equation because $\phi(t)$ is a linear combination of

    the solutions $y_1$, $y_2$, ..., $y_n$.



<!-- -->



\begin{itemize}
    \item Now note that the function $y = 0$ also satisfies the differential
\end{itemize}
    equation and the initial conditions. By the existence/uniqueness

    theorem for $n^{\mathrm{th}}$ order linear homogenous differential

    equations, this implies that $\phi(t) = 0$ for all $t \in I$, so:



$\begin{aligned} \quad 0 = k_1^* y_1(t) + k_2^* y_2(t) + ... + k_n^* y_n(t) \end{aligned}$



\begin{itemize}
    \item But $y_1$, $y_2$, ..., $y_n$ are linearly independent which implies
\end{itemize}
    that $k_1^* = k_2^* = ... = k_n^* = 0$. Thus $k_1^*$, $k_2^*$, ...,

    $k_n^*$ is a trivial solution to the system above, which is a

    contradiction. Therefore our assumption that $y_1$, $y_2$, ...,

    $y_n$ do not form a fundamental set of solutions was false.

    $\blacksquare$



\paragraph{Licensing}



Content obtained and/or adapted from:



\begin{itemize}
    \item \href{http://mathonline.wikidot.com/linear-independence-dependence-of-a-set-of-functions}{Linear Independence Dependence of a Set of Functions,
\end{itemize}
    }

    under a CC BY-SA license"

%%===============================================================
\section{Linear Dependence of Functions}
\label{sec:linear-dependence-of-functions}
%%===============================================================

"If we have an $n^{\mathrm{th}}$ order linear homogenous differential

equation

$\frac{d^ny}{dt^n} + p_1(t) \frac{d^{n-1}y}{dt^{n-1}} + ... + p_{n-1}(t) \frac{dy}{dt} + p_n(t)y = 0$

where $p_1$, $p_2$, ..., $p_n$ are continuous on an open interval $I$

and if $y = y_1(t)$, $y = y_2(t)$, ..., $y = y_n(t)$ are solutions to

this differential equation, then provided that

$W(y_1, y_2, ..., y_n) \neq 0$ for at least one point $t \in I$, then

$y_1$, $y_2$, ..., $y_n$ form a fundamental set of solutions to this

differential equation - that is, for constants $C_1$, $C_2$, ..., $C_n$,

then every solution to this differential equation can be written in the

form:



$\begin{aligned} \quad y = C_1y_1(t) + C_2y_2(t) + ... + C_ny_n(t) \end{aligned}$



We will now look at the connection between the solutions $y_1$, $y_2$,

..., $y_n$ forming a fundamental set of solutions and the linear

independence/dependence of such solutions. We first define linear

independence and linear dependence below.







``{=html}Definition:``{=html} The functions $f_1$,

$f_2$, ..., $f_n$ are said to be ``{=html}Linearly

Independent``{=html} on an interval $I$ if for constants $k_1$,

$k_2$, ..., $k_n$ we have that

$k_1f_1(t) + k_2f_2(t) + ... + k_nf_n(t) = 0$ implies that

$k_1 = k_2 = ... = k_n = 0$ for all $t \in I$. This set of functions is

said to be ``{=html}Linearly Dependent``{=html} if

$k_1f_1(t) + k_2f_2(t) + ... + k_nf_n(t) = 0$ where $k_1$, $k_2$, ...,

$k_n$ are not all zero for all $t \in I$.







Perhaps the simplest linearly independent sets of functions is that set

that contains $f_1(t) = 1$, $f_2(t) = t$, and $f_3(t) = t^2$. Let $k_1$,

$k_2$, and $k_3$ be constants and consider the following equation:



$\begin{aligned} \quad k_1f_1(t) + k_2f_2(t) + k_3f_3(t) = 0 \\\\ \quad k_1 + k_2t + k_3t^2 = 0 \end{aligned}$



It\'s not hard to see that equation above is satisfied if and only if

the constants $k_1 = k_2 = k_3 = 0$.



For another example, consider the functions $f_1(t) = \sin t$ and

$f_2(t) = \sin (t + \pi)$ defined on all of $\mathbb{R}$. This set of

functions is not linearly independent. To show this, let $k_1$ and $k_2$

be constants and consider the following equation:



$\begin{aligned} \quad k_1f_1(t) + k_2f_2(t) = 0 \\\\ \quad k_1 \sin t + k_2 \sin (t + \pi) = 0 \end{aligned}$



Now choose $t = \pi$. Then we have that:



$\begin{aligned} \quad k_1 \sin \pi + k_2 \sin 2\pi = 0 \\\\ \end{aligned}$



But the above equation is true for any choice of constants $k_1$ and

$k_2$ since $\sin \pi = \sin 2\pi = 0$, and thus $f_1$ and $f_2$ do not

form a linearly independent set on all of $\mathbb{R}$.



From the concept of linear independence/dependence, we obtain the

following theorem on fundamental sets of solutions for $n^{\mathrm{th}}$

order linear homogenous differential equations.







``{=html}Theorem 1:``{=html} Let

$\frac{d^ny}{dt^n} + p_1(t) \frac{d^{n-1}y}{dt^{n-1}} + ... + p_{n-1}(t) \frac{dy}{dt} + p_n(t)y = 0$

be an $n^{\mathrm{th}}$ order linear homogenous differential equation.

If $y = y_1(t)$, $y = y_2(t)$, ..., $y = y_n(t)$ are solutions to this

differential equation then $y_1$, $y_2$, ..., $y_n$ form a fundamental

set of solutions to this differential equation on the open interval $I$

if and only if $y_1$, $y_2$, ..., $y_n$ are linearly dependent on $I$.







\begin{itemize}
    \item ``{=html}Proof:``{=html} Consider the following
\end{itemize}
    $n^{\mathrm{th}}$ order linear homogenous differential equation:



$\begin{aligned} \quad \frac{d^ny}{dt^n} + p_1(t) \frac{d^{n-1}y}{dt^{n-1}} + ... + p_{n-1}(t) \frac{dy}{dt} + p_n(t)y = 0 \end{aligned}$



\begin{itemize}
    \item $\Rightarrow$ Suppose that $y = y_1(t)$, $y = y_2(t)$, ...,
\end{itemize}
    $y = y_n(t)$ form a fundamental set of solutions to this

    differential equation on the open interval $I$. Then this implies

    that for all $t_0 \in I$ we have that :



$\begin{aligned} \quad W(y_1, y_2, ..., y_n)  \bigg|_{t_0} \neq 0 \end{aligned}$



\begin{itemize}
    \item Thus this implies that following system of equations have only
\end{itemize}
    trivial solution $k_1 = k_2 = ... = k_n = 0$:



$\begin{aligned} \quad k_1y_1(t) + k_2y_2(t) + ... + k_ny_n(t) = 0 \\\\ \quad k_1y_1'(t) + k_2y_2'(t) + ... + k_ny_n'(t) = 0 \\\\ \quad \quad \quad \quad \quad \quad \vdots \quad \quad \quad \quad \quad \quad \\\\ \quad k_1y_1^{(n-2)}(t) + k_2y_2^{(n-2)}(t) + ... + k_ny_n^{(n-2)}(t) = 0 \\\\ \quad k_1y_1^{(n-1)}(t) + k_2y_2^{(n-1)}(t) + ... + k_ny_n^{(n-1)}(t) = 0 \end{aligned}$



\begin{itemize}
    \item Thus the equation $k_1y_1(t) + k_2y_2(t) + ... + k_ny_n(t) = 0$
\end{itemize}
    implies that $k_1 = k_2 = ... = k_n = 0$. Thus $y_1$, $y_2$, ...,

    $y_n$ are linearly independent on $I$.



<!-- -->



\begin{itemize}
    \item $\Leftarrow$ We will prove the converse of Theorem 1 by
\end{itemize}
    contradiction. Suppose that $y_1$, $y_2$, ..., $y_n$ are linearly

    independent on $I$, and assume that instead $y_1$, $y_2$, ..., $y_n$

    do NOT form a fundamental set of solutions on $I$. Then for some

    $t_0 \in I$, the Wronskian

    $W(y_1, y_2, ..., y_n)  \bigg|_{t_0} = 0$. Thus the system of

    equations above does not have only the trivial solution. Let the

    constants $k_1^*$, $k_2^*$, ..., $k_n^*$ be a nontrivial solution to

    this system. Define $\phi(t)$ as:



$\begin{aligned} \quad \phi(t) = k_1^* y_1(t) + k_2^* y_2(t) + ... + k_n^* y_n(t) \end{aligned}$



\begin{itemize}
    \item Note that $y = \phi(t)$ satisfies the initial conditions
\end{itemize}
    $y(t_0) = 0$, $y'(t_0) = 0$, ..., $y^{(n-1)} (t_0) = 0$, and

    $\phi(t)$ satisfies our $n^{\mathrm{th}}$ order linear homogenous

    differential equation because $\phi(t)$ is a linear combination of

    the solutions $y_1$, $y_2$, ..., $y_n$.



<!-- -->



\begin{itemize}
    \item Now note that the function $y = 0$ also satisfies the differential
\end{itemize}
    equation and the initial conditions. By the existence/uniqueness

    theorem for $n^{\mathrm{th}}$ order linear homogenous differential

    equations, this implies that $\phi(t) = 0$ for all $t \in I$, so:



$\begin{aligned} \quad 0 = k_1^* y_1(t) + k_2^* y_2(t) + ... + k_n^* y_n(t) \end{aligned}$



\begin{itemize}
    \item But $y_1$, $y_2$, ..., $y_n$ are linearly independent which implies
\end{itemize}
    that $k_1^* = k_2^* = ... = k_n^* = 0$. Thus $k_1^*$, $k_2^*$, ...,

    $k_n^*$ is a trivial solution to the system above, which is a

    contradiction. Therefore our assumption that $y_1$, $y_2$, ...,

    $y_n$ do not form a fundamental set of solutions was false.

    $\blacksquare$



\paragraph{Licensing}



Content obtained and/or adapted from:



\begin{itemize}
    \item \href{http://mathonline.wikidot.com/linear-independence-dependence-of-a-set-of-functions}{Linear Independence Dependence of a Set of Functions,
\end{itemize}
    }

    under a CC BY-SA license"

%%===============================================================
\section{Wronskian}
\label{sec:wronskian}
%%===============================================================

Wronskian



Let L(y)=0 be a homogeneous equation of degree n, and let

$y_1,y_2,...,y_n$ be linearly independent solutions of this equation.



The general solution is then $y=C_1y_1+C_2y_2+...+C_ny_n$.



Suppose instead that the solutions are not linearly independent. Then

there exists values for $C_1,C_2,...,C_n$ not all zero such that

$C_1y_1+C_2y_2+...+C_ny_n$=0.



Then the following equations also hold true:



$C_1y_1'+C_2y_2'+...+C_ny_n'=0$$C_1y_1''+C_2y_2''+...+C_ny_n''=0$\...$C_1y_1^{(n-1)}+C_2y_2^{(n-1)}+...+C_ny_n^{(n-1)}=0$



When there is to be a non-trivial solution to this system of homogeneous

linear equations, the column vectors corresponding to each coefficient

are linearly dependent. This is equivalent to saying that the following

determinant, called the \textbf{Wronskian} is 0:



$\begin{Vmatrix}

y_1 \& y_2 \& ... \& y_n \\\y_1' \& y_2' \& ... \& y_n' \\\y_1'' \& y_2'' \& ... \& y_n'' \\\... \& ... \& ... \& ... \\\y_1^{(n-1)} \& y_2^{(n-1)} \& ... \& y_n^{(n-1)} \\\\end{Vmatrix} = 0$



Thus, if the solutions are dependent, then their Wronskian is equal to

0, and conversely, and conversely if the Wronskian is equal to 0, then

the columns of that matrix must be linearly dependent, indicating that

the solutions are linearly dependent.



\paragraph{Resources}



\begin{itemize}
    \item \href{https://tutorial.math.lamar.edu/Classes/DE/Wronskian.aspx}{Notes on the
\end{itemize}
    Wronskian}.

    Produced by Paul Dawkins, Lamar University



\paragraph{Licensing}



Content obtained and/or adapted from:



\begin{itemize}
    \item \href{https://en.wikibooks.org/wiki/Ordinary_Differential_Equations/Linear_Equations}{Linear Equations, Wikibooks: Ordinary Differential
\end{itemize}
    Equations}

    under a CC BY-SA license"

%%===============================================================
\section{Reduction of the Order}
\label{sec:reduction-of-the-order}
%%===============================================================

"\textbf{Reduction of order} is a technique in mathematics for solving

second-order linear ordinary differential equations. It is employed when

one solution $y_1(x)$ is known and a second linearly independent

solution $y_2(x)$ is desired. The method also applies to \textit{n}-th order

equations. In this case the ansatz will yield an (\textit{n}-1)-th order

equation for $v$.



\paragraph{Second-order linear ordinary differential equations}



\paragraph{An example}



Consider the general, homogeneous, second-order linear constant

coefficient ordinary differential equation. (ODE)



$$a y''(x) + b y'(x) + c y(x) = 0,$$



where $a, b, c$ are real non-zero coefficients. Two linearly independent

solutions for this ODE can be straightforwardly found using

characteristic equations except for the case when the discriminant,

$b^2 - 4 a c$, vanishes. In this case,



$$a y''(x) + b y'(x) + \frac{b^2}{4a} y(x) = 0,$$



from which only one solution,



$$y_1(x) = e^{-\frac{b}{2a} x},$$



can be found using its characteristic equation.



The method of reduction of order is used to obtain a second linearly

independent solution to this differential equation using our one known

solution. To find a second solution we take as a guess



$$y_2(x) = v(x) y_1(x)$$



where $v(x)$ is an unknown function to be determined. Since $y_2(x)$

must satisfy the original ODE, we substitute it back in to get



$$a \left( v'' y_1 + 2 v' y_1' + v y_1'' \right) + b \left( v' y_1 + v y_1' \right) + \frac{b^2}{4a} v y_1 = 0.$$



Rearranging this equation in terms of the derivatives of $v(x)$ we get



$$\left(a y_1 \right) v'' + \left( 2 a y_1' + b y_1 \right) v' + \left( a y_1'' + b y_1' + \frac{b^2}{4a} y_1 \right) v = 0.$$



Since we know that $y_1(x)$ is a solution to the original problem, the

coefficient of the last term is equal to zero. Furthermore, substituting

$y_1(x)$ into the second term\'s coefficient yields (for that

coefficient)



$$2 a \left( - \frac{b}{2a} e^{-\frac{b}{2a} x} \right) + b e^{-\frac{b}{2a} x} = \left( -b + b \right) e^{-\frac{b}{2a} x} = 0.$$



Therefore, we are left with



$$a y_1 v'' = 0.$$



Since $a$ is assumed non-zero and $y_1(x)$ is an exponential function

(and thus always non-zero), we have



$$v'' = 0.$$



This can be integrated twice to yield



$$v(x) = c_1 x + c_2$$



where $c_1, c_2$ are constants of integration. We now can write our

second solution as



$$y_2(x) = ( c_1 x + c_2 ) y_1(x) = c_1 x y_1(x) + c_2 y_1(x).$$



Since the second term in $y_2(x)$ is a scalar multiple of the first

solution (and thus linearly dependent) we can drop that term, yielding a

final solution of



$$y_2(x) = x y_1(x) = x e^{-\frac{b}{2 a} x}.$$



Finally, we can prove that the second solution $y_2(x)$ found via this

method is linearly independent of the first solution by calculating the

Wronskian



$$W(y_1,y_2)(x) = \begin{vmatrix} y_1 & x y_1 \\\\ y_1' & y_1 + x y_1' \end{vmatrix} = y_1 ( y_1 + x y_1' ) - x y_1 y_1' = y_1^{2} + x y_1 y_1' - x y_1 y_1' = y_1^{2} = e^{-\frac{b}{a}x} \neq 0.$$



Thus $y_2(x)$ is the second linearly independent solution we were

looking for.



\paragraph{General method}



Given the general non-homogeneous linear differential equation



$$y'' + p(t)y' + q(t)y = r(t)$$



and a single solution $y_1(t)$ of the homogeneous equation \[$r(t)=0$\],

let us try a solution of the full non-homogeneous equation in the form:



$$y_2 = v(t)y_1(t)$$



where $v(t)$ is an arbitrary function. Thus



$$y_2' = v'(t)y_1(t) + v(t)y_1'(t)$$



and



$$y_2'' = v''(t)y_1(t) + 2v'(t)y_1'(t) + v(t)y_1''(t).$$



If these are substituted for $y$, $y'$, and $y''$ in the differential

equation, then



$$y_1(t)\,v'' + (2y_1'(t)+p(t)y_1(t))\,v' + (y_1''(t)+p(t)y_1'(t)+q(t)y_1(t))\,v = r(t).$$



Since $y_1(t)$ is a solution of the original homogeneous differential

equation, $y_1''(t)+p(t)y_1'(t)+q(t)y_1(t)=0$, so we can reduce to



$$y_1(t)\,v'' + (2y_1'(t)+p(t)y_1(t))\,v' = r(t)$$



which is a first-order differential equation for $v'(t)$ (reduction of

order). Divide by $y_1(t)$, obtaining



$$v''+\left(\frac{2y_1'(t)}{y_1(t)}+p(t)\right)\,v'=\frac{r(t)}{y_1(t)}.$$



The integrating factor is

$\mu(t)=e^{\int(\frac{2y_1'(t)}{y_1(t)}+p(t))dt}=y_1^2(t)e^{\int p(t) dt}$.



Multiplying the differential equation by the integrating factor

$\mu(t)$, the equation for $v(t)$ can be reduced to



$$\frac{d}{dt}\left(v'(t) y_1^2(t) e^{\int p(t) dt}\right) = y_1(t)r(t)e^{\int p(t) dt}.$$



After integrating the last equation, $v'(t)$ is found, containing one

constant of integration. Then, integrate $v'(t)$ to find the full

solution of the original non-homogeneous second-order equation,

exhibiting two constants of integration as it should:



$$y_2(t) = v(t)y_1(t).$$



\paragraph{Resources}



\begin{itemize}
    \item \href{https://tutorial.math.lamar.edu/Classes/DE/ReductionofOrder.aspx}{Reduction of the Order
\end{itemize}
    Notes}.

    Produced by Paul Dawkins, Lamar University



\paragraph{Licensing}



Content obtained and/or adapted from:



\begin{itemize}
    \item \href{https://en.wikipedia.org/wiki/Reduction_of_order}{Reduction of order,
\end{itemize}
    Wikipedia} under a

    CC BY-SA license"

%%===============================================================
\section{Linear Homogeneous Equations}
\label{sec:linear-homogeneous-equations}
%%===============================================================

"A differential equation can be \textbf{homogeneous} in either of two

respects.



A first order differential equation is said to be homogeneous if it may

be written



$$f(x,y) \, dy = g(x,y) \, dx,$$ where `{{mvar|f}}`{=mediawiki} and

`{{mvar|g}}`{=mediawiki} are homogeneous functions of the same degree of

`{{mvar|x}}`{=mediawiki} and `{{mvar|y}}`{=mediawiki}. In this case, the

change of variable `{{math|1=''y'' = ''ux''}}`{=mediawiki} leads to an

equation of the form



$$\frac{dx}{x} = h(u) \, du,$$ which is easy to solve by integration of

the two members.



Otherwise, a differential equation is homogeneous if it is a homogeneous

function of the unknown function and its derivatives. In the case of

linear differential equations, this means that there are no constant

terms. The solutions of any linear ordinary differential equation of any

order may be deduced by integration from the solution of the homogeneous

equation obtained by removing the constant term.



\paragraph{Homogeneous first-order differential equations}



A first-order ordinary differential equation in the form:



$$M(x,y)\,dx + N(x,y)\,dy = 0$$



is a homogeneous type if both functions

`{{math|''M''(''x'', ''y'')}}`{=mediawiki} and

`{{math|''N''(''x'', ''y'')}}`{=mediawiki} are homogeneous functions of

the same degree `{{mvar|n}}`{=mediawiki}. That is, multiplying each

variable by a parameter `{{math|?}}`{=mediawiki}, we find



$$M(\lambda x, \lambda y) = \lambda^n M(x,y) \quad \text{and} \quad N(\lambda x, \lambda y) = \lambda^n N(x,y)\,.$$



Thus,



$$\frac{M(\lambda x, \lambda y)}{N(\lambda x, \lambda y)} = \frac{M(x,y)}{N(x,y)}\,.$$



\paragraph{Solution method}



In the quotient $\frac{M(tx,ty)}{N(tx,ty)} = \frac{M(x,y)}{N(x,y)}$, we

can let $t = \frac{1}{x}$ to simplify this quotient to a function

`{{mvar|f}}`{=mediawiki} of the single variable $\frac{y}{x}$:



$$\frac{M(x,y)}{N(x,y)} = \frac{M(tx,ty)}{N(tx,ty)} = \frac{M(1,y/x)}{N(1,y/x)}=f(y/x)\,.$$

That is



$$\frac{dy}{dx} = -f(y/x).$$



Introduce the change of variables

`{{math|1=''y'' = ''ux''}}`{=mediawiki}; differentiate using the product

rule:



$$\frac{dy}{dx}=\frac{d(ux)}{dx} = x\frac{du}{dx} + u\frac{dx}{dx} = x\frac{du}{dx} + u.$$



This transforms the original differential equation into the separable

form



:   $x\frac{du}{dx} = -f(u) - u,$



or



:   $\frac 1x\frac{dx}{du} = \frac {-1}{f(u) + u},$



which can now be integrated directly: `{{math|ln ''x''}}`{=mediawiki}

equals the antiderivative of the right-hand side (see ordinary

differential equation).



\paragraph{Special case}



A first order differential equation of the form

(`{{mvar|a}}`{=mediawiki}, `{{mvar|b}}`{=mediawiki},

`{{mvar|c}}`{=mediawiki}, `{{mvar|e}}`{=mediawiki},

`{{mvar|f}}`{=mediawiki}, `{{mvar|g}}`{=mediawiki} are all constants)



$$\left(ax + by + c\right) dx + \left(ex + fy + g\right) dy = 0$$ where

`{{math|''af'' ? ''be''}}`{=mediawiki} can be transformed into a

homogeneous type by a linear transformation of both variables

(`{{mvar|a}}`{=mediawiki} and `{{mvar|ß}}`{=mediawiki} are constants):



$$t = x + \alpha; \; \; z = y + \beta \,.$$



\paragraph{Homogeneous linear differential equations}



A linear differential equation is \textbf{homogeneous} if it is a homogeneous

linear equation in the unknown function and its derivatives. It follows

that, if `{{math|''f''(''x'')}}`{=mediawiki} is a solution, so is

`{{math|''cf''(''x'')}}`{=mediawiki}, for any (non-zero) constant

`{{mvar|c}}`{=mediawiki}. In order for this condition to hold, each

nonzero term of the linear differential equation must depend on the

unknown function or any derivative of it. A linear differential equation

that fails this condition is called \textbf{inhomogeneous.}



A linear differential equation can be represented as a linear operator

acting on `{{math|''y''(''x'')}}`{=mediawiki} where

`{{mvar|x}}`{=mediawiki} is usually the independent variable and

`{{mvar|y}}`{=mediawiki} is the dependent variable. Therefore, the

general form of a linear homogeneous differential equation is



$$L(y) = 0$$



where `{{mvar|L}}`{=mediawiki} is a differential operator, a sum of

derivatives (defining the \""0th derivative\"" as the original,

non-differentiated function), each multiplied by a function

`{{math|''f''''i''}}`{=mediawiki} of

`{{mvar|x}}`{=mediawiki}:



$$L = \sum_{i=0}^n f_i(x)\frac{d^i}{dx^i} \, ,$$ where

`{{math|''f''''i''}}`{=mediawiki} may be constants, but not

all `{{math|''f''''i''}}`{=mediawiki} may be zero.



For example, the following linear differential equation is homogeneous:



$$\sin(x) \frac{d^2y}{dx^2} + 4 \frac{dy}{dx} + y = 0 \,,$$



whereas the following two are inhomogeneous:



$$2 x^2 \frac{d^2y}{dx^2} + 4 x \frac{dy}{dx} + y = \cos(x) \,;$$



$$2 x^2 \frac{d^2y}{dx^2} - 3 x \frac{dy}{dx} + y = 2 \,.$$ The

existence of a constant term is a sufficient condition for an equation

to be inhomogeneous, as in the above example.



\paragraph{Second order equations}



In order to show how these equations are solved, lets start with the

most basic case - a second order equation



$$A\frac{d^2y}{dx^2}+B\frac{dy}{dx}+Cy=0$$



where A, B, and C are constants.



\paragraph{The Auxiliary Quadratic}



For a DE



$$a\frac{d^2y}{dx^2}+b\frac{dy}{dx}+cy=0,$$



Make the following substitution:



$$y=e^{mx}\,$$



This gives also



$$\frac{dy}{dx}=me^{mx}\,$$



$$\frac{d^2y}{dx^2}=m^2e^{mx}\,$$



The DE is now



$$am^2e^{mx}+bme^{mx}+ce^{mx}=0 \,$$



Dividing by $e^{mx}$ gives (note$$e^{mx}$$ can never equal zero)



$$am^2+bm+c=0 \,$$



This is the \textbf{auxiliary quadratic} (AQ) of the DE. There are four

classes of outcomes to the auxiliary quadratic:



1.  $b^2-4ac > 0, \,$, giving two distinct, real, roots.

2.  $b^2-4ac = 0, \,$, giving two coincident real, roots.

3.  $b^2-4ac < 0, \,$, giving complex roots.



:   \textit{a}: Purely imaginary roots.

:   \textit{b}: Complex-conjugate pair.



The method of solution of the DE depends on the class of AQ.



\paragraph{Class 1: Distinct, Real Roots}



Consider the DE



$$\frac{d^2y}{dx^2}+\frac{dy}{dx}-6y=0$$



The AQ is



$$m^2+m-6=0 \,$$



$$(m+3)(m-2)=0 \,$$



This gives us the following roots:



$$m=2, \ m=-3$$



Going back to the substitution we made to obtain the AQ, we have



$$y=e^{2x},\ y=e^{-3x}$$



as two distinct solutions to the DE. According to the superposition

principle, the general solution is therefore



$$y=Ae^{2x}+Be^{-3x} \,$$



\paragraph{General Solution to Class 1 DEs}



:   Generalizing, for the Second Order DE



$$a\frac{d^2y}{dx^2}+b\frac{dy}{dx}+cy=0$$



:   with the auxiliary quadratic given by



$$am^2+bm+c=0 \,$$



:   with roots \textit{a} and \textit{ß}, the general solution is



$$y=Ae^{\alpha x}+Be^{\beta x} \,$$



\paragraph{Class 2: Coincident, Real Roots}



Consider the DE



$$\frac{d^2y}{dx^2}-4\frac{dy}{dx}+4y=0$$



The AQ is



$$m^2-4m+4=0 \,$$



$$(m-2)(m-2)=0 \,$$



so, $e^{2x}\,$ is a solution. However, we cannot have it as both

solutions as the factor of two produced will be absorbed into the

constant, leaving us with only one constant, and therefore a DE without

a full solution.



For the other solution we will use the Method of Reduction of Order. To

do so we assume that it is in the form of:



$$y_2=u(x)y_1=u(x)e^{2x}\,$$ At the end we will check if our assumption

is correct. We will now substitute this equation and solve for u(x)



$$y_2''-4y_2'+4y_2=0\,$$



$$u''(x)e^{2x}+2u'(x)e^{2x}+2u'(x)e^{2x}+4u(x)e^{2x}-4(u'(x)e^{2x}+2u(x)e^{2x})+4u(x)e^{2x}=0\,$$



$$u''(x)e^{2x}=0\,$$



$$e^{2x}\,$$ is always non-zero so the only way the product can equal

zero is if:



$$u''(x)=0\,$$ Integrating twice offers



$$u(x)=a_1x + a_2\,$$ Therefore



$$y_2=(a_1x+a_2)y_1\,$$



$$y_2=(a_1x+a_2)e^{2x}\,$$ Our general solution is:



$$y=C_1y_1+C_2y_2\,$$



$$y=C_1e^{2x}+C_2(a_1x+a_2)e^{2x}\,$$



$$y=(C_1+a_2)e^{2x}+C_2a_1xe^{2x}\,$$ Because each constant is arbitrary

we can simply write



$$y=C_1e^{2x}+C_2xe^{2x}\,$$



The Method of Reduction of Order can be used on different equations and

u(x) does not always equal x. You can see below that $y=xe^{x}$ is a

valid solution.



$$y=xe^{2x}\,$$



$$\frac{dy}{dx}=2xe^{2x}+e^{2x}=e^{2x}(2x+1)\,$$



$$\frac{d^2y}{dx^2}=2e^{2x}+2e^{2x}(2x+1)=e^{2x}(4x+4)\,$$



To check substitute these into the original DE:



$$e^{2x}(4x+4)-4e^{2x}(2x+1)+4xe^{2x}=0 \,$$



$$4xe^{2x}+4e^{2x}-8xe^{2x}-4e^{2x}+4xe^{2x}=0 \,$$



$$0=0\,$$



Therefore, $y=xe^{2x}\,$ is a solution as well.



\paragraph{General Solution to Class 2 DEs}



:   Generalizing, for the Second Order DE



$$a\frac{d^2y}{dx^2}+b\frac{dy}{dx}+cy=0$$



:   with the coincident root \textit{a} of the AQ, the general solution is



$$y=(A+Bx)e^{\alpha x} \,$$



\paragraph{Class 3a: Purely Imaginary Roots}



To have complex roots, the AQ must have a discriminant less than zero,

so



$$b^2-4ac < 0 \,$$



Also, for the solution to be purely imaginary, the value of \textit{b} must be

exactly zero.



Therefore,



$$4ac>0\,$$



This means that \textit{a} and \textit{c} have to have the same signs: either \textit{a} and

\textit{c} are both positive or they are both negative. If we consider our

general second-order DE:



$$a\frac{d^2y}{dx^2}+b\frac{dy}{dx}+cy=0$$



Setting \textit{b} to zero gives



$$a\frac{d^2y}{dx^2}+cy=0$$



Dividing through by \textit{a} gives



$$\frac{d^2y}{dx^2}+\frac{c}{a}y=0$$.



Therefore, the \textit{y} term is always positive, and this can be represented

by



$$\frac{d^2y}{dx^2}+\omega^2y=0$$.



(I\'m using \textit{?} here as it is used for simple harmonic motion, which is

the primary use of this DE). There are now two paths to the solution of

the DE. The first relies on us spotting that we can use the cyclical

nature of trig. functions when derived. Substitute the following



$$y=\cos \omega x \,$$



$$\frac{dy}{dx}=-\omega \sin \omega x\,$$



$$\frac{d^2y}{dx^2}=-\omega^2 \cos \omega x\,$$



And check in our DE:



$$\frac{d^2y}{dx^2}+\omega^2 y=-\omega^2 \cos \omega x+\omega^2 \cos \omega x=0$$



This checks out, so $\cos \omega x\,$ is a solution. A similar result

holds true for the substitution using $y=\sin \omega x \,$.



Our solutions are therefore



$$y=\cos \omega x \,$$



$$y=\sin \omega x \,$$



So the general solution is



$$y=A\cos \omega x+B\sin \omega x. \,$$



The other method of solving this equation is to use Euler\'s Formula:



$$e^{\omega ix}=\cos \omega x +i \sin x \,$$



:   and



$$e^{-\omega ix}=\cos \omega x -i \sin x \,$$



From our original DE, we have an AQ of



$$m^2+\omega^2=0 \,$$



giving us roots of



$$m=\pm i \omega \,$$



so the general solution, similar to the Class 1 DEs, is



$$y=Ae^{\omega ix}+Be^{-\omega ix} \,$$



$$y=A(\cos \omega x +i \sin \omega x)+B(\cos \omega x -i \sin \omega x)\,$$



$$y=(A+B)\cos \omega x+i(A-B)\sin \omega x \,$$



Since A and B are arbitrary, we can set new constants for convenience,

letting our new A equal A+B and our new B equal \textit{i}(A-B).



Thus we have as our general solution



$$y=A\cos \omega x+B\sin \omega x \,$$



\paragraph{General Solution to Class 3a DEs}



:   Generalizing, for the Second Order DE



$$\frac{d^2y}{dx^2}+\omega^2 y=0$$



:   the general solution is



$$y=A\cos \omega x+B\sin \omega x \,$$



\paragraph{Class 3b: Complex Conjugate Roots}



Since it is a proven theorem that complex roots of polynomials always

occur in conjugate pairs, the only remaining class of AQ is the one with

complex conjugates for solutions.



Given that the solutions are complex, we know that in the AQ



$$am^2+bm+c=0 \,$$



$$b^2-4ac,<0,\ b \ne 0 \,$$ (see Class 3a).



The roots of this are in the form



$$(p \pm iq)\,$$



The general solution is then



$$y=Ae^{(p+iq)x}+Be^{(p-iq)x}\,$$



$$y=e^{px} \left( Ae^{iqx}+Be^{-iqx}\right)\,$$



From Euler\'s Formulas, we can now get



$$y=e^{px}\left( A\left(\cos q x + i \sin q x \right)+B\left(\cos q x - i \sin q x \right) \right)\,$$



$$y=e^{px}\left( (A+B)\cos q x+i(A-B)\sin q x \right)\,$$



As A and B are arbitrary, we can collapse them as in Class 3a, so that

we have the general solution



$$y=e^{px}\left( A\cos q x+B\sin q x \right)\,$$



\paragraph{General Solution to Class 3b DEs}



:   Generalizing, for the Second Order DE



$$a\frac{d^2y}{dx^2}+b\frac{dy}{dx}+cy=0$$



:   with an AQ with roots



$$m=p \pm iq,$$



:   The general solution is



$$y=e^{px}\left( A\cos q x+B\sin q x \right)\,$$



We have now covered all possible types of homogeneous second-order

differential equation, and we didn\'t even have to integrate anything!

We will now have a look at higher order equivalents.



\paragraph{Licensing}



Content obtained and/or adapted from:



\begin{itemize}
    \item \href{https://en.wikipedia.org/wiki/Homogeneous_differential_equation}{Homogeneous differential equation,
\end{itemize}
    Wikipedia}

    under a CC BY-SA license

\begin{itemize}
    \item \href{https://en.wikibooks.org/wiki/Ordinary_Differential_Equations/Homogenous_1}{Homogeneous 1, Wikibooks: Ordinary Differential
\end{itemize}
    Equations}

    under a CC BY-SA license"

%%===============================================================
\section{Abel’s Theorem}
\label{sec:abel-s-theorem}
%%===============================================================

"In mathematics, \textbf{Abel\'s theorem} for power series relates a limit of

a power series to the sum of its coefficients. It is named after

Norwegian mathematician Niels Henrik Abel.



\paragraph{Theorem}



Let



$$G (x) = \sum_{k=0}^\infty a_k x^k$$



be a power series with real coefficients \textit{a}\textasciitilde{}\textit{k}\textasciitilde{} with radius of

convergence 1. Suppose that the series



$$\sum_{k=0}^\infty a_k$$



converges. Then \textit{G}(\textit{x}) is continuous from the left at $x=1$, i.e.



$$\lim_{x\to 1^-} G(x) = \sum_{k=0}^\infty a_k.$$



The same theorem holds for complex power series



$$G(z) = \sum_{k=0}^\infty a_k z^k,$$



provided that $z \to 1$ within a \textit{Stolz sector}, that is, a region of

the open unit disk where



:   $|1-z|\leq M(1-|z|)$



for some $M > 1$. Without this restriction, the limit may fail to exist:

for example, the power series



$$\sum_{n>0} \frac{z^{3^n}-z^{2\cdot 3^n}} n$$



converges to 0 at \textit{z} = 1, but is unbounded near any point of the form

\textit{e}^$\pi$i/3^\textit{n}^^, so the value at \textit{z} = 1 is not the limit as \textit{z}

tends to 1 in the whole open disk.



Note that \textit{G}(\textit{z}) is continuous on the real closed interval \[0, \textit{t}\]

for \textit{t} \< 1, by virtue of the uniform convergence of the series on

compact subsets of the disk of convergence. Abel\'s theorem allows us to

say more, namely that \textit{G}(\textit{z}) is continuous on \[0, 1\].



\paragraph{Remarks}



As an immediate consequence of this theorem, if \textit{z} is any nonzero

complex number for which the series



$$\sum_{k=0}^\infty a_k z^k$$



converges, then it follows that



$$\lim_{t\to 1^{-}} G(tz) = \sum_{k=0}^\infty a_kz^k$$



in which the limit is taken from below.



The theorem can also be generalized to account for sums which diverge to

infinity. If



$$\sum_{k=0}^\infty a_k = \infty$$



then



$$\lim_{z\to 1^{-}} G(z) \to \infty.$$



However, if the series is only known to be divergent, but for reasons

other than diverging to infinity, then the claim of the theorem may

fail: take, for example, the power series for



$$\frac{1}{1+z}.$$



At $z=1$ the series is equal to $1 - 1 + 1 - 1 + \cdots,$ but

$\tfrac{1}{1+1}=\tfrac{1}{2}.$



We also remark the theorem holds for radii of convergence other than

$R=1$: let



$$G(x)=\sum_{k=0}^\infty a_kx^k$$



be a power series with radius of convergence $R$, and suppose the series

converges at $x=R$. Then $G(x)$ is continuous from the left at $x=R$,

i.e.



$$\lim_{x\to R^-}G(x)=G(R)$$



\paragraph{Applications}



The utility of Abel\'s theorem is that it allows us to find the limit of

a power series as its argument (i.e. $z$) approaches 1 from below, even

in cases where the radius of convergence, $R$, of the power series is

equal to 1 and we cannot be sure whether the limit should be finite or

not. See e.g. the binomial series. Abel\'s theorem allows us to evaluate

many series in closed form. For example, when



$$a_k = \frac{(-1)^k}{k+1},$$



we obtain



$$G_a(z) = \frac{\ln(1+z)}{z}, \qquad 0 < z < 1,$$



by integrating the uniformly convergent geometric power series term by

term on $[-z, 0]$; thus the series



$$\sum_{k=0}^\infty \frac{(-1)^k}{k+1}$$



converges to $\ln(2)$ by Abel\'s theorem. Similarly,



$$\sum_{k=0}^\infty \frac{(-1)^k}{2k+1}$$



converges to $\arctan(1) = \tfrac{\pi}{4}.$



$G_a(z)$ is called the generating function of the sequence $a$. Abel\'s

theorem is frequently useful in dealing with generating functions of

real-valued and non-negative sequences, such as probability-generating

functions. In particular, it is useful in the theory of Galton--Watson

processes.



\paragraph{Outline of proof}



After subtracting a constant from $a_0$, we may assume that

$\sum_{k=0}^\infty a_k=0$. Let $s_n=\sum_{k=0}^n a_k\!$. Then

substituting $a_k=s_k-s_{k-1}$ and performing a simple manipulation of

the series (summation by parts) results in



$$G_a(z) = (1-z)\sum_{k=0}^{\infty} s_k z^k.$$



Given $\varepsilon > 0,$ pick \textit{n} large enough so that

$|s_k| < \varepsilon$ for all $k\ge n$ and note that



$$\left|(1-z)\sum_{k=n}^\infty s_kz^k \right| \le \varepsilon |1-z|\sum_{k=n}^\infty |z|^k = \varepsilon|1-z|\frac{|z|^n}{1-|z|} < \varepsilon M$$



when \textit{z} lies within the given Stolz angle. Whenever \textit{z} is sufficiently

close to 1 we have



$$\left|(1-z)\sum_{k=0}^{n-1} s_kz^k \right| < \varepsilon,$$



so that $|G_a(z)| < (M+1)\varepsilon$ when \textit{z} is both sufficiently

close to 1 and within the Stolz angle.



\paragraph{Related concepts}



Converses to a theorem like Abel\'s are called Tauberian theorems: There

is no exact converse, but results conditional on some hypothesis. The

field of divergent series, and their summation methods, contains many

theorems \textit{of abelian type} and \textit{of tauberian type}.



\paragraph{Licensing}



Content obtained and/or adapted from:



\begin{itemize}
    \item \href{https://en.wikipedia.org/wiki/Abel\%27s_theorem}{Abel\'s theorem,
\end{itemize}
    Wikipedia} under a

    CC BY-SA license"

%%===============================================================
\section{Fundamental Solutions}
\label{sec:fundamental-solutions}
%%===============================================================

Fundamental solutions to linear homogeneous equations



> 

>

> ``{=html}Theorem 1:``{=html} Let

> $\frac{d^2 y}{dt^2} + p(t) \frac{dy}{dt} + q(t)y = 0$ be a second

> order linear homogenous differential equation where $p$ and $q$ are

> continuous on an open interval $I$ such that $t_0 \in I$, and let

> $y = y_1(t)$ and $y = y_2(t)$ be two solutions to this differential

> equation. The set of all linear combinations of these two solutions,

> $y = Cy_1(t) + Dy_2(t)$ where $C$ and $D$ are constants contains all

> solutions to this differential equation if and only if there exists a

> point $t_0$ for which the Wronksian of $y_1$ and $y_2$ at $t_0$ is

> nonzero, that is $W(y_1, y_2) \bigg|_{t_0} \neq 0$.

>

> 



\begin{itemize}
    \item ``{=html}Proof:``{=html} Let $y = y_1(t)$ and
\end{itemize}
    $y = y_2(t)$ both be solutions to the differential equation

    $\frac{d^2 y}{dt^2} + p(t) \frac{dy}{dt} + q(t)y = 0$ and suppose

    that $y = \phi (t)$ is any arbitrary solution as well. We want to

    show that $\phi (t)$ is a linear combination of $y_1$ and $y_2$ for

    some constants $C$ and $D$.



<!-- -->



\begin{itemize}
    \item $\Leftarrow$ Let $t_0$ be such that the Wronskian of $y_1$ and $y_2$
\end{itemize}
    evaluated at $t_0$ is nonzero, that is:







$\begin{aligned} \quad W(y_1, y_2) \bigg|_{t_0} \neq 0 \end{aligned}$







\begin{itemize}
    \item Take this value $t_0$ and evaluate both $\phi$ and $\phi'$ at this
\end{itemize}
    point. Then $\phi (t_0) = y_0$ and $\phi'(t_0) = y'_0$ (since

    $y = \phi(t)$ is a solution to our differential equation). Now

    consider the initial value problem

    $\frac{d^2 y}{dt^2} + p(t) \frac{dy}{dt} + q(t) y = 0$ with the

    initial conditions $y(t_0) = y_0$ and $y'(t_0) = y'_0$. The function

    $\phi$ satisfies this differential equation. Since

    $W(y_1, y_2) \bigg|_{t_0} \neq 0$ then we have that there exists

    constants $C$ and $D$ such that $y = Cy_1(t) + Dy_2(t)$ satisfies

    this initial value problem. But since $p$ and $q$ are continuous on

    the open interval $I$ containing $t_0$ then this implies that a

    unique solution exists, and so:







$\begin{aligned} \phi(t) = Cy_1(t) + Dy_2(t) \end{aligned}$







\begin{itemize}
    \item So all solutions for this differential equation are a linear
\end{itemize}
    combination of the solutions $y = y_1(t)$ and $y = y_2(t)$.



<!-- -->



\begin{itemize}
    \item $\Rightarrow$ Suppose that every point $t_0 \in I$ is such that
\end{itemize}
    $W(y_1, y_2) \bigg|_{t_0} = 0$, that is, there exists no point $t_0$

    on $I$ where the Wronskian of $y_1$ and $y_2$ evaluated at $t_0$ is

    nonzero. Let $y_0$ and $y'_0$ be values for which the system

    $\left\{\begin{matrix} Cy_1(t_0) + Dy_2(t_0) = y_0 \\\\ Cy_1'(t_0) + Dy_2'(t_0) = y'_0 \end{matrix}\right.$

    has no solutions for a set of constants $C$ and $D$.



<!-- -->



\begin{itemize}
    \item Now since $p$ and $q$ are continuous on an open interval $I$
\end{itemize}
    containing $t_0$, such a solution $\phi(t)$ satisfies the initial

    conditions $y(t_0) = y_0$ and $y'(t_0) = y'_0$. Note though this

    solution is not a linear combination of $y_1$ and $y_2$ though which

    completes our proof. $\blacksquare$



Theorem 1 above implies that if we can find two solutions $y = y_1(t)$

and $y = y_2(t)$ for which the Wronskian $W(y_1, y_2) \neq 0$, then for

constants $C$ and $D$, all solutions of the second order linear

homogenous differential equation

$\frac{d^2 y}{dt^2} + p(t) \frac{dy}{dt} + q(t) y = 0$ are given by:







$\begin{aligned} \quad y = Cy_1(t) + Dy_2(t) \end{aligned}$







Also note that thus far we have not said that $y = y_1(t)$ and

$y = y_2(t)$ need to be distinct. However, with the Theorem above, we

see that if $y_1(t) = y_2(t)$ then the Wronskian

$W(y_1, y_2) = W(y_1, y_1) = W(y_2, y_2)$ is zero (as you should verify)

and so not all solutions to a second order linear homogenous

differential are given by the linear combination of just $y_1$.



> 

>

> ``{=html}Definition:``{=html} Let

> $\frac{d^2 y}{dt^2} + p(t) \frac{dy}{dt} + q(t)y = 0$ where $p$ and

> $q$ are continuous on an open interval $I$ such that $t_0 \in I$ and

> let $y = y_1(t)$ and $y = y_2(t)$ be solutions to this differential

> equation. If the Wronskian $W(y_1, y_2) \neq 0$ then the set of linear

> combinations of $y_1$ and $y_2$ is known as the

> ``{=html}Fundamental Set of Solutions``{=html} to

> this differential equation.

>

> 



From the definition above, we see that if we can find two solutions

$y = y_1(t)$ and $y = y_2(t)$ for which the Wronskian $W(y_1, y_2)$ is

nonzero, then $y_1$ and $y_2$ form a fundamental set of solutions. The

next question that we might pose is whether or not a second order linear

homogenous differential equation always has a fundamental set of

solutions.



> 

>

> ``{=html}Theorem 2:``{=html} Let

> $\frac{d^2 y}{dt} + p(t) \frac{dy}{dt} + q(t) y = 0$ be a second order

> linear homogenous differential equation where $p$ and $q$ are

> continuous on an open interval $I$ such that $t_0 \in I$. If

> $y = y_1(t)$ is a solution to this differential equation that

> satisfies the initial conditions $y_1(t_0) = 1$ and $y_1'(t_0) = 0$,

> and if $y = y_2(t)$ be a solution to this differential equation that

> satisfies the initial conditions $y_2(t_0) = 0$ and $y_2'(t_0) = 1$.

> Then $y_1$ and $y_2$ form a fundamental set of solutions for this

> differential equation.

>

> 



\begin{itemize}
    \item ``{=html}Proof:``{=html} We note that:
\end{itemize}





$\begin{aligned} \quad W(y_1, y_2) \bigg|_{t_0} = \begin{vmatrix} y_1(t_0)  y_2(t_0) \\\\ y_1'(t_0)  y_2'(t_0)\end{vmatrix} = \begin{vmatrix} 1  0\\\\ 0  1 \end{vmatrix} = 1 \neq 0 \end{aligned}$







\begin{itemize}
    \item Thus Theorem 1 implies that ALL solutions to this differential
\end{itemize}
    equation are given by $y = Cy_1(t) + Dy_2(t)$ where $C$ and $D$ are

    constants. Thus $y_1$ and $y_2$ form a fundamental set of solutions

    for this differential equation. $\blacksquare$



\paragraph{Fundamental sets of solutions}



> 

>

> ``{=html}Definition:``{=html} A

> ``{=html}Fundamental Set of Solutions``{=html} to the

> linear homogeneous system of first order ODEs

> $\mathbf{x}' = A(t) \mathbf{x}$ on $J = (a, b)$ is a set

> $\{ \phi^{[1]}, \phi^{[2]}, ..., \phi^{[n]} \}$ of linearly

> independent solutions to this system on $J$.

>

> 



For example, consider the following linear homogeneous system of $2$

first order ODEs:







$\begin{aligned} \quad \left\{\begin{matrix} x_1' = x_1\\\\ x_2' = 2x_2 \end{matrix}\right. \end{aligned}$







We can easily solve this system. For the first differential equation:







$\begin{aligned} \quad \frac{dx_1}{dt} &= x_1 \\\\ \frac{dx_1}{x_1} &= dt \\\\ \int \frac{1}{x_1}  dx &= \int  dt \\\\ \ln (x_1) &= t + C \\\\ x_1 &= e^{t + C} \\\\ x_1 &= C_1e^t \end{aligned}$







Where $C_1 = e^C > 0$.



For the second differential equation:







$\begin{aligned} \quad \frac{dx_2}{dt} &= 2x_2 \\\\ \frac{dx_2}{x_2} &= 2  dt \\\\ \int \frac{1}{x_2}  dx_2 &= \int 2  dt \\\\ \ln (x_2) &= 2t + C \\\\ x_2 &= e^{2t + C} \\\\ x_2 &= C_2e^{2t} \end{aligned}$







Where $C_2 = e^C > 0$.



Note that in fact $C_1, C_2$ can be any real numbers are we are not

simply restricted to $C_1, C_2 > 0$.



Now by taking $C_1 = 1$, $C_2 = 0$ we get that

$\phi^{[1]} = \begin{bmatrix} e^t\\\\ 0 \end{bmatrix}$ is a solution to

this system. Also, by taking $C_1 = 0$ and $C_2 = 1$ we get that

$\phi^{[2]} = \phi^{[1]} = \begin{bmatrix} 0\\\\ e^{2t} \end{bmatrix}$ is

a solution to this system.



We will now show that $\{ \phi^{[1]}, \phi^{[2]} \}$ is a Fundamental

set of solutions to this system on all of $\mathbb{R}$. Let

$\alpha, beta \in \mathbb{R}$ and consider the following equation:







$\begin{aligned} \alpha \phi^{[1]} + \beta \phi^{[2]} &= 0 \\\\ \alpha \begin{bmatrix} e^t\\\\ 0 \end{bmatrix} + \beta \begin{bmatrix} 0\\\\ e^{2t} \end{bmatrix} &= \begin{bmatrix} 0 \\\\ 0 \end{bmatrix} \\\\ \begin{bmatrix} \alpha e^t \\\\ \beta e^{2t} \end{bmatrix} &= \begin{bmatrix} 0 \\\\ 0 \end{bmatrix} \end{aligned}$







The equation above implies that $\alpha e^t = 0$ and $\beta e^{2t} = 0$

for all $t \in \mathbb{R}$. Since $e^t, e^{2t} > 0$ for all

$t \in \mathbb{R}$ this implies that $\alpha, \beta = 0$. So

$\{ \phi^{[1]}, \phi^{[2]} \}$ is a linearly independent set of

solutions to this system and so $\{ \phi^{[1]}, \phi^{[2]} \}$ is a

fundamental set of solutions to this system on $\mathbb{R}$.



\paragraph{Resources}



\begin{itemize}
    \item \href{https://tutorial.math.lamar.edu/Classes/DE/FundamentalSetsofSolutions.aspx}{Fundamental Sets of Solutions
\end{itemize}
    Notes}.

    Produced by Paul Dawkins, Lamar University

\begin{itemize}
    \item \href{https://www.cfm.brown.edu/people/dobrush/am33/Mathematica/ch4/fundset.html}{Fundamental Set of
\end{itemize}
    Solutions},

    Brown University



\paragraph{Licensing}



Content obtained and/or adapted from:



\begin{itemize}
    \item \href{http://mathonline.wikidot.com/fundamental-solutions-to-linear-homogenous-differential-equa}{Fundamental Solutions to Linear Homogenous Differential Equations,
\end{itemize}
    mathonline.wikidot.com}

    under a CC BY-SA license

\begin{itemize}
    \item \href{http://mathonline.wikidot.com/fundamental-sets-of-solutions-to-a-linear-homogeneous-system}{Fundamental Sets of Solutions,
\end{itemize}
    mathonline.wikidot.com}

    under a CC BY-SA license"

%%===============================================================
\section{Linear Non-homogeneous Equations}
\label{sec:linear-non-homogeneous-equations}
%%===============================================================

"In mathematics, a \textbf{linear differential equation} is a differential

equation that is defined by a linear polynomial in the unknown function

and its derivatives, that is an equation of the form



$$a_0(x)y + a_1(x)y' + a_2(x)y'' \cdots + a_n(x)y^{(n)} = b(x)$$ where

\textit{a}\textasciitilde{}0\textasciitilde{}(\textit{x}), ..., \textit{a}\textasciitilde{}\textit{n}\textasciitilde{}(\textit{x}) and `{{math|''b''(''x'')}}`{=mediawiki}

are arbitrary differentiable functions that do not need to be linear,

and `{{math|''y''', …, ''y''(''n'') }}`{=mediawiki} are the

successive derivatives of an unknown function `{{mvar|y}}`{=mediawiki}

of the variable `{{mvar|x}}`{=mediawiki}.



Such an equation is an ordinary differential equation (ODE). A *linear

differential equation* may also be a linear partial differential

equation (PDE), if the unknown function depends on several variables,

and the derivatives that appear in the equation are partial derivatives.



A linear differential equation or a system of linear equations such that

the associated homogeneous equations have constant coefficients may be

solved by quadrature, which means that the solutions may be expressed in

terms of integrals. This is also true for a linear equation of order

one, with non-constant coefficients. An equation of order two or higher

with non-constant coefficients cannot, in general, be solved by

quadrature. For order two, Kovacic\'s algorithm allows deciding whether

there are solutions in terms of integrals, and computing them if any.



The solutions of linear differential equations with polynomial

coefficients are called holonomic functions. This class of functions is

stable under sums, products, differentiation, integration, and contains

many usual functions and special functions such as exponential function,

logarithm, sine, cosine, inverse trigonometric functions, error

function, Bessel functions and hypergeometric functions. Their

representation by the defining differential equation and initial

conditions allows making algorithmic (on these functions) most

operations of calculus, such as computation of antiderivatives, limits,

asymptotic expansion, and numerical evaluation to any precision, with a

certified error bound.



\paragraph{Basic terminology}



The highest order of derivation that appears in a (linear) differential

equation is the \textit{order} of the equation. The term

`{{math|''b''(''x'')}}`{=mediawiki}, which does not depend on the

unknown function and its derivatives, is sometimes called the *constant

term* of the equation (by analogy with algebraic equations), even when

this term is a non-constant function. If the constant term is the zero

function, then the differential equation is said to be \textit{homogeneous}, as

it is a homogeneous polynomial in the unknown function and its

derivatives. The equation obtained by replacing, in a linear

differential equation, the constant term by the zero function is the

\textit{associated homogeneous equation}. A differential equation has *constant

coefficients* if only constant functions appear as coefficients in the

associated homogeneous equation.



A \textit{solution} of a differential equation is a function that satisfies the

equation. The solutions of a homogeneous linear differential equation

form a vector space. In the ordinary case, this vector space has a

finite dimension, equal to the order of the equation. All solutions of a

linear differential equation are found by adding to a particular

solution any solution of the associated homogeneous equation.



\paragraph{Linear differential operator}



A \textit{basic differential operator} of order `{{mvar|i}}`{=mediawiki} is a

mapping that maps any differentiable function to its

`{{mvar|i}}`{=mediawiki}th derivative, or, in the case of several

variables, to one of its partial derivatives of order

`{{mvar|i}}`{=mediawiki}. It is commonly denoted



$$\frac{d^i}{dx^i}$$ in the case of univariate functions, and



$$\frac{\partial^{i_1+\cdots +i_n}}{\partial x_1^{i_1}\cdots \partial x_n^{i_n}}$$

in the case of functions of `{{mvar|n}}`{=mediawiki} variables. The

basic differential operators include the derivative of order 0, which is

the identity mapping.



A \textbf{linear differential operator} (abbreviated, in this article, as

\textit{linear operator} or, simply, \textit{operator}) is a linear combination of

basic differential operators, with differentiable functions as

coefficients. In the univariate case, a linear operator has thus the

form



$$a_0(x)+a_1(x)\frac{d}{dx} + \cdots +a_n(x)\frac{d^n}{dx^n},$$ where

`{{math|''a''0(''x''), …, ''a''''n''(''x'')}}`{=mediawiki}

are differentiable functions, and the nonnegative integer

`{{mvar|n}}`{=mediawiki} is the \textit{order} of the operator (if

`{{math|''a''''n''(''x'')}}`{=mediawiki} is not the zero

function).



Let `{{mvar|L}}`{=mediawiki} be a linear differential operator. The

application of `{{mvar|L}}`{=mediawiki} to a function

`{{mvar|f}}`{=mediawiki} is usually denoted

`{{math|''Lf''}}`{=mediawiki} or `{{math|''Lf''(''X'')}}`{=mediawiki},

if one needs to specify the variable (this must not be confused with a

multiplication). A linear differential operator is a linear operator,

since it maps sums to sums and the product by a scalar to the product by

the same scalar.



As the sum of two linear operators is a linear operator, as well as the

product (on the left) of a linear operator by a differentiable function,

the linear differential operators form a vector space over the real

numbers or the complex numbers (depending on the nature of the functions

that are considered). They form also a free module over the ring of

differentiable functions.



The language of operators allows a compact writing for differentiable

equations: if



$$L=a_0(x)+a_1(x)\frac{d}{dx} + \cdots +a_n(x)\frac{d^n}{dx^n},$$ is a

linear differential operator, then the equation



$$a_0(x)y +a_1(x)y' + a_2(x)y'' +\cdots +a_n(x)y^{(n)}=b(x)$$ may be

rewritten



$$Ly=b(x).$$



There may be several variants to this notation; in particular the

variable of differentiation may appear explicitly or not in

`{{mvar|y}}`{=mediawiki} and the right-hand and of the equation, such as

`{{math|1=''Ly''(''x'') = ''b''(''x'')}}`{=mediawiki} or

`{{math|1=''Ly'' = ''b''}}`{=mediawiki}.



The \textit{kernel} of a linear differential operator is its kernel as a linear

mapping, that is the vector space of the solutions of the (homogeneous)

differential equation `{{math|1=''Ly'' = 0}}`{=mediawiki}.



In the case of an ordinary differential operator of order

`{{mvar|n}}`{=mediawiki}, Carathéodory\'s existence theorem implies

that, under very mild conditions, the kernel of `{{mvar|L}}`{=mediawiki}

is a vector space of dimension `{{mvar|n}}`{=mediawiki}, and that the

solutions of the equation

`{{math|1=''Ly''(''x'') = ''b''(''x'')}}`{=mediawiki} have the form



$$S_0(x) + c_1S_1(x) + \cdots +c_nS_n(x),$$ where

`{{math|''c''1, …, ''c''''n''}}`{=mediawiki} are

arbitrary numbers. Typically, the hypotheses of Carathéodory\'s theorem

are satisfied in an interval `{{mvar|I}}`{=mediawiki}, if the functions

`{{math|''b'', ''a''0, …, ''a''''n''}}`{=mediawiki}

are continuous in `{{mvar|I}}`{=mediawiki}, and there is a positive real

number `{{mvar|k}}`{=mediawiki} such that \|\textit{a}\textasciitilde{}\textit{n}\textasciitilde{}(\textit{x})\| \> \textit{k} for

every `{{mvar|x}}`{=mediawiki} in `{{mvar|I}}`{=mediawiki}.



\paragraph{Homogeneous equation with constant coefficients}



A homogeneous linear differential equation has \textit{constant coefficients}

if it has the form



$$a_0y + a_1y' + a_2y'' + \cdots + a_n y^{(n)} = 0$$ where

`{{math|''a''1, …, ''a''''n''}}`{=mediawiki} are

(real or complex) numbers. In other words, it has constant coefficients

if it is defined by a linear operator with constant coefficients.



The study of these differential equations with constant coefficients

dates back to Leonhard Euler, who introduced the exponential function

`{{math|''e''''x''}}`{=mediawiki}, which is the unique

solution of the equation `{{math|1=''f''' = ''f''}}`{=mediawiki} such

that `{{math|1=''f''(0) = 1}}`{=mediawiki}. It follows that the

`{{mvar|n}}`{=mediawiki}th derivative of

`{{math|''e''''cx'' }}`{=mediawiki} is

`{{math|''c''''n''''e''''cx''}}`{=mediawiki}, and

this allows solving homogeneous linear differential equations rather

easily.



Let



$$a_0y + a_1y' + a_2y'' + \cdots + a_ny^{(n)} = 0$$ be a homogeneous

linear differential equation with constant coefficients (that is

`{{math|''a''0, …, ''a''''n''}}`{=mediawiki} are

real or complex numbers).



Searching solutions of this equation that have the form

`{{math|''e''''ax''}}`{=mediawiki} is equivalent to searching

the constants `{{mvar|a}}`{=mediawiki} such that



$$a_0e^{\alpha x} + a_1\alpha e^{\alpha x} + a_2\alpha^2 e^{\alpha x}+\cdots + a_n\alpha^n e^{\alpha x} = 0.$$

Factoring out `{{math|''e''''ax''}}`{=mediawiki} (which is

never zero), shows that `{{mvar|a}}`{=mediawiki} must be a root of the

\textit{characteristic polynomial}



$$a_0 + a_1t + a_2t^2 + \cdots + a_nt^n$$ of the differential equation,

which is the left-hand side of the characteristic equation



$$a_0 + a_1t + a_2t^2 + \cdots + a_nt^n = 0.$$



When these roots are all distinct, one has `{{mvar|n}}`{=mediawiki}

distinct solutions that are not necessarily real, even if the

coefficients of the equation are real. These solutions can be shown to

be linearly independent, by considering the Vandermonde determinant of

the values of these solutions at

`{{math|1=''x'' = 0, …, ''n'' – 1}}`{=mediawiki}. Together they form a

basis of the vector space of solutions of the differential equation

(that is, the kernel of the differential operator).



+----------------------------------------------------------------------+

| Example                                                              |

+======================================================================+

| $$y''''-2y'''+2y''-2y'+y=0$$ has the characteristic equation         |

|                                                                      |

| :   $z^4-2z^3+2z^2-2z+1=0.$                                          |

|                                                                      |

| This has zeros, `{{mvar|i}}`{=mediawiki},                            |

| `{{math|-''i''}}`{=mediawiki}, and `{{math|1}}`{=mediawiki}          |

| (multiplicity 2). The solution basis is thus                         |

|                                                                      |

| :   $e^{ix}, \; e^{-ix}, \; e^x, \; xe^x.$                              |

|                                                                      |

| A real basis of solution is thus                                     |

|                                                                      |

| :   $\cos x, \; \sin x, \; e^x, \; xe^x.$                               |

+----------------------------------------------------------------------+



In the case where the characteristic polynomial has only simple roots,

the preceding provides a complete basis of the solutions vector space.

In the case of multiple roots, more linearly independent solutions are

needed for having a basis. These have the form



$$x^ke^{\alpha x},$$ where `{{mvar|k}}`{=mediawiki} is a nonnegative

integer, `{{mvar|a}}`{=mediawiki} is a root of the characteristic

polynomial of multiplicity `{{mvar|m}}`{=mediawiki}, and

`{{math|''k'' < ''m''}}`{=mediawiki}. For proving that these functions

are solutions, one may remark that if `{{mvar|a}}`{=mediawiki} is a root

of the characteristic polynomial of multiplicity

`{{mvar|m}}`{=mediawiki}, the characteristic polynomial may be factored

as

`{{math|''P''(''t'') (''t'' - ''a'')''m''}}`{=mediawiki}.

Thus, applying the differential operator of the equation is equivalent

with applying first `{{mvar|m}}`{=mediawiki} times the operator

$\frac{d}{dx} - \alpha$, and then the operator that has

`{{mvar|P}}`{=mediawiki} as characteristic polynomial. By the

exponential shift theorem,



$$\left(\frac{d}{dx}-\alpha\right)\left(x^ke^{\alpha x}\right)= kx^{k-1}e^{\alpha x},$$



and thus one gets zero after `{{math|''k'' + 1}}`{=mediawiki}

application of $\frac{d}{dx} - \alpha$.



As, by the fundamental theorem of algebra, the sum of the multiplicities

of the roots of a polynomial equals the degree of the polynomial, the

number of above solutions equals the order of the differential equation,

and these solutions form a base of the vector space of the solutions.



In the common case where the coefficients of the equation are real, it

is generally more convenient to have a basis of the solutions consisting

of real-valued functions. Such a basis may be obtained from the

preceding basis by remarking that, if

`{{math|''a'' + ''ib''}}`{=mediawiki} is a root of the characteristic

polynomial, then `{{math|''a'' – ''ib''}}`{=mediawiki} is also a root,

of the same multiplicity. Thus a real basis is obtained by using

Euler\'s formula, and replacing

`{{math|''x''''k''''e''(''a''+''ib'')''x''}}`{=mediawiki}

and

`{{math|''x''''k''''e''(''a''-''ib'')''x''}}`{=mediawiki}

by

`{{math|''x''''k''''e''''ax'' cos(''bx'')}}`{=mediawiki}

and

`{{math|''x''''k''''e''''ax'' sin(''bx'')}}`{=mediawiki}.



\paragraph{Second-order case}



A homogeneous linear differential equation of the second order may be

written



$$y'' + ay' + by = 0,$$ and its characteristic polynomial is



$$r^2 + ar + b.$$



If `{{mvar|a}}`{=mediawiki} and `{{mvar|b}}`{=mediawiki} are real, there

are three cases for the solutions, depending on the discriminant

`{{math|1=''D'' = ''a''2 - 4''b''}}`{=mediawiki}. In all

three cases, the general solution depends on two arbitrary constants

`{{math|''c''1}}`{=mediawiki} and

`{{math|''c''2}}`{=mediawiki}.



\begin{itemize}
    \item If `{{math|''D'' > 0}}`{=mediawiki}, the characteristic polynomial
\end{itemize}
    has two distinct real roots `{{mvar|a}}`{=mediawiki}, and

    `{{mvar|ß}}`{=mediawiki}. In this case, the general solution is



$$c_1 e^{\alpha x} + c_2 e^{\beta x}.$$



\begin{itemize}
    \item If `{{math|1=''D'' = 0}}`{=mediawiki}, the characteristic polynomial
\end{itemize}
    has a double root `{{math|-''a''/2}}`{=mediawiki}, and the general

    solution is



$$(c_1 + c_2 x) e^{-ax/2}.$$



\begin{itemize}
    \item If `{{math|''D'' < 0}}`{=mediawiki}, the characteristic polynomial
\end{itemize}
    has two complex conjugate roots

    `{{math|''a'' ± ''ßi''}}`{=mediawiki}, and the general solution is



$$c_1 e^{(\alpha + \beta i)x} + c_2 e^{(\alpha - \beta i)x},$$



:   which may be rewritten in real terms, using Euler\'s formula as

    $$e^{\alpha x}  (c_1\cos(\beta x) + c_2 \sin(\beta x)).$$



Finding the solution `{{math|''y''(''x'')}}`{=mediawiki} satisfying

`{{math|1=''y''(0) = ''d''1}}`{=mediawiki} and

`{{math|1=''y'''(0) = ''d''2}}`{=mediawiki}, one equates the

values of the above general solution at `{{math|0}}`{=mediawiki} and its

derivative there to `{{math|''d''1}}`{=mediawiki} and

`{{math|''d''2}}`{=mediawiki}, respectively. This results in

a linear system of two linear equations in the two unknowns

`{{math|''c''1}}`{=mediawiki} and

`{{math|''c''2}}`{=mediawiki}. Solving this system gives the

solution for a so-called Cauchy problem, in which the values at

`{{math|0}}`{=mediawiki} for the solution of the DEQ and its derivative

are specified.



\paragraph{Non-homogeneous equation with constant coefficients}



A non-homogeneous equation of order `{{mvar|n}}`{=mediawiki} with

constant coefficients may be written



$$y^{(n)}(x) + a_1 y^{(n-1)}(x) + \cdots + a_{n-1} y'(x)+ a_ny(x) = f(x),$$

where `{{math|''a''1, …, ''a''''n''}}`{=mediawiki}

are real or complex numbers, `{{mvar|f}}`{=mediawiki} is a given

function of `{{mvar|x}}`{=mediawiki}, and `{{mvar|y}}`{=mediawiki} is

the unknown function (for sake of simplicity,

\""`{{math|(''x'')}}`{=mediawiki}\"" will be omitted in the following).



There are several methods for solving such an equation. The best method

depends on the nature of the function `{{mvar|f}}`{=mediawiki} that

makes the equation non-homogeneous. If `{{mvar|f}}`{=mediawiki} is a

linear combination of exponential and sinusoidal functions, then the

exponential response formula may be used. If, more generally,

`{{mvar|f}}`{=mediawiki} is a linear combination of functions of the

form `{{math|''x''''n''''e''''ax''}}`{=mediawiki},

`{{math|''x''''n'' cos(''ax'')}}`{=mediawiki}, and

`{{math|''x''''n'' sin(''ax'')}}`{=mediawiki}, where

`{{mvar|n}}`{=mediawiki} is a nonnegative integer, and

`{{mvar|a}}`{=mediawiki} a constant (which need not be the same in each

term), then the method of undetermined coefficients may be used. Still

more general, the annihilator method applies when

`{{mvar|f}}`{=mediawiki} satisfies a homogeneous linear differential

equation, typically, a holonomic function.



The most general method is the variation of constants, which is

presented here.



The general solution of the associated homogeneous equation



$$y^{(n)} + a_1 y^{(n-1)} + \cdots + a_{n-1} y'+ a_ny = 0$$ is



$$y=u_1y_1+\cdots+ u_ny_n,$$ where

`{{math|(''y''1, …, ''y''''n'')}}`{=mediawiki} is

a basis of the vector space of the solutions and

`{{math|''u''1, …, ''u''''n''}}`{=mediawiki} are

arbitrary constants. The method of variation of constants takes its name

from the following idea. Instead of considering

`{{math|''u''1, …, ''u''''n''}}`{=mediawiki} as

constants, they can considered as unknown functions that have to be

determined for making `{{mvar|y}}`{=mediawiki} a solution of the

non-homogeneous equation. For this purpose, one adds the constraints



$$\begin{aligned}

0 &= u'_1y_1 + u'_2y_2 + \cdots+u'_ny_n \\\0 &= u'_1y'_1 + u'_2y'_2 + \cdots + u'_n y'_n \\\  & \; \;\vdots \\\0 &= u'_1y^{(n-2)}_1+u'_2y^{(n-2)}_2 + \cdots + u'_n y^{(n-2)}_n,

\end{aligned}$$ which imply (by product rule and induction)



$$y^{(i)} = u_1 y_1^{(i)} + \cdots + u_n y_n^{(i)}$$ for

`{{math|1=''i'' = 1, …, ''n'' – 1}}`{=mediawiki}, and



$$y^{(n)} = u_1 y_1^{(n)} + \cdots + u_n y_n^{(n)} +u'_1y_1^{(n-1)}+u'_2y_2^{(n-1)}+\cdots+u'_ny_n^{(n-1)}.$$



Replacing in the original equation `{{mvar|y}}`{=mediawiki} and its

derivatives by these expressions, and using the fact that

`{{math|''y''1, …, ''y''''n''}}`{=mediawiki} are

solutions of the original homogeneous equation, one gets



$$f=u'_1y_1^{(n-1)} + \cdots + u'_ny_n^{(n-1)}.$$



This equation and the above ones with `{{math|0}}`{=mediawiki} as

left-hand side form a system of `{{mvar|n}}`{=mediawiki} linear

equations in

`{{math|''u'''1, …, ''u'''''n''}}`{=mediawiki}

whose coefficients are known functions (`{{mvar|f}}`{=mediawiki}, the

$y_i$, and their derivatives). This system can be solved by any method

of linear algebra. The computation of antiderivatives gives

`{{math|''u''1, …, ''u''''n''}}`{=mediawiki}, and

then

`{{math|1=''y'' = ''u''1''y''1 + ? + ''u''''n''''y''''n''}}`{=mediawiki}.



As antiderivatives are defined up to the addition of a constant, one

finds again that the general solution of the non-homogeneous equation is

the sum of an arbitrary solution and the general solution of the

associated homogeneous equation.



\paragraph{First-order linear equation}



The general form of a linear ordinary differential equation of order 1,

after dividing out the coefficient of

`{{math|''y'''(''x'')}}`{=mediawiki}, is:



$$y'(x) = f(x) y(x) + g(x).$$



If the equation is homogeneous, i.e.

`{{math|1=''g''(''x'') = 0}}`{=mediawiki}, one may rewrite and

integrate:



$$\frac{y'}{y}= f, \qquad \log y = k +F,$$ where

`{{mvar|k}}`{=mediawiki} is an arbitrary constant of integration and

$F=\textstyle\int f\,dx$ is any antiderivative of

`{{mvar|f}}`{=mediawiki}. Thus, the general solution of the homogeneous

equation is



$$y=ce^F,$$ where `{{math|1=''c'' = ''e''''k''}}`{=mediawiki}

is an arbitrary constant.



For the general non-homogeneous equation, one may multiply it by the

reciprocal `{{math|''e''-''F''}}`{=mediawiki} of a solution

of the homogeneous equation. This gives



$$y'e^{-F}-yfe^{-F}= ge^{-F}.$$ As

$-fe^{-F} = \frac{d}{dx} \left(e^{-F}\right),$ the product rule allows

rewriting the equation as



$$\frac{d}{dx}\left(ye^{-F}\right)= ge^{-F}.$$ Thus, the general

solution is



$$y=ce^F + e^F\int ge^{-F}dx,$$ where `{{mvar|c}}`{=mediawiki} is a

constant of integration, and `{{mvar|F}}`{=mediawiki} is any

antiderivative of `{{mvar|f}}`{=mediawiki} (changing of antiderivative

amounts to change the constant of integration).



\paragraph{Example}



Solving the equation



:   $y'(x) + \frac{y(x)}{x} = 3x.$



The associated homogeneous equation $y'(x) + \frac{y(x)}{x} = 0$ gives



$$\frac{y'}{y}=-\frac{1}{x},$$ that is



$$y=\frac{c}{x}.$$



Dividing the original equation by one of these solutions gives



:   $xy'+y=3x^2.$



That is



$$(xy)'=3x^2,$$



$$xy=x^3 +c,$$ and



$$y(x)=x^2+c/x.$$ For the initial condition



:   $y(1)=\alpha,$



one gets the particular solution



$$y(x)=x^2+\frac{\alpha-1}{x}.$$



\paragraph{Making Linear Equations from Non-Linear Equations}



Sometimes a non-linear equation, which is not solvable like this, can be

made linear, and more easily solvable, by applying a substitution.



\paragraph{Example}



$$y\frac{dy}{dx}+xy^2=5x$$



:   Let\'s make the following substitution:



$$v=y^2 \,$$



$$\frac{dv}{dx}=2y\frac{dy}{dx}$$



:   Plugging in, we get



$$\frac{1}{2}\frac{dv}{dx}+xv=5x$$



$$\frac{dv}{dx}+2xv=10x$$



:   We can then solve as a linear equation in \textit{v}, using the

    step-by-step method above:



<!-- -->



:   \textbf{Step 1}: Find the integrating factor:



$$e^{\int P(x)dx}$$



$$e^{\int 2x dx}=e^{x^2+C}$$



$$e^{\int P(x)dx}=Ce^{x^2}$$



:   Letting C=1 for convenience, we get $e^{x^2}$ as our integrating

    factor.



<!-- -->



:   \textbf{Step 2}: Multiply through



$$e^{x^2}\frac{dv}{dx}+e^{x^2}xv=10e^{x^2}x$$



:   \textbf{Step 3}: Recognize that the left hand is

    $\frac{d}{dx} e^{\int P(x)dx}v$



$$\frac{d}{dx} e^{x^2}v=10e^{x^2}x$$



:   \textbf{Step 4}: Integrate both sides w.r.t.\textit{x}.



$$\int (\frac{d}{dx} e^{x^2}v)dx=\int 10e^{x^2}xdx$$



$$e^{x^2}v=5e^{x^2}+C$$



:   \textbf{Step 5}: Solve for \textit{v}.



$$v=5+\frac{C}{e^{x^2}}$$



:   Now that we have \textit{v}, solve for \textit{y}.



$$v=y^2 \,$$



$$y^2=5+\frac{C}{e^{x^2}} \,$$



\paragraph{Higher order linear equations}



A linear ordinary equation of order one with variable coefficients may

be solved by quadrature, which means that the solutions may be expressed

in terms of integrals. This is not the case for order at least two. This

is the main result of Picard--Vessiot theory which was initiated by

Émile Picard and Ernest Vessiot, and whose recent developments are

called differential Galois theory.



The impossibility of solving by quadrature can be compared with the

Abel--Ruffini theorem, which states that an algebraic equation of degree

at least five cannot, in general, be solved by radicals. This analogy

extends to the proof methods and motivates the denomination of

differential Galois theory.



Similarly to the algebraic case, the theory allows deciding which

equations may be solved by quadrature, and if possible solving them.

However, for both theories, the necessary computations are extremely

difficult, even with the most powerful computers.



Nevertheless, the case of order two with rational coefficients has been

completely solved by Kovacic\'s algorithm.



\paragraph{Cauchy--Euler equation}



Cauchy--Euler equations are examples of equations of any order, with

variable coefficients, that can be solved explicitly. These are the

equations of the form



$$x^n y^{(n)}(x) + a_{n-1} x^{n-1} y^{(n-1)}(x) + \cdots + a_0 y(x) = 0,$$

where $a_0, \ldots, a_{n-1}$ are constant coefficients.



\paragraph{Licensing}



Content obtained and/or adapted from:



\begin{itemize}
    \item \href{https://en.wikipedia.org/wiki/Linear_differential_equation}{Linear differential equation,
\end{itemize}
    Wikipedia}

    under a CC BY-SA license

\begin{itemize}
    \item \href{https://en.wikibooks.org/wiki/Ordinary_Differential_Equations/First_Order_Linear_1}{First Order Linear 1, Wikibooks: Differential
\end{itemize}
    Equations}

    under a CC BY-SA license"

%%===============================================================
\section{Variation of Parameters (2nd Order)}
\label{sec:variation-of-parameters-2nd-order}
%%===============================================================

"In mathematics, \textbf{variation of parameters}, also known as **variation

of constants**, is a general method to solve inhomogeneous linear

ordinary differential equations.



For first-order inhomogeneous linear differential equations it is

usually possible to find solutions via integrating factors or

undetermined coefficients with considerably less effort, although those

methods leverage heuristics that involve guessing and do not work for

all inhomogeneous linear differential equations.



Variation of parameters extends to linear partial differential equations

as well, specifically to inhomogeneous problems for linear evolution

equations like the heat equation, wave equation, and vibrating plate

equation. In this setting, the method is more often known as Duhamel\'s

principle, named after Jean-Marie Duhamel (1797--1872) who first applied

the method to solve the inhomogeneous heat equation. Sometimes variation

of parameters itself is called Duhamel\'s principle and vice versa.



\paragraph{Intuitive explanation}



Consider the equation of the forced dispersionless spring, in suitable

units:



$$x''(t) + x(t) = F(t).$$ Here `{{math|''x''}}`{=mediawiki} is the

displacement of the spring from the equilibrium $x = 0$, and

`{{math|''F''(''t'')}}`{=mediawiki} is an external applied force that

depends on time. When the external force is zero, this is the

homogeneous equation (whose solutions are linear combinations of sines

and cosines, corresponding to the spring oscillating with constant total

energy).



We can construct the solution physically, as follows. Between times

$t=s$ and $t=s+ds$, the momentum corresponding to the solution has a net

change $F(s)\,ds$ (see: Impulse (physics)). A solution to the

inhomogeneous equation, at the present time

`{{math|''t'' > 0}}`{=mediawiki}, is obtained by linearly superposing

the solutions obtained in this manner, for `{{math|''s''}}`{=mediawiki}

going between 0 and `{{math|t}}`{=mediawiki}.



The homogeneous initial-value problem, representing a small impulse

$F(s)\,ds$ being added to the solution at time $t=s$, is



$$x''(t)+x(t)=0,\quad x(s)=0,\ x'(s)=F(s)\,ds.$$ The unique solution to

this problem is easily seen to be $x(t) = F(s)\sin(t-s)\,ds$. The linear

superposition of all of these solutions is given by the integral:



$$x(t) = \int_0^t F(s)\sin(t-s)\,ds.$$



To verify that this satisfies the required equation:



$$x'(t)=\int_0^t F(s)\cos(t-s)\,ds$$



$$x''(t) = F(t) - \int_0^tF(s)\sin(t-s)\,ds = F(t)-x(t),$$ as required

(see: Leibniz integral rule).



The general method of variation of parameters allows for solving an

inhomogeneous linear equation



$$Lx(t)=F(t)$$ by means of considering the second-order linear

differential operator \textit{L} to be the net force, thus the total impulse

imparted to a solution between time \textit{s} and \textit{s}+\textit{ds} is \textit{F}(\textit{s})\textit{ds}.

Denote by $x_s$ the solution of the homogeneous initial value problem



$$Lx(t)=0, \quad x(s)=0,\ x'(s)=F (s)\,ds.$$ Then a particular solution

of the inhomogeneous equation is



$$x (t)=\int_0^t x_s (t)\,ds,$$ the result of linearly superposing the

infinitesimal homogeneous solutions. There are generalizations to higher

order linear differential operators.



In practice, variation of parameters usually involves the fundamental

solution of the homogeneous problem, the infinitesimal solutions $x_s$

then being given in terms of explicit linear combinations of linearly

independent fundamental solutions. In the case of the forced

dispersionless spring, the kernel

$\sin(t-s)=\sin t\cos s - \sin s\cos t$ is the associated decomposition

into fundamental solutions.



\paragraph{Description of method (Second Order)}



Consider a general second order linear nonhomogeneous differential

equation whose coefficient functions $p$, $q$, and $g$ are continuous:



$\begin{aligned} \quad \frac{d^2y}{dt^2} + p(t) \frac{dy}{dt} + q(t) y = g(t) \end{aligned}$



The corresponding second order linear homogeneous differential equation

is thus $\frac{d^2y}{dt^2} + p(t) \frac{dy}{dt} + q(t) y = 0$. Suppose

that we know the general solution to the corresponding second order

linear homogeneous differential equation in terms of two functions

$y_1(t)$ and $y_2(t)$ which form a fundamental set of solutions to the

corresponding second order linear homogeneous differential equation, say

$y_h(t) = Cy_1(t) + Dy_2(t)$.



We will now replace the constants $C$ and $D$ with functions, $u_1(t)$

and $u_2(t)$ to get:



$\begin{aligned} \quad Y(t) = u_1(t) y_1(t) + u_2(t) y_2(t) \end{aligned}$



We then want to try and determine what functions $u_1(t)$ and $u_2(t)$

make $Y(t) = u_1(t)y_1(t) + y_2(t)y_2(t)$ a particular solution to our

original second order linear nonhomogeneous differential equation. We

first differentiate $y$ and apply the product rule where appropriate to

get:



$\begin{aligned} \quad y' = u_1'(t)y_1(t) + u_1(t)y_1'(t) + u_2'(t)y_2(t) + u_2(t)y_2'(t) \end{aligned}$



Now we will set the terms containing the derivatives of the functions

$u_1'(t)$ and $u_2'(t)$ to equal zero, that is

$u_1'(t)y_1(t) + u_2'(t)y_2(t) = 0$. Note that this is a rather hefty

assumption, however, this assumption is not rash as we\'re looking only

for a particular solution, namely one for which this property holds.

We\'ll see that making this assumption does not lead to any

contradictions, and so:



$\begin{aligned} \quad y' = u_1(t)y_1'(t) + u_2(t)y_2'(t) \end{aligned}$



We now differentiate again by applying the product rule where

appropriate to get the second derivative of $y$:



$\begin{aligned} \quad y'' = u_1'(t)y_1'(t) + u_1(t)y_1''(t) + u_2'(t)y_2'(t) + u_2(t)y_2''(t) \end{aligned}$



We will now plug in $y$, $y'$, and $y''$ into our second order linear

nonhomogeneous differential equation to get that:



$\begin{aligned} \quad \frac{d^2y}{dt^2} + p(t) \frac{dy}{dt} + q(t)y = g(t) \\\\ \quad \left [ u_1'(t)y_1'(t) + u_1(t)y_1''(t) + u_2'(t)y_2'(t) + u_2(t)y_2''(t) \right ] + p(t) \left [ u_1(t)y_1'(t) + u_2(t)y_2'(t) \right ] + q(t) \left [ u_1(t) y_1(t) + u_2(t) y_2(t) \right ] = g(t) \\\\ \quad u_1(t) \underbrace{\left [ y_1''(t) + p(t)y_1'(t) + q(t)y_1(t) \right ]}_{=0} + u_2(t) \underbrace{\left [ y_2''(t) + p(t)y_2'(t) + q(t)y_2(t) \right ]}_{=0} + \left [ u_1'(t)y_1'(t) + u_2'(t)y_2'(t) \right ] = g(t) \\\\ \quad u_1'(t)y_1'(t) + u_2'(t)y_2'(t) = g(t) \end{aligned}$



Now recall that we supposed that $u_1'(t)y_1(t) + u_2'(t)y_2(t) = 0$. To

solve for $u_1'(t)$ and $u_2'(t)$, then all we need to do is solve the

following system of equations:



$\begin{aligned} u_1'(t)y_1(t) + u_2'(t)y_2(t) = 0 \\\\ u_1'(t)y_1'(t) + u_2'(t)y_2'(t) = g(t) \end{aligned}$



Recall that a unique solution exists provided that the determinant

$\begin{vmatrix} y_1(t) & y_2(t)\\\\ y_1'(t) & y_2'(t) \end{vmatrix}$

is nonzero. But this determinant is identically the Wronskian

$W(y_1, y_2)$, and it is assumed that this Wronskian is nonzero since

$y_1$ and $y_2$ form the general solution of the corresponding second

order linear homogeneous differential equation, and so by apply

Cramer\'s rule, the values of $u_1'(t)$ and $u_2'(t)$ are:



$\begin{aligned} \quad u_1'(t) = \frac{\begin{vmatrix}0 & y_2(t)\\\\ g(t) & y_2'(t) \end{vmatrix}}{W(y_1, y_2)} = - \frac{y_2(t)g(t)}{W(y_1, y_2)} \quad , \quad u_2'(t) = \frac{\begin{vmatrix} y_1(t) & 0\\\\ y_1'(t) & g(t) \end{vmatrix}}{W(y_1, y_2)} = \frac{y_1(t)g(t)}{W(y_1, y_2)} \end{aligned}$



We now integrate both sides of each equation above, and for constants

$A$ and $B$, we get $u_1(t)$ and $u_2(t)$:



$\begin{aligned} \quad u_1(t) = -\int \frac{y_2(t)g(t)}{W(y_1, y_2)}  dt + A \quad , \quad u_2(t) = \int \frac{y_1(t)g(t)}{W(y_1, y_2)}  dt + B \end{aligned}$



Therefore, a particular solution to our second order linear

nonhomogeneous differential equation is:



$\begin{aligned} \quad Y(t) = u_1(t)y_1(t) + u_2(t)y_2(t) \\\\ \quad Y(t) = y_1(t) \left ( -\int \frac{y_2(t)g(t)}{W(y_1, y_2)}  dt + A \right ) + y_2(t) \left ( \int \frac{y_1(t)g(t)}{W(y_1, y_2)}  dt + B \right ) \\\\ \quad Y(t) = -y_1(t) \int \frac{y_2(t)g(t)}{W(y_1, y_2)}  dt + y_2(t) \int \frac{y_1(t)g(t)}{W(y_1, y_2)}  dt  \underbrace{-Ay_1(t) + By_2(t)}_{=0} \\\\ \quad Y(t) = -y_1(t) \int \frac{y_2(t)g(t)}{W(y_1, y_2)}  dt + y_2(t) \int \frac{y_1(t)g(t)}{W(y_1, y_2)}  dt \end{aligned}$



And finally, the general solution to our differential equation will be:



$\begin{aligned} \quad y = Cy_1(t) + Dy_2(t) + Y(t) \\\\ \quad y = Cy_1(t) + Dy_2(t) -y_1(t) \int \frac{y_2(t)g(t)}{W(y_1, y_2)}  dt + y_2(t) \int \frac{y_1(t)g(t)}{W(y_1, y_2)}  dt \end{aligned}$



Note that the method of variation of parameters is useful provided that

the general solution to the corresponding second order linear

homogeneous differential equation is easy to solve, and provided that

the two integrals in the formula above are relatively simply to compute.



\paragraph{Examples}



\paragraph{First-order equation}



$$y' + p(x)y = q(x)$$ The general solution of the corresponding

homogeneous equation (written below) is the complementary solution to

our original (inhomogeneous) equation:



:   $y' + p(x)y = 0$.



This homogeneous differential equation can be solved by different

methods, for example separation of variables:



$$\frac{d}{dx} y + p(x)y = 0$$



$$\frac{dy}{dx}=-p(x)y$$



$${dy \over y} = -{p(x)\,dx},$$



$$\int \frac{1}{ y} \, dy = -\int p(x) \, dx$$



$$\ln |y| = -\int p(x) \, dx + C$$



$$y = \pm e^{-\int p(x) \, dx +C } = C_0 e^{-\int p(x) \, dx}$$ The

complementary solution to our original equation is therefore:



$$y_c = C_0 e^{-\int p(x) \, dx}$$ Now we return to solving the

non-homogeneous equation:



:   $y' + p(x)y = q(x)$



Using the method variation of parameters, the particular solution is

formed by multiplying the complementary solution by an unknown function

\textit{C}(\textit{x}):



$$y_p = C(x) e^{-\int p(x) \, dx}$$ By substituting the particular

solution into the non-homogeneous equation, we can find \textit{C}(\textit{x}):



:   $C' (x) e^{-\int p(x) \, dx} - C(x) p(x) e^{-\int p(x) \, dx} + p(x) C(x) e^{-\int p(x) \, dx} = q(x)$



<!-- -->



:   $C' (x) e^{-\int p(x) \, dx} = q(x)$



<!-- -->



:   $C' (x) = q(x) e^{\int p(x) \, dx}$



<!-- -->



:   $C(x) =\int q(x) e^{\int p(x) \, dx} \, dx + C_1$



We only need a single particular solution, so we arbitrarily select

$C_1=0$ for simplicity. Therefore the particular solution is:



$$y_p =e^{-\int p(x) \, dx} \int q(x) e^{\int p(x) \, dx} \, dx$$ The

final solution of the differential equation is:



$$\begin{aligned}

y &= y_c + y_p\\\&=C_0 e^{-\int p(x) \, dx} + e^{-\int p(x) \, dx} \int q(x) e^{\int p(x) \, dx} \, dx

\end{aligned}$$



This recreates the method of integrating factors.



\paragraph{Specific second-order equation}



Let us solve



$$y''+4y'+4y = \cosh x$$



We want to find the general solution to the differential equation, that

is, we want to find solutions to the homogeneous differential equation



:   $y''+4y'+4y=0.$



The characteristic equation is:



:   $\lambda^2+4\lambda+4=(\lambda+2)^2=0$



Since $\lambda=-2$ is a repeated root, we have to introduce a factor of

\textit{x} for one solution to ensure linear independence: \textit{u}\textasciitilde{}1\textasciitilde{} = \textit{e}^-2\textit{x}^

and \textit{u}\textasciitilde{}2\textasciitilde{} = \textit{xe}^-2\textit{x}^. The Wronskian of these two functions is



:   $W=\begin{vmatrix}

      e^{-2x} \& xe^{-2x} \\\    -2e^{-2x} \& -e^{-2x}(2x-1)\\\    \end{vmatrix} = -e^{-2x}e^{-2x}(2x-1)+2xe^{-2x}e^{-2x} = e^{-4x}.$



Because the Wronskian is non-zero, the two functions are linearly

independent, so this is in fact the general solution for the homogeneous

differential equation (and not a mere subset of it).



We seek functions \textit{A}(\textit{x}) and \textit{B}(\textit{x}) so

\textit{A}(\textit{x})\textit{u}\textasciitilde{}1\textasciitilde{} + \textit{B}(\textit{x})\textit{u}\textasciitilde{}2\textasciitilde{} is a particular solution of the

non-homogeneous equation. We need only calculate the integrals



$$A(x) = - \int {1\over W} u_2(x) b(x)\,\mathrm dx, \; B(x) = \int {1 \over W} u_1(x)b(x)\,\mathrm dx$$



Recall that for this example



$$b(x) = \cosh x$$



That is,



$$A(x) = - \int {1\over e^{-4x}} xe^{-2x} \cosh x \,\mathrm dx = - \int xe^{2x}\cosh x \,\mathrm dx = -{1\over 18}e^x\left(9(x-1)+e^{2x}(3x-1)\right)+C_1$$



$$B(x) = \int {1 \over e^{-4x}} e^{-2x} \cosh x \,\mathrm dx = \int e^{2x}\cosh x\,\mathrm dx ={1\over 6}e^x\left(3+e^{2x}\right)+C_2$$



where $C_1$ and $C_2$ are constants of integration.



\paragraph{General second-order equation}



We have a differential equation of the form



$$u''+p(x)u'+q(x)u=f(x)$$



and we define the linear operator



$$L=D^2+p(x)D+q(x)$$



where \textit{D} represents the differential operator. We therefore have to

solve the equation $L u(x)=f(x)$ for $u(x)$, where $L$ and $f(x)$ are

known.



We must solve first the corresponding homogeneous equation:



$$u''+p(x)u'+q(x)u=0$$



by the technique of our choice. Once we\'ve obtained two linearly

independent solutions to this homogeneous differential equation (because

this ODE is second-order) --- call them \textit{u}\textasciitilde{}1\textasciitilde{} and \textit{u}\textasciitilde{}2\textasciitilde{} --- we can

proceed with variation of parameters.



Now, we seek the general solution to the differential equation $u_G(x)$

which we assume to be of the form



$$u_G(x)=A(x)u_1(x)+B(x)u_2(x).$$



Here, $A(x)$ and $B(x)$ are unknown and $u_1(x)$ and $u_2(x)$ are the

solutions to the homogeneous equation. (Observe that if $A(x)$ and

$B(x)$ are constants, then $Lu_G(x)=0$.) Since the above is only one

equation and we have two unknown functions, it is reasonable to impose a

second condition. We choose the following:



$$A'(x)u_1(x)+B'(x)u_2(x)=0.$$



Now,



$$\begin{aligned}

u_G'(x) &= \left (A(x)u_1(x)+B(x)u_2(x) \right )' \\\&= \left (A(x)u_1(x) \right )'+ \left (B(x)u_2(x) \right )'\\\&=A'(x)u_1(x)+A(x)u_1'(x)+B'(x)u_2(x)+B(x)u_2'(x)\\\&=A'(x)u_1(x)+B'(x)u_2(x)+A(x)u_1'(x)+B(x)u_2'(x) \\\&= A(x)u_1'(x)+B(x)u_2'(x)

\end{aligned}$$



Differentiating again (omitting intermediary steps)



$$u_G''(x)=A(x)u_1''(x)+B(x)u_2''(x)+A'(x)u_1'(x)+B'(x)u_2'(x).$$



Now we can write the action of \textit{L} upon \textit{u}\textasciitilde{}\textit{G}\textasciitilde{} as



$$Lu_G=A(x)Lu_1(x)+B(x)Lu_2(x)+A'(x)u_1'(x)+B'(x)u_2'(x).$$



Since \textit{u}\textasciitilde{}1\textasciitilde{} and \textit{u}\textasciitilde{}2\textasciitilde{} are solutions, then



$$Lu_G=A'(x)u_1'(x)+B'(x)u_2'(x).$$



We have the system of equations



$$\begin{bmatrix}

u_1(x)  & u_2(x) \\\u_1'(x) & u_2'(x) \end{bmatrix}

\begin{bmatrix}

A'(x) \\\B'(x)\end{bmatrix} =

\begin{bmatrix} 0 \\\\ f \end{bmatrix}.$$



Expanding,



$$\begin{bmatrix}

A'(x)u_1(x)+B'(x)u_2(x)\\\A'(x)u_1'(x)+B'(x)u_2'(x) \end{bmatrix}

= \begin{bmatrix} 0\\\\f\end{bmatrix}.$$



So the above system determines precisely the conditions



$$A'(x)u_1(x)+B'(x)u_2(x)=0.$$



$$A'(x)u_1'(x)+B'(x)u_2'(x)=Lu_G=f.$$



We seek \textit{A}(\textit{x}) and \textit{B}(\textit{x}) from these conditions, so, given



$$\begin{bmatrix}

u_1(x)  & u_2(x) \\\u_1'(x) & u_2'(x)

\end{bmatrix}

\begin{bmatrix}

A'(x) \\\B'(x)\end{bmatrix} =

\begin{bmatrix}

0\\\f\end{bmatrix}$$



we can solve for (\textit{A}'(\textit{x}), \textit{B}'(\textit{x}))^T^, so



$$\begin{bmatrix} A'(x) \\\\ B'(x) \end{bmatrix} =

\begin{bmatrix}

u_1(x)  & u_2(x) \\\u_1'(x) & u_2'(x)

\end{bmatrix}^{-1}

\begin{bmatrix} 0\\\\ f \end{bmatrix} =\frac{1}{W} \begin{bmatrix}

u_2'(x)  & -u_2(x) \\\-u_1'(x) & u_1(x) \end{bmatrix}

\begin{bmatrix} 0\\\\ f \end{bmatrix},$$



where \textit{W} denotes the Wronskian of \textit{u}\textasciitilde{}1\textasciitilde{} and \textit{u}\textasciitilde{}2\textasciitilde{}. (We know that \textit{W}

is nonzero, from the assumption that \textit{u}\textasciitilde{}1\textasciitilde{} and \textit{u}\textasciitilde{}2\textasciitilde{} are linearly

independent.) So,



$$\begin{aligned}

A'(x) &= - {1\over W} u_2(x) f(x), & B'(x) &= {1 \over W} u_1(x)f(x) \\\A(x)  &= - \int {1\over W} u_2(x) f(x)\,\mathrm dx, & B(x) &= \int {1 \over W} u_1(x)f(x)\,\mathrm dx

\end{aligned}$$



While homogeneous equations are relatively easy to solve, this method

allows the calculation of the coefficients of the general solution of

the \textit{in} homogeneous equation, and thus the complete general solution of

the inhomogeneous equation can be determined.



Note that $A(x)$ and $B(x)$ are each determined only up to an arbitrary

additive constant (the constant of integration). Adding a constant to

$A(x)$ or $B(x)$ does not change the value of $Lu_G(x)$ because the

extra term is just a linear combination of \textit{u}\textasciitilde{}1\textasciitilde{} and \textit{u}\textasciitilde{}2\textasciitilde{}, which is a

solution of $L$ by definition.



\paragraph{Higher Order}



The method of variation of parameters can also be applied to higher

order differential equations. The process can be derived similarly.

Suppose that we have a higher order differential equation of the

following form:



$\begin{aligned} \quad \frac{d^ny}{dt^n} + p_1(t) \frac{d^{n-1}}{dt^{n-1}} + ... + p_n(t) y = g(t) \end{aligned}$



We first solve the corresponding homogeneous differential equation to

get $y_h(t) = C_1 y_1(t) + C_2y_2(t) + ... + C_ny_n(t)$. The functions

$y_1(t)$, $y_2(t)$, ..., $y_n(t)$ form a fundamental set. Assume a

particular solution to the nonhomogeneous differential equation is of

the form:



$\begin{aligned} \quad Y(t) = u_1(t)y_1(t) + u_2(t)y_2(t) + ... + u_n(t)y_n(t) \end{aligned}$



We then solve the following system of equations for the functions

$u_1'(t)$, $u_2'(t)$, ..., $u_n'(t)$.



$\begin{aligned} \quad \left\{\begin{matrix} u_1'(t)y_1(t) + u_2'(t)y_2(t) + ... + u_n'(t)y_n(t) = 0 \\\\ u_1'(t)y_1'(t) + u_2'(t)y_2'(t) + ... + u_n'(t)y_n'(t) = 0 \\\\ \vdots \\\\ u_1'(t)y_1^{(n-1)}(t) + u_2'(t)y_2^{(n-1)}(t) + ... + u_n'(t)y_n^{(n-1)}(t) = g(t) \end{matrix}\right. \end{aligned}$



Once again, a unique solution is guaranteed since $y_1(t)$, $y_2(t)$,

..., $y_n(t)$ form a fundamental set of solutions implies the Wronskian

$W(y_1, y_2, ..., y_n)$ is nonzero. Furthermore, it should be noted that

the system above can be solved for with row reduction (if the process is

simple) or more commonly by applying Cramer\'s rule once again.



We then integrate each of the functions $u_1'(t)$, $u_2'(t)$, ...,

$u_n'(t)$ to obtain $u_1(t)$, $u_2(t)$, ..., $u_n(t)$. Lastly, we obtain

that $Y(t) = u_1(t)y_1(t) + u_2(t)y_2(t) + ... + u_n(t)y_n(t)$ as our

particular solution.



We will now look at an example of using the method of variation of

parameters for higher order nonhomogeneous differential equations.



\paragraph{Example 1}



``{=html}Solve the third order linear nonhomogeneous

differential equation $y''' + y' = \tan t$ using the method of variation

of parameters.``{=html}



For the corresponding homogeneous differential equation $y''' + y' = 0$

we have that the characteristic equation is $r^3 + r = 0$ which can be

factored as $r(r^2 + 1) = 0$ and so $r = 0$ or $r = \pm i$. The solution

to the corresponding homogeneous differential equation is therefore:



$\begin{aligned} \quad y(t) = P + Q\cos t + R \sin t \end{aligned}$



Furthermore, $y_1 = e^{0t} = 1$, $y_2 = e^{0t} \cos t= \cos t$, and

$y_3 = e^{0t} \sin t = \sin t$.



Now let\'s try to find a particular solution to this differential

equation. We have that:



$\begin{aligned} \quad y_p = u_1y_1 + u_2y_2 + u_3y_3 \\\\ \quad y_p = u_1 + u_2 \cos t + u_3 \sin t \end{aligned}$



Thus we want to solve the following system of equations:



$\begin{aligned} \quad u_1' (1) + u_2' \cos t + u_3' \sin t = 0 \\\\ \quad u_1' (0) + u_2' (-\sin t) + u_3' (\cos t) = 0 \\\\ \quad u_1' (0) + u_2' (\cos t) + u_3' (\sin t) = \tan t \end{aligned}$



We will use Cramer\'s rule in order to solve for $u_1'$, $u_2'$, and

$u_3'$. We first find the corresponding Wronskian:



$\begin{aligned} \quad W = \begin{bmatrix} 1 & \cos t & \sin t \\\\ 0 & -\sin t & \cos t \\\\ 0 & -\cos t & -\sin t \end{bmatrix} = \begin{bmatrix} -\sin t & \cos t \\\\ -\cos t & -\sin t \end{bmatrix} = 1 \end{aligned}$



Now we have that:



$\begin{aligned} \quad u_1' = \frac{\begin{bmatrix} 0 & \cos t & \sin t \\\\ 0 & -\sin t & \cos t \\\\ \tan t & -\cos t & -\sin t \end{bmatrix}}{1} = \tan t \begin{bmatrix} \cos t & \sin t \\\\ -\sin t & \cos t \end{bmatrix} = \tan t \end{aligned}$



$\begin{aligned} \quad u_2' = \frac{\begin{bmatrix} 1 & 0 & \sin t \\\\ 0 & 0 & \cos t \\\\ 0 & \tan t & -\sin t \end{bmatrix}}{1} = - \begin{bmatrix} 1 & 0 & \sin t \\\\ 0 & \tan t & -\sin t \\\\ 0 & 0 & \cos t \end{bmatrix} = -\tan t \cos t = -\frac{\sin t}{\cos t} \cos t = -\sin t \end{aligned}$



$\begin{aligned} \quad u_3' = \frac{\begin{bmatrix} 1 & \cos t & 0 \\\\ 0 & -\sin t & 0 \\\\ 0 & -\cos t & \tan t \end{bmatrix}}{1} = \begin{bmatrix} -\sin t & 0 \\\\ -\cos t & \tan t \end{bmatrix} = -\sin t \tan t = -\frac{\sin^2 t}{\cos t} \end{aligned}$



We will now integrate $u_1'$, $u_2'$, and $u_3'$ to get:



$\begin{aligned} \quad \int u_1'(t)   dt = \int \tan t   dt = \ln \mid \sec t \mid + C \end{aligned}$



$\begin{aligned} \quad \int u_2'(t)   dt = \int - \sin t   dt = \cos t + D \end{aligned}$



$\begin{aligned} \quad \int u_3'(t)   dt = \int -\frac{\sin^2 t}{\cos t}   dt = \int \frac{\cos^2 t - 1}{\cos t}   dt = \int \left ( \cos t - \sec t \right )   dt = \sin t + \ln \mid \sec t + \tan t \mid + E \end{aligned}$



Thus we have that:



$\begin{aligned} \quad y(t) = \left ( \ln \mid \sec t \mid + C \right ) (1) + \left ( \cos t + D \right ) \cos t + \left ( \sin t + \ln \mid \sec t + \tan t \mid + E \right ) \sin t \end{aligned}$



\paragraph{Licensing}



Content obtained and/or adapted from:



\begin{itemize}
    \item \href{https://en.wikipedia.org/wiki/Variation_of_parameters}{Variation of parameters,
\end{itemize}
    Wikipedia}

    under a CC BY-SA license

\begin{itemize}
    \item \href{https://en.wikipedia.org/wiki/Variation_of_parameters}{The Method of Variation of Parameters,
\end{itemize}
    mathonline.wikidot.com}

    under a CC BY-SA license

\begin{itemize}
    \item \href{http://mathonline.wikidot.com/the-method-of-variation-of-parameters-for-higher-order-nonho}{Variation of Parameters for Higher Order,
\end{itemize}
    mathonline.wikidot.com}

    under a CC BY-SA license"

%%===============================================================
\section{Method of Undetermined Coefficients (2nd Order)}
\label{sec:method-of-undetermined-coefficients-2nd-order}
%%===============================================================

"In mathematics, the \textbf{method of undetermined coefficients} is an

approach to finding a particular solution to certain nonhomogeneous

ordinary differential equations and recurrence relations. It is closely

related to the annihilator method, but instead of using a particular

kind of differential operator (the annihilator) in order to find the

best possible form of the particular solution, a \""guess\"" is made as to

the appropriate form, which is then tested by differentiating the

resulting equation. For complex equations, the annihilator method or

variation of parameters is less time-consuming to perform.



Undetermined coefficients is not as general a method as variation of

parameters, since it only works for differential equations that follow

certain forms.



\paragraph{Description of the method}



Consider a linear non-homogeneous ordinary differential equation of the

form



$$\sum_{i=0}^n c_i y^{(i)} + y^{(n+1)} = g(x)$$



:   where $y^{(i)}$ denotes the i-th derivative of $y$, and $c_i$

    denotes a function of $x$.



The method of undetermined coefficients provides a straightforward

method of obtaining the solution to this ODE when two criteria are met:



1.  $c_i$ are constants.

2.  \textit{g}(\textit{x}) is a constant, a polynomial function, exponential function

    $e^{\alpha x}$, sine or cosine functions $\sin{\beta x}$ or

    $\cos{\beta x}$, or finite sums and products of these functions

    (${\alpha}$, ${\beta}$ constants).



The method consists of finding the general homogeneous solution $y_c$

for the complementary linear homogeneous differential equation



$$\sum_{i=0}^n c_i y^{(i)} + y^{(n+1)} = 0,$$



and a particular integral $y_p$ of the linear non-homogeneous ordinary

differential equation based on $g(x)$. Then the general solution $y$ to

the linear non-homogeneous ordinary differential equation would be



$$y = y_c + y_p.$$



If $g(x)$ consists of the sum of two functions $h(x) + w(x)$ and we say

that $y_{p_1}$ is the solution based on $h(x)$ and $y_{p_2}$ the

solution based on $w(x)$. Then, using a superposition principle, we can

say that the particular integral $y_p$ is



$$y_p = y_{p_1} + y_{p_2}.$$



\paragraph{Typical forms of the particular integral}



In order to find the particular integral, we need to \'guess\' its form,

with some coefficients left as variables to be solved for. This takes

the form of the first derivative of the complementary function. Below is

a table of some typical functions and the solution to guess for them.



  Function of \textit{x}                                                                                                           Form for \textit{y}

  ------------------------------------------------------------------------------------------------------------------------- ------------------------------------------------------------------------------------------------------------------

  $k e^{a x}\!$                                                                                                             $C e^{a x}\!$

  $k x^n, \; n = 0, 1, 2,\ldots\!$                                                                                           $\sum_{i=0}^n K_i x^i \!$

  $k \cos(a x) \text{ or } k \sin(a x) \!$                                                                                  $K \cos(a x) + M \sin(a x) \!$

  $k e^{a x} \cos(b x) \text{ or } ke^{a x} \sin(b x) \!$                                                                   $e^{a x} (K \cos(b x) + M \sin(b x)) \!$

  $\left(\sum_{i=0}^n k_i x^i\right) \cos(b x) \text{ or }\ \left(\sum_{i=0}^n k_i x^i\right) \sin(b x) \!$                 $\left(\sum_{i=0}^n Q_i x^i\right) \cos(b x) + \left(\sum_{i=0}^n R_i x^i\right) \sin(b x)$

  $\left(\sum_{i=0}^n k_i x^i\right) e^{a x} \cos(b x) \text{ or } \left(\sum_{i=0}^n k_i x^i\right) e^{a x} \sin(b x)\!$   $e^{a x} \left(\left(\sum_{i=0}^n Q_i x^i\right) \cos(b x) + \left(\sum_{i=0}^n R_i x^i\right) \sin(b x)\right)$



If a term in the above particular integral for \textit{y} appears in the

homogeneous solution, it is necessary to multiply by a sufficiently

large power of \textit{x} in order to make the solution independent. If the

function of \textit{x} is a sum of terms in the above table, the particular

integral can be guessed using a sum of the corresponding terms for \textit{y}.



\paragraph{Examples}



\paragraph{Example 1}



Find a particular integral of the equation



$$y'' + y = t \cos t.$$



The right side \textit{t} cos \textit{t} has the form



$$P_n e^{\alpha t} \cos{\beta t}$$



with \textit{n} = 2, \textit{a} = 0, and \textit{ß} = 1.



Since \textit{a} + \textit{iß} = \textit{i} is \textit{a simple root} of the characteristic equation



$$\lambda^2 + 1 = 0$$



we should try a particular integral of the form



$$\begin{aligned}

y_p &= t \left [F_1 (t) e^{\alpha t} \cos{\beta t} + G_1 (t) e^{\alpha t} \sin{\beta t} \right ] \\\&= t \left [F_1 (t) \cos t + G_1 (t) \sin t \right ] \\\&= t \left [ \left (A_0 t + A_1 \right ) \cos t + \left (B_0 t + B_1 \right ) \sin t \right ] \\\&= \left (A_0 t^2 + A_1 t \right ) \cos t + \left (B_0 t^2 + B_1 t \right) \sin t.

\end{aligned}$$



Substituting \textit{y}\textasciitilde{}\textit{p}\textasciitilde{} into the differential equation, we have the

identity



$$\begin{aligned}

t \cos t &= y_p'' + y_p \\\&= \left  [ \left(A_0 t^2 + A_1 t \right ) \cos t + \left (B_0 t^2 + B_1 t \right ) \sin t \right ]'' + \left[\left(A_0 t^2 + A_1 t \right ) \cos t + \left(B_0 t^2 + B_1 t \right ) \sin t \right ] \\\&=  \left [2A_0 \cos t + 2 \left (2A_0 t + A_1 \right )(-\sin t) + \left (A_0 t^2 + A_1 t \right )(-\cos t) + 2B_0 \sin t + 2 \left (2B_0 t + B_1 \right ) \cos t + \left (B_0 t^2 + B_1 t \right )(- \sin t) \right ] \\\&\qquad +\left[\left(A_0 t^2 + A_1 t \right ) \cos t + \left(B_0 t^2 + B_1 t \right ) \sin t \right ] \\\&= [4B_0 t + (2A_0 + 2B_1)] \cos t + [-4A_0 t + (-2A_1 + 2B_0)] \sin t.

\end{aligned}$$



Comparing both sides, we have



$$\begin{cases} 

1 = 4B_0\\\0 = 2A_0 + 2B_1  \\\0 = -4A_0 \\\0 = -2A_1 + 2B_0

\end{cases}$$



which has the solution



$$A_0 = 0, \quad A_1 = B_0 = \frac{1}{4}, \quad  B_1 = 0.$$



We then have a particular integral



$$y_p = \frac {1} {4} t \cos t + \frac {1}{4} t^2 \sin t.$$



\paragraph{Example 2}



Consider the following linear nonhomogeneous differential equation:



:   $\frac{dy}{dx} = y + e^x.$



This is like the first example above, except that the nonhomogeneous

part ($e^x$) is \textit{not} linearly independent to the general solution of

the homogeneous part ($c_1 e^x$); as a result, we have to multiply our

guess by a sufficiently large power of \textit{x} to make it linearly

independent.



Here our guess becomes:



$$y_p = A x e^x.$$



By substituting this function and its derivative into the differential

equation, one can solve for \textit{A}:



$$\frac{d}{dx} \left( A x e^x \right) = A x e^x + e^x$$



$$A x e^x + A e^x = A x e^x + e^x$$



$$A = 1.$$



So, the general solution to this differential equation is:



$$y = c_1 e^x + xe^x.$$



\paragraph{Example 3}



Find the general solution of the equation:



$$\frac{dy}{dt} = t^2 - y$$



$t^2$ is a polynomial of degree 2, so we look for a solution using the

same form,



$$y_p = A t^2 + B t + C,$$



Plugging this particular function into the original equation yields,



$$2 A t + B = t^2 - (A t^2 + B t + C),$$



$$2 A t + B =(1-A)t^2 -Bt -C,$$



$$(A-1)t^2 + (2A+B)t + (B+C) = 0.$$



which gives:



$$A-1 = 0, \quad 2A+B =0, \quad  B+C=0.$$



Solving for constants we get:



$$y_p = t^2 - 2 t + 2$$



To solve for the general solution,



$$y= y_p + y_c$$



where $y_c$ is the homogeneous solution $y_c = c_1 e^{-t}$, therefore,

the general solution is:



$$y= t^2 - 2 t + 2 + c_1 e^{-t}$$



\paragraph{Higher Order}



If we have a second order linear nonhomogenous differential equation

whose coefficients were constant, that is

$a \frac{d^2y}{dt^2} + b \frac{dy}{dt} + cy = g(t)$, then to solve this

differential equation, all we need to do is solve the corresponding

second order linear homogenous differential equation

$a \frac{d^2y}{dt^2} + b \frac{dy}{dt} + cy = 0$, and then find a

partial solution by assuming the form of the particular solution. More

precisely, if $g(t)$ was a polynomial, exponential, or sine/cosine

function (or a combination of these), then we could assume a form for

the particular solutions (see the linked page above for more details)

and solve for the coefficients of this form to obtain a particular

solution.



We will now look at this method for higher order linear nonhomogenous

differential equations. Consider the following $n^{\mathrm{th}}$ order

linear nonhomogenous differential equation with coefficients

$a_0, a_1, ..., a_n \in \mathbb{R}$:



$\begin{aligned} \quad a_0 \frac{d^ny}{dt^n} + a_1 \frac{d^{n-1}y}{dt^{n-1}} + ... + a_{n-1}\frac{dy}{dt} + a_n y = g(t) \end{aligned}$



Suppose that $g(t)$ is of a form containing a polynomial, exponential

function, or a sine/cosine function (like with when we were dealing with

the method of undetermined coefficients for second order linear

nonhomogenous differential equations). We can then find a particular

solution $Y(t)$ if we extend the assumed forms of combinations of

polynomials, exponential functions, or sine/cosine functions, and

multiplying by powers of $t$ to ensure our particular solution does not

contain part of a solution to the corresponding higher order linear

homogenous differential equation. We can then solve for these constants

by plugging our assumed form $Y(t)$ into our differential equation

above.



Providing a list of possible forms is trivial, so we will instead look

at some examples of apply this method.



\paragraph{Example 1}



``{=html}Find the general solution to the differential equation

$\frac{d^3y}{dt^3} + \frac{d^2y}{dt^2} + \frac{dy}{dt} + y = e^{-t} + 4t$.``{=html}



We will need to first solve the corresponding third order linear

homogenous differential equation

$\frac{d^3y}{dt^3} + \frac{d^2y}{dt^2} + \frac{dy}{dt} + y = 0$. This

characteristic equation to this differential equation is:



$\begin{aligned} \quad r^3 + r^2 + r + 1 = 0 \end{aligned}$



We can (by trial and error) see that $r_1 = -1$ is a solution to this

characteristic equation. Applying long division with the factor

$(r + 1)$ and we have that our characteristic equation can be written

as:



$\begin{aligned} \quad (r + 1)(r^2 + 1) = 0 \end{aligned}$



Therefore we can see that $r_2 = i$ and $r_3 = -i$. Therefore the

general solution to our third order linear homogenous differential

equation is:



$\begin{aligned} \quad y_h(t) = C_1e^{-t} + C_2\cos(t) + C_3\sin(t) \end{aligned}$



Now we note that $g(t) = e^{-t} + 4t$ has an exponential term and a

cosine term, so we expect the form of our particular solution to be

$Y(t) = Ae^{-t} ...$. We now compute the first, second, and third

derivatives of $Y$.



\paragraph{Licensing}



Content obtained and/or adapted from:



\begin{itemize}
    \item \href{https://en.wikipedia.org/wiki/Method_of_undetermined_coefficients}{Method of undetermined coefficients,
\end{itemize}
    Wikipedia}

    under a CC BY-SA license"

%%===============================================================
\section{Non-linear 2nd Order ODEs}
\label{sec:non-linear-2nd-order-odes}
%%===============================================================

Non-linear 2nd Order ODEs

%%===============================================================
\section{Variation of Parameters (Higher Order)}
\label{sec:variation-of-parameters-higher-order}
%%===============================================================

"In mathematics, \textbf{variation of parameters}, also known as **variation

of constants**, is a general method to solve inhomogeneous linear

ordinary differential equations.



For first-order inhomogeneous linear differential equations it is

usually possible to find solutions via integrating factors or

undetermined coefficients with considerably less effort, although those

methods leverage heuristics that involve guessing and do not work for

all inhomogeneous linear differential equations.



Variation of parameters extends to linear partial differential equations

as well, specifically to inhomogeneous problems for linear evolution

equations like the heat equation, wave equation, and vibrating plate

equation. In this setting, the method is more often known as Duhamel\'s

principle, named after Jean-Marie Duhamel (1797--1872) who first applied

the method to solve the inhomogeneous heat equation. Sometimes variation

of parameters itself is called Duhamel\'s principle and vice versa.



\paragraph{Intuitive explanation}



Consider the equation of the forced dispersionless spring, in suitable

units:



$$x''(t) + x(t) = F(t).$$ Here `{{math|''x''}}`{=mediawiki} is the

displacement of the spring from the equilibrium $x = 0$, and

`{{math|''F''(''t'')}}`{=mediawiki} is an external applied force that

depends on time. When the external force is zero, this is the

homogeneous equation (whose solutions are linear combinations of sines

and cosines, corresponding to the spring oscillating with constant total

energy).



We can construct the solution physically, as follows. Between times

$t=s$ and $t=s+ds$, the momentum corresponding to the solution has a net

change $F(s)\,ds$ (see: Impulse (physics)). A solution to the

inhomogeneous equation, at the present time

`{{math|''t'' > 0}}`{=mediawiki}, is obtained by linearly superposing

the solutions obtained in this manner, for `{{math|''s''}}`{=mediawiki}

going between 0 and `{{math|t}}`{=mediawiki}.



The homogeneous initial-value problem, representing a small impulse

$F(s)\,ds$ being added to the solution at time $t=s$, is



$$x''(t)+x(t)=0,\quad x(s)=0,\ x'(s)=F(s)\,ds.$$ The unique solution to

this problem is easily seen to be $x(t) = F(s)\sin(t-s)\,ds$. The linear

superposition of all of these solutions is given by the integral:



$$x(t) = \int_0^t F(s)\sin(t-s)\,ds.$$



To verify that this satisfies the required equation:



$$x'(t)=\int_0^t F(s)\cos(t-s)\,ds$$



$$x''(t) = F(t) - \int_0^tF(s)\sin(t-s)\,ds = F(t)-x(t),$$ as required

(see: Leibniz integral rule).



The general method of variation of parameters allows for solving an

inhomogeneous linear equation



$$Lx(t)=F(t)$$ by means of considering the second-order linear

differential operator \textit{L} to be the net force, thus the total impulse

imparted to a solution between time \textit{s} and \textit{s}+\textit{ds} is \textit{F}(\textit{s})\textit{ds}.

Denote by $x_s$ the solution of the homogeneous initial value problem



$$Lx(t)=0, \quad x(s)=0,\ x'(s)=F (s)\,ds.$$ Then a particular solution

of the inhomogeneous equation is



$$x (t)=\int_0^t x_s (t)\,ds,$$ the result of linearly superposing the

infinitesimal homogeneous solutions. There are generalizations to higher

order linear differential operators.



In practice, variation of parameters usually involves the fundamental

solution of the homogeneous problem, the infinitesimal solutions $x_s$

then being given in terms of explicit linear combinations of linearly

independent fundamental solutions. In the case of the forced

dispersionless spring, the kernel

$\sin(t-s)=\sin t\cos s - \sin s\cos t$ is the associated decomposition

into fundamental solutions.



\paragraph{Description of method (Second Order)}



Consider a general second order linear nonhomogeneous differential

equation whose coefficient functions $p$, $q$, and $g$ are continuous:



$\begin{aligned} \quad \frac{d^2y}{dt^2} + p(t) \frac{dy}{dt} + q(t) y = g(t) \end{aligned}$



The corresponding second order linear homogeneous differential equation

is thus $\frac{d^2y}{dt^2} + p(t) \frac{dy}{dt} + q(t) y = 0$. Suppose

that we know the general solution to the corresponding second order

linear homogeneous differential equation in terms of two functions

$y_1(t)$ and $y_2(t)$ which form a fundamental set of solutions to the

corresponding second order linear homogeneous differential equation, say

$y_h(t) = Cy_1(t) + Dy_2(t)$.



We will now replace the constants $C$ and $D$ with functions, $u_1(t)$

and $u_2(t)$ to get:



$\begin{aligned} \quad Y(t) = u_1(t) y_1(t) + u_2(t) y_2(t) \end{aligned}$



We then want to try and determine what functions $u_1(t)$ and $u_2(t)$

make $Y(t) = u_1(t)y_1(t) + y_2(t)y_2(t)$ a particular solution to our

original second order linear nonhomogeneous differential equation. We

first differentiate $y$ and apply the product rule where appropriate to

get:



$\begin{aligned} \quad y' = u_1'(t)y_1(t) + u_1(t)y_1'(t) + u_2'(t)y_2(t) + u_2(t)y_2'(t) \end{aligned}$



Now we will set the terms containing the derivatives of the functions

$u_1'(t)$ and $u_2'(t)$ to equal zero, that is

$u_1'(t)y_1(t) + u_2'(t)y_2(t) = 0$. Note that this is a rather hefty

assumption, however, this assumption is not rash as we\'re looking only

for a particular solution, namely one for which this property holds.

We\'ll see that making this assumption does not lead to any

contradictions, and so:



$\begin{aligned} \quad y' = u_1(t)y_1'(t) + u_2(t)y_2'(t) \end{aligned}$



We now differentiate again by applying the product rule where

appropriate to get the second derivative of $y$:



$\begin{aligned} \quad y'' = u_1'(t)y_1'(t) + u_1(t)y_1''(t) + u_2'(t)y_2'(t) + u_2(t)y_2''(t) \end{aligned}$



We will now plug in $y$, $y'$, and $y''$ into our second order linear

nonhomogeneous differential equation to get that:



$\begin{aligned} \quad \frac{d^2y}{dt^2} + p(t) \frac{dy}{dt} + q(t)y = g(t) \\\\ \quad \left [ u_1'(t)y_1'(t) + u_1(t)y_1''(t) + u_2'(t)y_2'(t) + u_2(t)y_2''(t) \right ] + p(t) \left [ u_1(t)y_1'(t) + u_2(t)y_2'(t) \right ] + q(t) \left [ u_1(t) y_1(t) + u_2(t) y_2(t) \right ] = g(t) \\\\ \quad u_1(t) \underbrace{\left [ y_1''(t) + p(t)y_1'(t) + q(t)y_1(t) \right ]}_{=0} + u_2(t) \underbrace{\left [ y_2''(t) + p(t)y_2'(t) + q(t)y_2(t) \right ]}_{=0} + \left [ u_1'(t)y_1'(t) + u_2'(t)y_2'(t) \right ] = g(t) \\\\ \quad u_1'(t)y_1'(t) + u_2'(t)y_2'(t) = g(t) \end{aligned}$



Now recall that we supposed that $u_1'(t)y_1(t) + u_2'(t)y_2(t) = 0$. To

solve for $u_1'(t)$ and $u_2'(t)$, then all we need to do is solve the

following system of equations:



$\begin{aligned} u_1'(t)y_1(t) + u_2'(t)y_2(t) = 0 \\\\ u_1'(t)y_1'(t) + u_2'(t)y_2'(t) = g(t) \end{aligned}$



Recall that a unique solution exists provided that the determinant

$\begin{vmatrix} y_1(t) & y_2(t)\\\\ y_1'(t) & y_2'(t) \end{vmatrix}$

is nonzero. But this determinant is identically the Wronskian

$W(y_1, y_2)$, and it is assumed that this Wronskian is nonzero since

$y_1$ and $y_2$ form the general solution of the corresponding second

order linear homogeneous differential equation, and so by apply

Cramer\'s rule, the values of $u_1'(t)$ and $u_2'(t)$ are:



$\begin{aligned} \quad u_1'(t) = \frac{\begin{vmatrix}0 & y_2(t)\\\\ g(t) & y_2'(t) \end{vmatrix}}{W(y_1, y_2)} = - \frac{y_2(t)g(t)}{W(y_1, y_2)} \quad , \quad u_2'(t) = \frac{\begin{vmatrix} y_1(t) & 0\\\\ y_1'(t) & g(t) \end{vmatrix}}{W(y_1, y_2)} = \frac{y_1(t)g(t)}{W(y_1, y_2)} \end{aligned}$



We now integrate both sides of each equation above, and for constants

$A$ and $B$, we get $u_1(t)$ and $u_2(t)$:



$\begin{aligned} \quad u_1(t) = -\int \frac{y_2(t)g(t)}{W(y_1, y_2)}  dt + A \quad , \quad u_2(t) = \int \frac{y_1(t)g(t)}{W(y_1, y_2)}  dt + B \end{aligned}$



Therefore, a particular solution to our second order linear

nonhomogeneous differential equation is:



$\begin{aligned} \quad Y(t) = u_1(t)y_1(t) + u_2(t)y_2(t) \\\\ \quad Y(t) = y_1(t) \left ( -\int \frac{y_2(t)g(t)}{W(y_1, y_2)}  dt + A \right ) + y_2(t) \left ( \int \frac{y_1(t)g(t)}{W(y_1, y_2)}  dt + B \right ) \\\\ \quad Y(t) = -y_1(t) \int \frac{y_2(t)g(t)}{W(y_1, y_2)}  dt + y_2(t) \int \frac{y_1(t)g(t)}{W(y_1, y_2)}  dt  \underbrace{-Ay_1(t) + By_2(t)}_{=0} \\\\ \quad Y(t) = -y_1(t) \int \frac{y_2(t)g(t)}{W(y_1, y_2)}  dt + y_2(t) \int \frac{y_1(t)g(t)}{W(y_1, y_2)}  dt \end{aligned}$



And finally, the general solution to our differential equation will be:



$\begin{aligned} \quad y = Cy_1(t) + Dy_2(t) + Y(t) \\\\ \quad y = Cy_1(t) + Dy_2(t) -y_1(t) \int \frac{y_2(t)g(t)}{W(y_1, y_2)}  dt + y_2(t) \int \frac{y_1(t)g(t)}{W(y_1, y_2)}  dt \end{aligned}$



Note that the method of variation of parameters is useful provided that

the general solution to the corresponding second order linear

homogeneous differential equation is easy to solve, and provided that

the two integrals in the formula above are relatively simply to compute.



\paragraph{Examples}



\paragraph{First-order equation}



$$y' + p(x)y = q(x)$$ The general solution of the corresponding

homogeneous equation (written below) is the complementary solution to

our original (inhomogeneous) equation:



:   $y' + p(x)y = 0$.



This homogeneous differential equation can be solved by different

methods, for example separation of variables:



$$\frac{d}{dx} y + p(x)y = 0$$



$$\frac{dy}{dx}=-p(x)y$$



$${dy \over y} = -{p(x)\,dx},$$



$$\int \frac{1}{ y} \, dy = -\int p(x) \, dx$$



$$\ln |y| = -\int p(x) \, dx + C$$



$$y = \pm e^{-\int p(x) \, dx +C } = C_0 e^{-\int p(x) \, dx}$$ The

complementary solution to our original equation is therefore:



$$y_c = C_0 e^{-\int p(x) \, dx}$$ Now we return to solving the

non-homogeneous equation:



:   $y' + p(x)y = q(x)$



Using the method variation of parameters, the particular solution is

formed by multiplying the complementary solution by an unknown function

\textit{C}(\textit{x}):



$$y_p = C(x) e^{-\int p(x) \, dx}$$ By substituting the particular

solution into the non-homogeneous equation, we can find \textit{C}(\textit{x}):



:   $C' (x) e^{-\int p(x) \, dx} - C(x) p(x) e^{-\int p(x) \, dx} + p(x) C(x) e^{-\int p(x) \, dx} = q(x)$



<!-- -->



:   $C' (x) e^{-\int p(x) \, dx} = q(x)$



<!-- -->



:   $C' (x) = q(x) e^{\int p(x) \, dx}$



<!-- -->



:   $C(x) =\int q(x) e^{\int p(x) \, dx} \, dx + C_1$



We only need a single particular solution, so we arbitrarily select

$C_1=0$ for simplicity. Therefore the particular solution is:



$$y_p =e^{-\int p(x) \, dx} \int q(x) e^{\int p(x) \, dx} \, dx$$ The

final solution of the differential equation is:



$$\begin{aligned}

y &= y_c + y_p\\\&=C_0 e^{-\int p(x) \, dx} + e^{-\int p(x) \, dx} \int q(x) e^{\int p(x) \, dx} \, dx

\end{aligned}$$



This recreates the method of integrating factors.



\paragraph{Specific second-order equation}



Let us solve



$$y''+4y'+4y = \cosh x$$



We want to find the general solution to the differential equation, that

is, we want to find solutions to the homogeneous differential equation



:   $y''+4y'+4y=0.$



The characteristic equation is:



:   $\lambda^2+4\lambda+4=(\lambda+2)^2=0$



Since $\lambda=-2$ is a repeated root, we have to introduce a factor of

\textit{x} for one solution to ensure linear independence: \textit{u}\textasciitilde{}1\textasciitilde{} = \textit{e}^-2\textit{x}^

and \textit{u}\textasciitilde{}2\textasciitilde{} = \textit{xe}^-2\textit{x}^. The Wronskian of these two functions is



:   $W=\begin{vmatrix}

      e^{-2x} \& xe^{-2x} \\\    -2e^{-2x} \& -e^{-2x}(2x-1)\\\    \end{vmatrix} = -e^{-2x}e^{-2x}(2x-1)+2xe^{-2x}e^{-2x} = e^{-4x}.$



Because the Wronskian is non-zero, the two functions are linearly

independent, so this is in fact the general solution for the homogeneous

differential equation (and not a mere subset of it).



We seek functions \textit{A}(\textit{x}) and \textit{B}(\textit{x}) so

\textit{A}(\textit{x})\textit{u}\textasciitilde{}1\textasciitilde{} + \textit{B}(\textit{x})\textit{u}\textasciitilde{}2\textasciitilde{} is a particular solution of the

non-homogeneous equation. We need only calculate the integrals



$$A(x) = - \int {1\over W} u_2(x) b(x)\,\mathrm dx, \; B(x) = \int {1 \over W} u_1(x)b(x)\,\mathrm dx$$



Recall that for this example



$$b(x) = \cosh x$$



That is,



$$A(x) = - \int {1\over e^{-4x}} xe^{-2x} \cosh x \,\mathrm dx = - \int xe^{2x}\cosh x \,\mathrm dx = -{1\over 18}e^x\left(9(x-1)+e^{2x}(3x-1)\right)+C_1$$



$$B(x) = \int {1 \over e^{-4x}} e^{-2x} \cosh x \,\mathrm dx = \int e^{2x}\cosh x\,\mathrm dx ={1\over 6}e^x\left(3+e^{2x}\right)+C_2$$



where $C_1$ and $C_2$ are constants of integration.



\paragraph{General second-order equation}



We have a differential equation of the form



$$u''+p(x)u'+q(x)u=f(x)$$



and we define the linear operator



$$L=D^2+p(x)D+q(x)$$



where \textit{D} represents the differential operator. We therefore have to

solve the equation $L u(x)=f(x)$ for $u(x)$, where $L$ and $f(x)$ are

known.



We must solve first the corresponding homogeneous equation:



$$u''+p(x)u'+q(x)u=0$$



by the technique of our choice. Once we\'ve obtained two linearly

independent solutions to this homogeneous differential equation (because

this ODE is second-order) --- call them \textit{u}\textasciitilde{}1\textasciitilde{} and \textit{u}\textasciitilde{}2\textasciitilde{} --- we can

proceed with variation of parameters.



Now, we seek the general solution to the differential equation $u_G(x)$

which we assume to be of the form



$$u_G(x)=A(x)u_1(x)+B(x)u_2(x).$$



Here, $A(x)$ and $B(x)$ are unknown and $u_1(x)$ and $u_2(x)$ are the

solutions to the homogeneous equation. (Observe that if $A(x)$ and

$B(x)$ are constants, then $Lu_G(x)=0$.) Since the above is only one

equation and we have two unknown functions, it is reasonable to impose a

second condition. We choose the following:



$$A'(x)u_1(x)+B'(x)u_2(x)=0.$$



Now,



$$\begin{aligned}

u_G'(x) &= \left (A(x)u_1(x)+B(x)u_2(x) \right )' \\\&= \left (A(x)u_1(x) \right )'+ \left (B(x)u_2(x) \right )'\\\&=A'(x)u_1(x)+A(x)u_1'(x)+B'(x)u_2(x)+B(x)u_2'(x)\\\&=A'(x)u_1(x)+B'(x)u_2(x)+A(x)u_1'(x)+B(x)u_2'(x) \\\&= A(x)u_1'(x)+B(x)u_2'(x)

\end{aligned}$$



Differentiating again (omitting intermediary steps)



$$u_G''(x)=A(x)u_1''(x)+B(x)u_2''(x)+A'(x)u_1'(x)+B'(x)u_2'(x).$$



Now we can write the action of \textit{L} upon \textit{u}\textasciitilde{}\textit{G}\textasciitilde{} as



$$Lu_G=A(x)Lu_1(x)+B(x)Lu_2(x)+A'(x)u_1'(x)+B'(x)u_2'(x).$$



Since \textit{u}\textasciitilde{}1\textasciitilde{} and \textit{u}\textasciitilde{}2\textasciitilde{} are solutions, then



$$Lu_G=A'(x)u_1'(x)+B'(x)u_2'(x).$$



We have the system of equations



$$\begin{bmatrix}

u_1(x)  & u_2(x) \\\u_1'(x) & u_2'(x) \end{bmatrix}

\begin{bmatrix}

A'(x) \\\B'(x)\end{bmatrix} =

\begin{bmatrix} 0 \\\\ f \end{bmatrix}.$$



Expanding,



$$\begin{bmatrix}

A'(x)u_1(x)+B'(x)u_2(x)\\\A'(x)u_1'(x)+B'(x)u_2'(x) \end{bmatrix}

= \begin{bmatrix} 0\\\\f\end{bmatrix}.$$



So the above system determines precisely the conditions



$$A'(x)u_1(x)+B'(x)u_2(x)=0.$$



$$A'(x)u_1'(x)+B'(x)u_2'(x)=Lu_G=f.$$



We seek \textit{A}(\textit{x}) and \textit{B}(\textit{x}) from these conditions, so, given



$$\begin{bmatrix}

u_1(x)  & u_2(x) \\\u_1'(x) & u_2'(x)

\end{bmatrix}

\begin{bmatrix}

A'(x) \\\B'(x)\end{bmatrix} =

\begin{bmatrix}

0\\\f\end{bmatrix}$$



we can solve for (\textit{A}'(\textit{x}), \textit{B}'(\textit{x}))^T^, so



$$\begin{bmatrix} A'(x) \\\\ B'(x) \end{bmatrix} =

\begin{bmatrix}

u_1(x)  & u_2(x) \\\u_1'(x) & u_2'(x)

\end{bmatrix}^{-1}

\begin{bmatrix} 0\\\\ f \end{bmatrix} =\frac{1}{W} \begin{bmatrix}

u_2'(x)  & -u_2(x) \\\-u_1'(x) & u_1(x) \end{bmatrix}

\begin{bmatrix} 0\\\\ f \end{bmatrix},$$



where \textit{W} denotes the Wronskian of \textit{u}\textasciitilde{}1\textasciitilde{} and \textit{u}\textasciitilde{}2\textasciitilde{}. (We know that \textit{W}

is nonzero, from the assumption that \textit{u}\textasciitilde{}1\textasciitilde{} and \textit{u}\textasciitilde{}2\textasciitilde{} are linearly

independent.) So,



$$\begin{aligned}

A'(x) &= - {1\over W} u_2(x) f(x), & B'(x) &= {1 \over W} u_1(x)f(x) \\\A(x)  &= - \int {1\over W} u_2(x) f(x)\,\mathrm dx, & B(x) &= \int {1 \over W} u_1(x)f(x)\,\mathrm dx

\end{aligned}$$



While homogeneous equations are relatively easy to solve, this method

allows the calculation of the coefficients of the general solution of

the \textit{in} homogeneous equation, and thus the complete general solution of

the inhomogeneous equation can be determined.



Note that $A(x)$ and $B(x)$ are each determined only up to an arbitrary

additive constant (the constant of integration). Adding a constant to

$A(x)$ or $B(x)$ does not change the value of $Lu_G(x)$ because the

extra term is just a linear combination of \textit{u}\textasciitilde{}1\textasciitilde{} and \textit{u}\textasciitilde{}2\textasciitilde{}, which is a

solution of $L$ by definition.



\paragraph{Higher Order}



The method of variation of parameters can also be applied to higher

order differential equations. The process can be derived similarly.

Suppose that we have a higher order differential equation of the

following form:



$\begin{aligned} \quad \frac{d^ny}{dt^n} + p_1(t) \frac{d^{n-1}}{dt^{n-1}} + ... + p_n(t) y = g(t) \end{aligned}$



We first solve the corresponding homogeneous differential equation to

get $y_h(t) = C_1 y_1(t) + C_2y_2(t) + ... + C_ny_n(t)$. The functions

$y_1(t)$, $y_2(t)$, ..., $y_n(t)$ form a fundamental set. Assume a

particular solution to the nonhomogeneous differential equation is of

the form:



$\begin{aligned} \quad Y(t) = u_1(t)y_1(t) + u_2(t)y_2(t) + ... + u_n(t)y_n(t) \end{aligned}$



We then solve the following system of equations for the functions

$u_1'(t)$, $u_2'(t)$, ..., $u_n'(t)$.



$\begin{aligned} \quad \left\{\begin{matrix} u_1'(t)y_1(t) + u_2'(t)y_2(t) + ... + u_n'(t)y_n(t) = 0 \\\\ u_1'(t)y_1'(t) + u_2'(t)y_2'(t) + ... + u_n'(t)y_n'(t) = 0 \\\\ \vdots \\\\ u_1'(t)y_1^{(n-1)}(t) + u_2'(t)y_2^{(n-1)}(t) + ... + u_n'(t)y_n^{(n-1)}(t) = g(t) \end{matrix}\right. \end{aligned}$



Once again, a unique solution is guaranteed since $y_1(t)$, $y_2(t)$,

..., $y_n(t)$ form a fundamental set of solutions implies the Wronskian

$W(y_1, y_2, ..., y_n)$ is nonzero. Furthermore, it should be noted that

the system above can be solved for with row reduction (if the process is

simple) or more commonly by applying Cramer\'s rule once again.



We then integrate each of the functions $u_1'(t)$, $u_2'(t)$, ...,

$u_n'(t)$ to obtain $u_1(t)$, $u_2(t)$, ..., $u_n(t)$. Lastly, we obtain

that $Y(t) = u_1(t)y_1(t) + u_2(t)y_2(t) + ... + u_n(t)y_n(t)$ as our

particular solution.



We will now look at an example of using the method of variation of

parameters for higher order nonhomogeneous differential equations.



\paragraph{Example 1}



``{=html}Solve the third order linear nonhomogeneous

differential equation $y''' + y' = \tan t$ using the method of variation

of parameters.``{=html}



For the corresponding homogeneous differential equation $y''' + y' = 0$

we have that the characteristic equation is $r^3 + r = 0$ which can be

factored as $r(r^2 + 1) = 0$ and so $r = 0$ or $r = \pm i$. The solution

to the corresponding homogeneous differential equation is therefore:



$\begin{aligned} \quad y(t) = P + Q\cos t + R \sin t \end{aligned}$



Furthermore, $y_1 = e^{0t} = 1$, $y_2 = e^{0t} \cos t= \cos t$, and

$y_3 = e^{0t} \sin t = \sin t$.



Now let\'s try to find a particular solution to this differential

equation. We have that:



$\begin{aligned} \quad y_p = u_1y_1 + u_2y_2 + u_3y_3 \\\\ \quad y_p = u_1 + u_2 \cos t + u_3 \sin t \end{aligned}$



Thus we want to solve the following system of equations:



$\begin{aligned} \quad u_1' (1) + u_2' \cos t + u_3' \sin t = 0 \\\\ \quad u_1' (0) + u_2' (-\sin t) + u_3' (\cos t) = 0 \\\\ \quad u_1' (0) + u_2' (\cos t) + u_3' (\sin t) = \tan t \end{aligned}$



We will use Cramer\'s rule in order to solve for $u_1'$, $u_2'$, and

$u_3'$. We first find the corresponding Wronskian:



$\begin{aligned} \quad W = \begin{bmatrix} 1 & \cos t & \sin t \\\\ 0 & -\sin t & \cos t \\\\ 0 & -\cos t & -\sin t \end{bmatrix} = \begin{bmatrix} -\sin t & \cos t \\\\ -\cos t & -\sin t \end{bmatrix} = 1 \end{aligned}$



Now we have that:



$\begin{aligned} \quad u_1' = \frac{\begin{bmatrix} 0 & \cos t & \sin t \\\\ 0 & -\sin t & \cos t \\\\ \tan t & -\cos t & -\sin t \end{bmatrix}}{1} = \tan t \begin{bmatrix} \cos t & \sin t \\\\ -\sin t & \cos t \end{bmatrix} = \tan t \end{aligned}$



$\begin{aligned} \quad u_2' = \frac{\begin{bmatrix} 1 & 0 & \sin t \\\\ 0 & 0 & \cos t \\\\ 0 & \tan t & -\sin t \end{bmatrix}}{1} = - \begin{bmatrix} 1 & 0 & \sin t \\\\ 0 & \tan t & -\sin t \\\\ 0 & 0 & \cos t \end{bmatrix} = -\tan t \cos t = -\frac{\sin t}{\cos t} \cos t = -\sin t \end{aligned}$



$\begin{aligned} \quad u_3' = \frac{\begin{bmatrix} 1 & \cos t & 0 \\\\ 0 & -\sin t & 0 \\\\ 0 & -\cos t & \tan t \end{bmatrix}}{1} = \begin{bmatrix} -\sin t & 0 \\\\ -\cos t & \tan t \end{bmatrix} = -\sin t \tan t = -\frac{\sin^2 t}{\cos t} \end{aligned}$



We will now integrate $u_1'$, $u_2'$, and $u_3'$ to get:



$\begin{aligned} \quad \int u_1'(t)   dt = \int \tan t   dt = \ln \mid \sec t \mid + C \end{aligned}$



$\begin{aligned} \quad \int u_2'(t)   dt = \int - \sin t   dt = \cos t + D \end{aligned}$



$\begin{aligned} \quad \int u_3'(t)   dt = \int -\frac{\sin^2 t}{\cos t}   dt = \int \frac{\cos^2 t - 1}{\cos t}   dt = \int \left ( \cos t - \sec t \right )   dt = \sin t + \ln \mid \sec t + \tan t \mid + E \end{aligned}$



Thus we have that:



$\begin{aligned} \quad y(t) = \left ( \ln \mid \sec t \mid + C \right ) (1) + \left ( \cos t + D \right ) \cos t + \left ( \sin t + \ln \mid \sec t + \tan t \mid + E \right ) \sin t \end{aligned}$



\paragraph{Licensing}



Content obtained and/or adapted from:



\begin{itemize}
    \item \href{https://en.wikipedia.org/wiki/Variation_of_parameters}{Variation of parameters,
\end{itemize}
    Wikipedia}

    under a CC BY-SA license

\begin{itemize}
    \item \href{https://en.wikipedia.org/wiki/Variation_of_parameters}{The Method of Variation of Parameters,
\end{itemize}
    mathonline.wikidot.com}

    under a CC BY-SA license

\begin{itemize}
    \item \href{http://mathonline.wikidot.com/the-method-of-variation-of-parameters-for-higher-order-nonho}{Variation of Parameters for Higher Order,
\end{itemize}
    mathonline.wikidot.com}

    under a CC BY-SA license"

%%===============================================================
\section{Method of Undetermined Coefficients (Higher Order)}
\label{sec:method-of-undetermined-coefficients-higher-order}
%%===============================================================

"In mathematics, the \textbf{method of undetermined coefficients} is an

approach to finding a particular solution to certain nonhomogeneous

ordinary differential equations and recurrence relations. It is closely

related to the annihilator method, but instead of using a particular

kind of differential operator (the annihilator) in order to find the

best possible form of the particular solution, a \""guess\"" is made as to

the appropriate form, which is then tested by differentiating the

resulting equation. For complex equations, the annihilator method or

variation of parameters is less time-consuming to perform.



Undetermined coefficients is not as general a method as variation of

parameters, since it only works for differential equations that follow

certain forms.



\paragraph{Description of the method}



Consider a linear non-homogeneous ordinary differential equation of the

form



$$\sum_{i=0}^n c_i y^{(i)} + y^{(n+1)} = g(x)$$



:   where $y^{(i)}$ denotes the i-th derivative of $y$, and $c_i$

    denotes a function of $x$.



The method of undetermined coefficients provides a straightforward

method of obtaining the solution to this ODE when two criteria are met:



1.  $c_i$ are constants.

2.  \textit{g}(\textit{x}) is a constant, a polynomial function, exponential function

    $e^{\alpha x}$, sine or cosine functions $\sin{\beta x}$ or

    $\cos{\beta x}$, or finite sums and products of these functions

    (${\alpha}$, ${\beta}$ constants).



The method consists of finding the general homogeneous solution $y_c$

for the complementary linear homogeneous differential equation



$$\sum_{i=0}^n c_i y^{(i)} + y^{(n+1)} = 0,$$



and a particular integral $y_p$ of the linear non-homogeneous ordinary

differential equation based on $g(x)$. Then the general solution $y$ to

the linear non-homogeneous ordinary differential equation would be



$$y = y_c + y_p.$$



If $g(x)$ consists of the sum of two functions $h(x) + w(x)$ and we say

that $y_{p_1}$ is the solution based on $h(x)$ and $y_{p_2}$ the

solution based on $w(x)$. Then, using a superposition principle, we can

say that the particular integral $y_p$ is



$$y_p = y_{p_1} + y_{p_2}.$$



\paragraph{Typical forms of the particular integral}



In order to find the particular integral, we need to \'guess\' its form,

with some coefficients left as variables to be solved for. This takes

the form of the first derivative of the complementary function. Below is

a table of some typical functions and the solution to guess for them.



  Function of \textit{x}                                                                                                           Form for \textit{y}

  ------------------------------------------------------------------------------------------------------------------------- ------------------------------------------------------------------------------------------------------------------

  $k e^{a x}\!$                                                                                                             $C e^{a x}\!$

  $k x^n, \; n = 0, 1, 2,\ldots\!$                                                                                           $\sum_{i=0}^n K_i x^i \!$

  $k \cos(a x) \text{ or } k \sin(a x) \!$                                                                                  $K \cos(a x) + M \sin(a x) \!$

  $k e^{a x} \cos(b x) \text{ or } ke^{a x} \sin(b x) \!$                                                                   $e^{a x} (K \cos(b x) + M \sin(b x)) \!$

  $\left(\sum_{i=0}^n k_i x^i\right) \cos(b x) \text{ or }\ \left(\sum_{i=0}^n k_i x^i\right) \sin(b x) \!$                 $\left(\sum_{i=0}^n Q_i x^i\right) \cos(b x) + \left(\sum_{i=0}^n R_i x^i\right) \sin(b x)$

  $\left(\sum_{i=0}^n k_i x^i\right) e^{a x} \cos(b x) \text{ or } \left(\sum_{i=0}^n k_i x^i\right) e^{a x} \sin(b x)\!$   $e^{a x} \left(\left(\sum_{i=0}^n Q_i x^i\right) \cos(b x) + \left(\sum_{i=0}^n R_i x^i\right) \sin(b x)\right)$



If a term in the above particular integral for \textit{y} appears in the

homogeneous solution, it is necessary to multiply by a sufficiently

large power of \textit{x} in order to make the solution independent. If the

function of \textit{x} is a sum of terms in the above table, the particular

integral can be guessed using a sum of the corresponding terms for \textit{y}.



\paragraph{Examples}



\paragraph{Example 1}



Find a particular integral of the equation



$$y'' + y = t \cos t.$$



The right side \textit{t} cos \textit{t} has the form



$$P_n e^{\alpha t} \cos{\beta t}$$



with \textit{n} = 2, \textit{a} = 0, and \textit{ß} = 1.



Since \textit{a} + \textit{iß} = \textit{i} is \textit{a simple root} of the characteristic equation



$$\lambda^2 + 1 = 0$$



we should try a particular integral of the form



$$\begin{aligned}

y_p &= t \left [F_1 (t) e^{\alpha t} \cos{\beta t} + G_1 (t) e^{\alpha t} \sin{\beta t} \right ] \\\&= t \left [F_1 (t) \cos t + G_1 (t) \sin t \right ] \\\&= t \left [ \left (A_0 t + A_1 \right ) \cos t + \left (B_0 t + B_1 \right ) \sin t \right ] \\\&= \left (A_0 t^2 + A_1 t \right ) \cos t + \left (B_0 t^2 + B_1 t \right) \sin t.

\end{aligned}$$



Substituting \textit{y}\textasciitilde{}\textit{p}\textasciitilde{} into the differential equation, we have the

identity



$$\begin{aligned}

t \cos t &= y_p'' + y_p \\\&= \left  [ \left(A_0 t^2 + A_1 t \right ) \cos t + \left (B_0 t^2 + B_1 t \right ) \sin t \right ]'' + \left[\left(A_0 t^2 + A_1 t \right ) \cos t + \left(B_0 t^2 + B_1 t \right ) \sin t \right ] \\\&=  \left [2A_0 \cos t + 2 \left (2A_0 t + A_1 \right )(-\sin t) + \left (A_0 t^2 + A_1 t \right )(-\cos t) + 2B_0 \sin t + 2 \left (2B_0 t + B_1 \right ) \cos t + \left (B_0 t^2 + B_1 t \right )(- \sin t) \right ] \\\&\qquad +\left[\left(A_0 t^2 + A_1 t \right ) \cos t + \left(B_0 t^2 + B_1 t \right ) \sin t \right ] \\\&= [4B_0 t + (2A_0 + 2B_1)] \cos t + [-4A_0 t + (-2A_1 + 2B_0)] \sin t.

\end{aligned}$$



Comparing both sides, we have



$$\begin{cases} 

1 = 4B_0\\\0 = 2A_0 + 2B_1  \\\0 = -4A_0 \\\0 = -2A_1 + 2B_0

\end{cases}$$



which has the solution



$$A_0 = 0, \quad A_1 = B_0 = \frac{1}{4}, \quad  B_1 = 0.$$



We then have a particular integral



$$y_p = \frac {1} {4} t \cos t + \frac {1}{4} t^2 \sin t.$$



\paragraph{Example 2}



Consider the following linear nonhomogeneous differential equation:



:   $\frac{dy}{dx} = y + e^x.$



This is like the first example above, except that the nonhomogeneous

part ($e^x$) is \textit{not} linearly independent to the general solution of

the homogeneous part ($c_1 e^x$); as a result, we have to multiply our

guess by a sufficiently large power of \textit{x} to make it linearly

independent.



Here our guess becomes:



$$y_p = A x e^x.$$



By substituting this function and its derivative into the differential

equation, one can solve for \textit{A}:



$$\frac{d}{dx} \left( A x e^x \right) = A x e^x + e^x$$



$$A x e^x + A e^x = A x e^x + e^x$$



$$A = 1.$$



So, the general solution to this differential equation is:



$$y = c_1 e^x + xe^x.$$



\paragraph{Example 3}



Find the general solution of the equation:



$$\frac{dy}{dt} = t^2 - y$$



$t^2$ is a polynomial of degree 2, so we look for a solution using the

same form,



$$y_p = A t^2 + B t + C,$$



Plugging this particular function into the original equation yields,



$$2 A t + B = t^2 - (A t^2 + B t + C),$$



$$2 A t + B =(1-A)t^2 -Bt -C,$$



$$(A-1)t^2 + (2A+B)t + (B+C) = 0.$$



which gives:



$$A-1 = 0, \quad 2A+B =0, \quad  B+C=0.$$



Solving for constants we get:



$$y_p = t^2 - 2 t + 2$$



To solve for the general solution,



$$y= y_p + y_c$$



where $y_c$ is the homogeneous solution $y_c = c_1 e^{-t}$, therefore,

the general solution is:



$$y= t^2 - 2 t + 2 + c_1 e^{-t}$$



\paragraph{Higher Order}



If we have a second order linear nonhomogenous differential equation

whose coefficients were constant, that is

$a \frac{d^2y}{dt^2} + b \frac{dy}{dt} + cy = g(t)$, then to solve this

differential equation, all we need to do is solve the corresponding

second order linear homogenous differential equation

$a \frac{d^2y}{dt^2} + b \frac{dy}{dt} + cy = 0$, and then find a

partial solution by assuming the form of the particular solution. More

precisely, if $g(t)$ was a polynomial, exponential, or sine/cosine

function (or a combination of these), then we could assume a form for

the particular solutions (see the linked page above for more details)

and solve for the coefficients of this form to obtain a particular

solution.



We will now look at this method for higher order linear nonhomogenous

differential equations. Consider the following $n^{\mathrm{th}}$ order

linear nonhomogenous differential equation with coefficients

$a_0, a_1, ..., a_n \in \mathbb{R}$:



$\begin{aligned} \quad a_0 \frac{d^ny}{dt^n} + a_1 \frac{d^{n-1}y}{dt^{n-1}} + ... + a_{n-1}\frac{dy}{dt} + a_n y = g(t) \end{aligned}$



Suppose that $g(t)$ is of a form containing a polynomial, exponential

function, or a sine/cosine function (like with when we were dealing with

the method of undetermined coefficients for second order linear

nonhomogenous differential equations). We can then find a particular

solution $Y(t)$ if we extend the assumed forms of combinations of

polynomials, exponential functions, or sine/cosine functions, and

multiplying by powers of $t$ to ensure our particular solution does not

contain part of a solution to the corresponding higher order linear

homogenous differential equation. We can then solve for these constants

by plugging our assumed form $Y(t)$ into our differential equation

above.



Providing a list of possible forms is trivial, so we will instead look

at some examples of apply this method.



\paragraph{Example 1}



``{=html}Find the general solution to the differential equation

$\frac{d^3y}{dt^3} + \frac{d^2y}{dt^2} + \frac{dy}{dt} + y = e^{-t} + 4t$.``{=html}



We will need to first solve the corresponding third order linear

homogenous differential equation

$\frac{d^3y}{dt^3} + \frac{d^2y}{dt^2} + \frac{dy}{dt} + y = 0$. This

characteristic equation to this differential equation is:



$\begin{aligned} \quad r^3 + r^2 + r + 1 = 0 \end{aligned}$



We can (by trial and error) see that $r_1 = -1$ is a solution to this

characteristic equation. Applying long division with the factor

$(r + 1)$ and we have that our characteristic equation can be written

as:



$\begin{aligned} \quad (r + 1)(r^2 + 1) = 0 \end{aligned}$



Therefore we can see that $r_2 = i$ and $r_3 = -i$. Therefore the

general solution to our third order linear homogenous differential

equation is:



$\begin{aligned} \quad y_h(t) = C_1e^{-t} + C_2\cos(t) + C_3\sin(t) \end{aligned}$



Now we note that $g(t) = e^{-t} + 4t$ has an exponential term and a

cosine term, so we expect the form of our particular solution to be

$Y(t) = Ae^{-t} ...$. We now compute the first, second, and third

derivatives of $Y$.



\paragraph{Licensing}



Content obtained and/or adapted from:



\begin{itemize}
    \item \href{https://en.wikipedia.org/wiki/Method_of_undetermined_coefficients}{Method of undetermined coefficients,
\end{itemize}
    Wikipedia}

    under a CC BY-SA license"

%%===============================================================
\section{Linear Differential Equations (Higher Order)}
\label{sec:linear-differential-equations-higher-order}
%%===============================================================

"In mathematics, a \textbf{linear differential equation} is a differential

equation that is defined by a linear polynomial in the unknown function

and its derivatives, that is an equation of the form



$$a_0(x)y + a_1(x)y' + a_2(x)y'' \cdots + a_n(x)y^{(n)} = b(x)$$ where

\textit{a}\textasciitilde{}0\textasciitilde{}(\textit{x}), ..., \textit{a}\textasciitilde{}\textit{n}\textasciitilde{}(\textit{x}) and `{{math|''b''(''x'')}}`{=mediawiki}

are arbitrary differentiable functions that do not need to be linear,

and `{{math|''y''', …, ''y''(''n'') }}`{=mediawiki} are the

successive derivatives of an unknown function `{{mvar|y}}`{=mediawiki}

of the variable `{{mvar|x}}`{=mediawiki}.



Such an equation is an ordinary differential equation (ODE). A *linear

differential equation* may also be a linear partial differential

equation (PDE), if the unknown function depends on several variables,

and the derivatives that appear in the equation are partial derivatives.



A linear differential equation or a system of linear equations such that

the associated homogeneous equations have constant coefficients may be

solved by quadrature, which means that the solutions may be expressed in

terms of integrals. This is also true for a linear equation of order

one, with non-constant coefficients. An equation of order two or higher

with non-constant coefficients cannot, in general, be solved by

quadrature. For order two, Kovacic\'s algorithm allows deciding whether

there are solutions in terms of integrals, and computing them if any.



The solutions of linear differential equations with polynomial

coefficients are called holonomic functions. This class of functions is

stable under sums, products, differentiation, integration, and contains

many usual functions and special functions such as exponential function,

logarithm, sine, cosine, inverse trigonometric functions, error

function, Bessel functions and hypergeometric functions. Their

representation by the defining differential equation and initial

conditions allows making algorithmic (on these functions) most

operations of calculus, such as computation of antiderivatives, limits,

asymptotic expansion, and numerical evaluation to any precision, with a

certified error bound.



\paragraph{Basic terminology}



The highest order of derivation that appears in a (linear) differential

equation is the \textit{order} of the equation. The term

`{{math|''b''(''x'')}}`{=mediawiki}, which does not depend on the

unknown function and its derivatives, is sometimes called the *constant

term* of the equation (by analogy with algebraic equations), even when

this term is a non-constant function. If the constant term is the zero

function, then the differential equation is said to be \textit{homogeneous}, as

it is a homogeneous polynomial in the unknown function and its

derivatives. The equation obtained by replacing, in a linear

differential equation, the constant term by the zero function is the

\textit{associated homogeneous equation}. A differential equation has *constant

coefficients* if only constant functions appear as coefficients in the

associated homogeneous equation.



A \textit{solution} of a differential equation is a function that satisfies the

equation. The solutions of a homogeneous linear differential equation

form a vector space. In the ordinary case, this vector space has a

finite dimension, equal to the order of the equation. All solutions of a

linear differential equation are found by adding to a particular

solution any solution of the associated homogeneous equation.



\paragraph{Linear differential operator}



A \textit{basic differential operator} of order `{{mvar|i}}`{=mediawiki} is a

mapping that maps any differentiable function to its

`{{mvar|i}}`{=mediawiki}th derivative, or, in the case of several

variables, to one of its partial derivatives of order

`{{mvar|i}}`{=mediawiki}. It is commonly denoted



$$\frac{d^i}{dx^i}$$ in the case of univariate functions, and



$$\frac{\partial^{i_1+\cdots +i_n}}{\partial x_1^{i_1}\cdots \partial x_n^{i_n}}$$

in the case of functions of `{{mvar|n}}`{=mediawiki} variables. The

basic differential operators include the derivative of order 0, which is

the identity mapping.



A \textbf{linear differential operator} (abbreviated, in this article, as

\textit{linear operator} or, simply, \textit{operator}) is a linear combination of

basic differential operators, with differentiable functions as

coefficients. In the univariate case, a linear operator has thus the

form



$$a_0(x)+a_1(x)\frac{d}{dx} + \cdots +a_n(x)\frac{d^n}{dx^n},$$ where

`{{math|''a''0(''x''), …, ''a''''n''(''x'')}}`{=mediawiki}

are differentiable functions, and the nonnegative integer

`{{mvar|n}}`{=mediawiki} is the \textit{order} of the operator (if

`{{math|''a''''n''(''x'')}}`{=mediawiki} is not the zero

function).



Let `{{mvar|L}}`{=mediawiki} be a linear differential operator. The

application of `{{mvar|L}}`{=mediawiki} to a function

`{{mvar|f}}`{=mediawiki} is usually denoted

`{{math|''Lf''}}`{=mediawiki} or `{{math|''Lf''(''X'')}}`{=mediawiki},

if one needs to specify the variable (this must not be confused with a

multiplication). A linear differential operator is a linear operator,

since it maps sums to sums and the product by a scalar to the product by

the same scalar.



As the sum of two linear operators is a linear operator, as well as the

product (on the left) of a linear operator by a differentiable function,

the linear differential operators form a vector space over the real

numbers or the complex numbers (depending on the nature of the functions

that are considered). They form also a free module over the ring of

differentiable functions.



The language of operators allows a compact writing for differentiable

equations: if



$$L=a_0(x)+a_1(x)\frac{d}{dx} + \cdots +a_n(x)\frac{d^n}{dx^n},$$ is a

linear differential operator, then the equation



$$a_0(x)y +a_1(x)y' + a_2(x)y'' +\cdots +a_n(x)y^{(n)}=b(x)$$ may be

rewritten



$$Ly=b(x).$$



There may be several variants to this notation; in particular the

variable of differentiation may appear explicitly or not in

`{{mvar|y}}`{=mediawiki} and the right-hand and of the equation, such as

`{{math|1=''Ly''(''x'') = ''b''(''x'')}}`{=mediawiki} or

`{{math|1=''Ly'' = ''b''}}`{=mediawiki}.



The \textit{kernel} of a linear differential operator is its kernel as a linear

mapping, that is the vector space of the solutions of the (homogeneous)

differential equation `{{math|1=''Ly'' = 0}}`{=mediawiki}.



In the case of an ordinary differential operator of order

`{{mvar|n}}`{=mediawiki}, Carathéodory\'s existence theorem implies

that, under very mild conditions, the kernel of `{{mvar|L}}`{=mediawiki}

is a vector space of dimension `{{mvar|n}}`{=mediawiki}, and that the

solutions of the equation

`{{math|1=''Ly''(''x'') = ''b''(''x'')}}`{=mediawiki} have the form



$$S_0(x) + c_1S_1(x) + \cdots +c_nS_n(x),$$ where

`{{math|''c''1, …, ''c''''n''}}`{=mediawiki} are

arbitrary numbers. Typically, the hypotheses of Carathéodory\'s theorem

are satisfied in an interval `{{mvar|I}}`{=mediawiki}, if the functions

`{{math|''b'', ''a''0, …, ''a''''n''}}`{=mediawiki}

are continuous in `{{mvar|I}}`{=mediawiki}, and there is a positive real

number `{{mvar|k}}`{=mediawiki} such that \|\textit{a}\textasciitilde{}\textit{n}\textasciitilde{}(\textit{x})\| \> \textit{k} for

every `{{mvar|x}}`{=mediawiki} in `{{mvar|I}}`{=mediawiki}.



\paragraph{Homogeneous equation with constant coefficients}



A homogeneous linear differential equation has \textit{constant coefficients}

if it has the form



$$a_0y + a_1y' + a_2y'' + \cdots + a_n y^{(n)} = 0$$ where

`{{math|''a''1, …, ''a''''n''}}`{=mediawiki} are

(real or complex) numbers. In other words, it has constant coefficients

if it is defined by a linear operator with constant coefficients.



The study of these differential equations with constant coefficients

dates back to Leonhard Euler, who introduced the exponential function

`{{math|''e''''x''}}`{=mediawiki}, which is the unique

solution of the equation `{{math|1=''f''' = ''f''}}`{=mediawiki} such

that `{{math|1=''f''(0) = 1}}`{=mediawiki}. It follows that the

`{{mvar|n}}`{=mediawiki}th derivative of

`{{math|''e''''cx'' }}`{=mediawiki} is

`{{math|''c''''n''''e''''cx''}}`{=mediawiki}, and

this allows solving homogeneous linear differential equations rather

easily.



Let



$$a_0y + a_1y' + a_2y'' + \cdots + a_ny^{(n)} = 0$$ be a homogeneous

linear differential equation with constant coefficients (that is

`{{math|''a''0, …, ''a''''n''}}`{=mediawiki} are

real or complex numbers).



Searching solutions of this equation that have the form

`{{math|''e''''ax''}}`{=mediawiki} is equivalent to searching

the constants `{{mvar|a}}`{=mediawiki} such that



$$a_0e^{\alpha x} + a_1\alpha e^{\alpha x} + a_2\alpha^2 e^{\alpha x}+\cdots + a_n\alpha^n e^{\alpha x} = 0.$$

Factoring out `{{math|''e''''ax''}}`{=mediawiki} (which is

never zero), shows that `{{mvar|a}}`{=mediawiki} must be a root of the

\textit{characteristic polynomial}



$$a_0 + a_1t + a_2t^2 + \cdots + a_nt^n$$ of the differential equation,

which is the left-hand side of the characteristic equation



$$a_0 + a_1t + a_2t^2 + \cdots + a_nt^n = 0.$$



When these roots are all distinct, one has `{{mvar|n}}`{=mediawiki}

distinct solutions that are not necessarily real, even if the

coefficients of the equation are real. These solutions can be shown to

be linearly independent, by considering the Vandermonde determinant of

the values of these solutions at

`{{math|1=''x'' = 0, …, ''n'' – 1}}`{=mediawiki}. Together they form a

basis of the vector space of solutions of the differential equation

(that is, the kernel of the differential operator).



+----------------------------------------------------------------------+

| Example                                                              |

+======================================================================+

| $$y''''-2y'''+2y''-2y'+y=0$$ has the characteristic equation         |

|                                                                      |

| :   $z^4-2z^3+2z^2-2z+1=0.$                                          |

|                                                                      |

| This has zeros, `{{mvar|i}}`{=mediawiki},                            |

| `{{math|-''i''}}`{=mediawiki}, and `{{math|1}}`{=mediawiki}          |

| (multiplicity 2). The solution basis is thus                         |

|                                                                      |

| :   $e^{ix}, \; e^{-ix}, \; e^x, \; xe^x.$                              |

|                                                                      |

| A real basis of solution is thus                                     |

|                                                                      |

| :   $\cos x, \; \sin x, \; e^x, \; xe^x.$                               |

+----------------------------------------------------------------------+



In the case where the characteristic polynomial has only simple roots,

the preceding provides a complete basis of the solutions vector space.

In the case of multiple roots, more linearly independent solutions are

needed for having a basis. These have the form



$$x^ke^{\alpha x},$$ where `{{mvar|k}}`{=mediawiki} is a nonnegative

integer, `{{mvar|a}}`{=mediawiki} is a root of the characteristic

polynomial of multiplicity `{{mvar|m}}`{=mediawiki}, and

`{{math|''k'' < ''m''}}`{=mediawiki}. For proving that these functions

are solutions, one may remark that if `{{mvar|a}}`{=mediawiki} is a root

of the characteristic polynomial of multiplicity

`{{mvar|m}}`{=mediawiki}, the characteristic polynomial may be factored

as

`{{math|''P''(''t'') (''t'' - ''a'')''m''}}`{=mediawiki}.

Thus, applying the differential operator of the equation is equivalent

with applying first `{{mvar|m}}`{=mediawiki} times the operator

$\frac{d}{dx} - \alpha$, and then the operator that has

`{{mvar|P}}`{=mediawiki} as characteristic polynomial. By the

exponential shift theorem,



$$\left(\frac{d}{dx}-\alpha\right)\left(x^ke^{\alpha x}\right)= kx^{k-1}e^{\alpha x},$$



and thus one gets zero after `{{math|''k'' + 1}}`{=mediawiki}

application of $\frac{d}{dx} - \alpha$.



As, by the fundamental theorem of algebra, the sum of the multiplicities

of the roots of a polynomial equals the degree of the polynomial, the

number of above solutions equals the order of the differential equation,

and these solutions form a base of the vector space of the solutions.



In the common case where the coefficients of the equation are real, it

is generally more convenient to have a basis of the solutions consisting

of real-valued functions. Such a basis may be obtained from the

preceding basis by remarking that, if

`{{math|''a'' + ''ib''}}`{=mediawiki} is a root of the characteristic

polynomial, then `{{math|''a'' – ''ib''}}`{=mediawiki} is also a root,

of the same multiplicity. Thus a real basis is obtained by using

Euler\'s formula, and replacing

`{{math|''x''''k''''e''(''a''+''ib'')''x''}}`{=mediawiki}

and

`{{math|''x''''k''''e''(''a''-''ib'')''x''}}`{=mediawiki}

by

`{{math|''x''''k''''e''''ax'' cos(''bx'')}}`{=mediawiki}

and

`{{math|''x''''k''''e''''ax'' sin(''bx'')}}`{=mediawiki}.



\paragraph{Second-order case}



A homogeneous linear differential equation of the second order may be

written



$$y'' + ay' + by = 0,$$ and its characteristic polynomial is



$$r^2 + ar + b.$$



If `{{mvar|a}}`{=mediawiki} and `{{mvar|b}}`{=mediawiki} are real, there

are three cases for the solutions, depending on the discriminant

`{{math|1=''D'' = ''a''2 - 4''b''}}`{=mediawiki}. In all

three cases, the general solution depends on two arbitrary constants

`{{math|''c''1}}`{=mediawiki} and

`{{math|''c''2}}`{=mediawiki}.



\begin{itemize}
    \item If `{{math|''D'' > 0}}`{=mediawiki}, the characteristic polynomial
\end{itemize}
    has two distinct real roots `{{mvar|a}}`{=mediawiki}, and

    `{{mvar|ß}}`{=mediawiki}. In this case, the general solution is



$$c_1 e^{\alpha x} + c_2 e^{\beta x}.$$



\begin{itemize}
    \item If `{{math|1=''D'' = 0}}`{=mediawiki}, the characteristic polynomial
\end{itemize}
    has a double root `{{math|-''a''/2}}`{=mediawiki}, and the general

    solution is



$$(c_1 + c_2 x) e^{-ax/2}.$$



\begin{itemize}
    \item If `{{math|''D'' < 0}}`{=mediawiki}, the characteristic polynomial
\end{itemize}
    has two complex conjugate roots

    `{{math|''a'' ± ''ßi''}}`{=mediawiki}, and the general solution is



$$c_1 e^{(\alpha + \beta i)x} + c_2 e^{(\alpha - \beta i)x},$$



:   which may be rewritten in real terms, using Euler\'s formula as

    $$e^{\alpha x}  (c_1\cos(\beta x) + c_2 \sin(\beta x)).$$



Finding the solution `{{math|''y''(''x'')}}`{=mediawiki} satisfying

`{{math|1=''y''(0) = ''d''1}}`{=mediawiki} and

`{{math|1=''y'''(0) = ''d''2}}`{=mediawiki}, one equates the

values of the above general solution at `{{math|0}}`{=mediawiki} and its

derivative there to `{{math|''d''1}}`{=mediawiki} and

`{{math|''d''2}}`{=mediawiki}, respectively. This results in

a linear system of two linear equations in the two unknowns

`{{math|''c''1}}`{=mediawiki} and

`{{math|''c''2}}`{=mediawiki}. Solving this system gives the

solution for a so-called Cauchy problem, in which the values at

`{{math|0}}`{=mediawiki} for the solution of the DEQ and its derivative

are specified.



\paragraph{Non-homogeneous equation with constant coefficients}



A non-homogeneous equation of order `{{mvar|n}}`{=mediawiki} with

constant coefficients may be written



$$y^{(n)}(x) + a_1 y^{(n-1)}(x) + \cdots + a_{n-1} y'(x)+ a_ny(x) = f(x),$$

where `{{math|''a''1, …, ''a''''n''}}`{=mediawiki}

are real or complex numbers, `{{mvar|f}}`{=mediawiki} is a given

function of `{{mvar|x}}`{=mediawiki}, and `{{mvar|y}}`{=mediawiki} is

the unknown function (for sake of simplicity,

\""`{{math|(''x'')}}`{=mediawiki}\"" will be omitted in the following).



There are several methods for solving such an equation. The best method

depends on the nature of the function `{{mvar|f}}`{=mediawiki} that

makes the equation non-homogeneous. If `{{mvar|f}}`{=mediawiki} is a

linear combination of exponential and sinusoidal functions, then the

exponential response formula may be used. If, more generally,

`{{mvar|f}}`{=mediawiki} is a linear combination of functions of the

form `{{math|''x''''n''''e''''ax''}}`{=mediawiki},

`{{math|''x''''n'' cos(''ax'')}}`{=mediawiki}, and

`{{math|''x''''n'' sin(''ax'')}}`{=mediawiki}, where

`{{mvar|n}}`{=mediawiki} is a nonnegative integer, and

`{{mvar|a}}`{=mediawiki} a constant (which need not be the same in each

term), then the method of undetermined coefficients may be used. Still

more general, the annihilator method applies when

`{{mvar|f}}`{=mediawiki} satisfies a homogeneous linear differential

equation, typically, a holonomic function.



The most general method is the variation of constants, which is

presented here.



The general solution of the associated homogeneous equation



$$y^{(n)} + a_1 y^{(n-1)} + \cdots + a_{n-1} y'+ a_ny = 0$$ is



$$y=u_1y_1+\cdots+ u_ny_n,$$ where

`{{math|(''y''1, …, ''y''''n'')}}`{=mediawiki} is

a basis of the vector space of the solutions and

`{{math|''u''1, …, ''u''''n''}}`{=mediawiki} are

arbitrary constants. The method of variation of constants takes its name

from the following idea. Instead of considering

`{{math|''u''1, …, ''u''''n''}}`{=mediawiki} as

constants, they can considered as unknown functions that have to be

determined for making `{{mvar|y}}`{=mediawiki} a solution of the

non-homogeneous equation. For this purpose, one adds the constraints



$$\begin{aligned}

0 &= u'_1y_1 + u'_2y_2 + \cdots+u'_ny_n \\\0 &= u'_1y'_1 + u'_2y'_2 + \cdots + u'_n y'_n \\\  & \; \;\vdots \\\0 &= u'_1y^{(n-2)}_1+u'_2y^{(n-2)}_2 + \cdots + u'_n y^{(n-2)}_n,

\end{aligned}$$ which imply (by product rule and induction)



$$y^{(i)} = u_1 y_1^{(i)} + \cdots + u_n y_n^{(i)}$$ for

`{{math|1=''i'' = 1, …, ''n'' – 1}}`{=mediawiki}, and



$$y^{(n)} = u_1 y_1^{(n)} + \cdots + u_n y_n^{(n)} +u'_1y_1^{(n-1)}+u'_2y_2^{(n-1)}+\cdots+u'_ny_n^{(n-1)}.$$



Replacing in the original equation `{{mvar|y}}`{=mediawiki} and its

derivatives by these expressions, and using the fact that

`{{math|''y''1, …, ''y''''n''}}`{=mediawiki} are

solutions of the original homogeneous equation, one gets



$$f=u'_1y_1^{(n-1)} + \cdots + u'_ny_n^{(n-1)}.$$



This equation and the above ones with `{{math|0}}`{=mediawiki} as

left-hand side form a system of `{{mvar|n}}`{=mediawiki} linear

equations in

`{{math|''u'''1, …, ''u'''''n''}}`{=mediawiki}

whose coefficients are known functions (`{{mvar|f}}`{=mediawiki}, the

$y_i$, and their derivatives). This system can be solved by any method

of linear algebra. The computation of antiderivatives gives

`{{math|''u''1, …, ''u''''n''}}`{=mediawiki}, and

then

`{{math|1=''y'' = ''u''1''y''1 + ? + ''u''''n''''y''''n''}}`{=mediawiki}.



As antiderivatives are defined up to the addition of a constant, one

finds again that the general solution of the non-homogeneous equation is

the sum of an arbitrary solution and the general solution of the

associated homogeneous equation.



\paragraph{First-order linear equation}



The general form of a linear ordinary differential equation of order 1,

after dividing out the coefficient of

`{{math|''y'''(''x'')}}`{=mediawiki}, is:



$$y'(x) = f(x) y(x) + g(x).$$



If the equation is homogeneous, i.e.

`{{math|1=''g''(''x'') = 0}}`{=mediawiki}, one may rewrite and

integrate:



$$\frac{y'}{y}= f, \qquad \log y = k +F,$$ where

`{{mvar|k}}`{=mediawiki} is an arbitrary constant of integration and

$F=\textstyle\int f\,dx$ is any antiderivative of

`{{mvar|f}}`{=mediawiki}. Thus, the general solution of the homogeneous

equation is



$$y=ce^F,$$ where `{{math|1=''c'' = ''e''''k''}}`{=mediawiki}

is an arbitrary constant.



For the general non-homogeneous equation, one may multiply it by the

reciprocal `{{math|''e''-''F''}}`{=mediawiki} of a solution

of the homogeneous equation. This gives



$$y'e^{-F}-yfe^{-F}= ge^{-F}.$$ As

$-fe^{-F} = \frac{d}{dx} \left(e^{-F}\right),$ the product rule allows

rewriting the equation as



$$\frac{d}{dx}\left(ye^{-F}\right)= ge^{-F}.$$ Thus, the general

solution is



$$y=ce^F + e^F\int ge^{-F}dx,$$ where `{{mvar|c}}`{=mediawiki} is a

constant of integration, and `{{mvar|F}}`{=mediawiki} is any

antiderivative of `{{mvar|f}}`{=mediawiki} (changing of antiderivative

amounts to change the constant of integration).



\paragraph{Example}



Solving the equation



:   $y'(x) + \frac{y(x)}{x} = 3x.$



The associated homogeneous equation $y'(x) + \frac{y(x)}{x} = 0$ gives



$$\frac{y'}{y}=-\frac{1}{x},$$ that is



$$y=\frac{c}{x}.$$



Dividing the original equation by one of these solutions gives



:   $xy'+y=3x^2.$



That is



$$(xy)'=3x^2,$$



$$xy=x^3 +c,$$ and



$$y(x)=x^2+c/x.$$ For the initial condition



:   $y(1)=\alpha,$



one gets the particular solution



$$y(x)=x^2+\frac{\alpha-1}{x}.$$



\paragraph{Making Linear Equations from Non-Linear Equations}



Sometimes a non-linear equation, which is not solvable like this, can be

made linear, and more easily solvable, by applying a substitution.



\paragraph{Example}



$$y\frac{dy}{dx}+xy^2=5x$$



:   Let\'s make the following substitution:



$$v=y^2 \,$$



$$\frac{dv}{dx}=2y\frac{dy}{dx}$$



:   Plugging in, we get



$$\frac{1}{2}\frac{dv}{dx}+xv=5x$$



$$\frac{dv}{dx}+2xv=10x$$



:   We can then solve as a linear equation in \textit{v}, using the

    step-by-step method above:



<!-- -->



:   \textbf{Step 1}: Find the integrating factor:



$$e^{\int P(x)dx}$$



$$e^{\int 2x dx}=e^{x^2+C}$$



$$e^{\int P(x)dx}=Ce^{x^2}$$



:   Letting C=1 for convenience, we get $e^{x^2}$ as our integrating

    factor.



<!-- -->



:   \textbf{Step 2}: Multiply through



$$e^{x^2}\frac{dv}{dx}+e^{x^2}xv=10e^{x^2}x$$



:   \textbf{Step 3}: Recognize that the left hand is

    $\frac{d}{dx} e^{\int P(x)dx}v$



$$\frac{d}{dx} e^{x^2}v=10e^{x^2}x$$



:   \textbf{Step 4}: Integrate both sides w.r.t.\textit{x}.



$$\int (\frac{d}{dx} e^{x^2}v)dx=\int 10e^{x^2}xdx$$



$$e^{x^2}v=5e^{x^2}+C$$



:   \textbf{Step 5}: Solve for \textit{v}.



$$v=5+\frac{C}{e^{x^2}}$$



:   Now that we have \textit{v}, solve for \textit{y}.



$$v=y^2 \,$$



$$y^2=5+\frac{C}{e^{x^2}} \,$$



\paragraph{Higher order linear equations}



A linear ordinary equation of order one with variable coefficients may

be solved by quadrature, which means that the solutions may be expressed

in terms of integrals. This is not the case for order at least two. This

is the main result of Picard--Vessiot theory which was initiated by

Émile Picard and Ernest Vessiot, and whose recent developments are

called differential Galois theory.



The impossibility of solving by quadrature can be compared with the

Abel--Ruffini theorem, which states that an algebraic equation of degree

at least five cannot, in general, be solved by radicals. This analogy

extends to the proof methods and motivates the denomination of

differential Galois theory.



Similarly to the algebraic case, the theory allows deciding which

equations may be solved by quadrature, and if possible solving them.

However, for both theories, the necessary computations are extremely

difficult, even with the most powerful computers.



Nevertheless, the case of order two with rational coefficients has been

completely solved by Kovacic\'s algorithm.



\paragraph{Cauchy--Euler equation}



Cauchy--Euler equations are examples of equations of any order, with

variable coefficients, that can be solved explicitly. These are the

equations of the form



$$x^n y^{(n)}(x) + a_{n-1} x^{n-1} y^{(n-1)}(x) + \cdots + a_0 y(x) = 0,$$

where $a_0, \ldots, a_{n-1}$ are constant coefficients.



\paragraph{Licensing}



Content obtained and/or adapted from:



\begin{itemize}
    \item \href{https://en.wikipedia.org/wiki/Linear_differential_equation}{Linear differential equation,
\end{itemize}
    Wikipedia}

    under a CC BY-SA license

\begin{itemize}
    \item \href{https://en.wikibooks.org/wiki/Ordinary_Differential_Equations/First_Order_Linear_1}{First Order Linear 1, Wikibooks: Differential
\end{itemize}
    Equations}

    under a CC BY-SA license"

%%===============================================================
\section{Power Series Solutions}
\label{sec:power-series-solutions}
%%===============================================================

"In mathematics, the \textbf{power series method} is used to seek a power

series solution to certain differential equations. In general, such a

solution assumes a power series with unknown coefficients, then

substitutes that solution into the differential equation to find a

recurrence relation for the coefficients.



\paragraph{Method}



Consider the second-order linear differential equation



:   $a_2(z)f''(z)+a_1(z)f'(z)+a_0(z)f(z)=0. \; \!$



Suppose \textit{a}\textasciitilde{}2\textasciitilde{} is nonzero for all \textit{z}. Then we can divide throughout to

obtain



:   $f''+{a_1(z)\over a_2(z)}f'+{a_0(z)\over a_2(z)}f=0.$



Suppose further that \textit{a}\textasciitilde{}1\textasciitilde{}/\textit{a}\textasciitilde{}2\textasciitilde{} and \textit{a}\textasciitilde{}0\textasciitilde{}/\textit{a}\textasciitilde{}2\textasciitilde{} are analytic

functions.



The power series method calls for the construction of a power series

solution



$$f=\sum_{k=0}^\infty A_kz^k.$$



If \textit{a}\textasciitilde{}2\textasciitilde{} is zero for some \textit{z}, then the Frobenius method, a variation

on this method, is suited to deal with so called \""singular points\"".

The method works analogously for higher order equations as well as for

systems.



\paragraph{Example usage}



Let us look at the Hermite differential equation,



:   $f''-2zf'+\lambda f=0; \; \lambda=1$



We can try to construct a series solution



:   $f=\sum_{k=0}^\infty A_kz^k$

:   $f'=\sum_{k=1}^\infty kA_kz^{k-1}$

:   $f''=\sum_{k=2}^\infty k(k-1)A_kz^{k-2}$



Substituting these in the differential equation



:   $\begin{aligned}

    \& \sum_{k=2}^\infty k(k-1)A_kz^{k-2}-2z\sum_{k=1}^\infty kA_kz^{k-1}+\sum_{k=0}^\infty A_kz^k=0 \\\    = \& \sum_{k=2}^\infty k(k-1)A_kz^{k-2}-\sum_{k=1}^\infty 2kA_kz^k+\sum_{k=0}^\infty A_kz^k

    \end{aligned}$



Making a shift on the first sum



:   $\begin{aligned}

    \& = \sum_{k=0}^\infty (k+2)(k+1) A_{k+2} z^k - \sum_{k=1}^\infty 2k A_k z^k + \sum_{k=0}^\infty A_k z^k \\\    \& = 2 A_2 + \sum_{k=1}^\infty (k+2)(k+1) A_{k+2} z^k - \sum_{k=1}^\infty 2k A_k z^k + A_0 + \sum_{k=1}^\infty A_k z^k \\\    \& = 2 A_2 + A_0 + \sum_{k=1}^\infty \left( (k+2) (k+1) A_{k+2} + (-2k+1) A_k \right) z^k

    \end{aligned}$



If this series is a solution, then all these coefficients must be zero,

so for both \textit{k}=0 and \textit{k}\>0:



:   $(k+2)(k+1)A_{k+2}+(-2k+1)A_k=0 \; \!$



We can rearrange this to get a recurrence relation for \textit{A}\textasciitilde{}\textit{k}+2\textasciitilde{}.



:   $(k+2)(k+1)A_{k+2}=-(-2k+1)A_k \; \!$



<!-- -->



:   $A_{k+2}={(2k-1)\over (k+2)(k+1)}A_k \; \!$



Now, we have



:   $A_2 = {-1 \over (2)(1)}A_0={-1\over 2}A_0,\, A_3 = {1 \over (3)(2)} A_1={1\over 6}A_1$



We can determine \textit{A}\textasciitilde{}0\textasciitilde{} and \textit{A}\textasciitilde{}1\textasciitilde{} if there are initial conditions, i.e.

if we have an initial value problem.



So we have



:   $\begin{aligned}

    A_4 \& ={1\over 4}A_2 = \left({1\over 4}\right)\left({-1 \over 2}\right)A_0 = {-1 \over 8}A_0 \\\\[8pt]

    A_5 \& ={1\over 4}A_3  = \left({1\over 4}\right)\left({1 \over 6}\right)A_1 = {1 \over 24}A_1 \\\\[8pt]

    A_6 \& = {7\over 30}A_4 = \left({7\over 30}\right)\left({-1 \over 8}\right)A_0 = {-7 \over 240}A_0 \\\\[8pt]

    A_7 \& = {3\over 14}A_5 = \left({3\over 14}\right)\left({1 \over 24}\right)A_1 = {1 \over 112}A_1

    \end{aligned}$



and the series solution is



:   $\begin{aligned}

    f \& = A_0z^0 + A_1z^1 +A_2z^2 +A_3z^3 +A_4z^4 +A_5z^5 + A_6z^6 + A_7z^7+\cdots \\\\[8pt]

    \& = A_0z^0 + A_1z^1 + {-1\over 2}A_0z^2 + {1\over 6}A_1z^3 + {-1 \over 8}A_0z^4 + {1 \over 24}A_1z^5 + {-7 \over 240}A_0z^6 + {1 \over 112}A_1z^7 + \cdots \\\\[8pt]

    \& = A_0z^0 + {-1\over 2}A_0z^2 + {-1 \over 8}A_0z^4 + {-7 \over 240}A_0z^6 + A_1z + {1\over 6}A_1z^3 + {1 \over 24}A_1z^5 + {1 \over 112}A_1z^7 + \cdots

    \end{aligned}$



which we can break up into the sum of two linearly independent series

solutions:



:   $f=A_0 \left(1+{-1\over 2}z^2+{-1 \over 8}z^4+{-7 \over 240}z^6+\cdots\right) + A_1\left(z+{1\over 6}z^3+{1 \over 24}z^5+{1 \over 112}z^7+\cdots\right)$



which can be further simplified by the use of hypergeometric series.



\paragraph{A simpler way using Taylor series}



A much simpler way of solving this equation (and power series solution

in general) using the Taylor series form of the expansion. Here we

assume the answer is of the form



:   $f=\sum_{k=0}^\infty {A_kz^k\over {k!}}$



If we do this, the general rule for obtaining the recurrence

relationship for the coefficients is



$$y^{[n]} \to A_{k+n}$$ and



$$x^m y^{[n]} \to (k)(k-1)\cdots(k-m+1)A_{k+n-m}$$



In this case we can solve the Hermite equation in fewer steps:



:   $f''-2zf'+\lambda f=0; \; \lambda=1$



becomes



:   $A_{k+2} -2kA_k +\lambda A_k=0$



or



:   $A_{k+2} = (2k-\lambda) A_k$



in the series



:   $f=\sum_{k=0}^\infty {A_kz^k\over {k!}}$



\paragraph{Nonlinear equations}



The power series method can be applied to certain nonlinear differential

equations, though with less flexibility. A very large class of nonlinear

equations can be solved analytically by using the Parker--Sochacki

method. Since the Parker--Sochacki method involves an expansion of the

original system of ordinary differential equations through auxiliary

equations, it is not simply referred to as the power series method. The

Parker--Sochacki method is done before the power series method to make

the power series method possible on many nonlinear problems. An ODE

problem can be expanded with the auxiliary variables which make the

power series method trivial for an equivalent, larger system. Expanding

the ODE problem with auxiliary variables produces the same coefficients

(since the power series for a function is unique) at the cost of also

calculating the coefficients of auxiliary equations. Many times, without

using auxiliary variables, there is no known way to get the power series

for the solution to a system, hence the power series method alone is

difficult to apply to most nonlinear equations.



The power series method will give solutions only to initial value

problems (opposed to boundary value problems), this is not an issue when

dealing with linear equations since the solution may turn up multiple

linearly independent solutions which may be combined (by superposition)

to solve boundary value problems as well. A further restriction is that

the series coefficients will be specified by a nonlinear recurrence (the

nonlinearities are inherited from the differential equation).



In order for the solution method to work, as in linear equations, it is

necessary to express every term in the nonlinear equation as a power

series so that all of the terms may be combined into one power series.



As an example, consider the initial value problem



$$F F'' + 2 F'^2 + \eta F' = 0 \quad ; \quad F(1) = 0 \ , \ F'(1) = -\frac{1}{2}$$



which describes a solution to capillary-driven flow in a groove. There

are two nonlinearities: the first and second terms involve products. The

initial values are given at $\eta = 1$, which hints that the power

series must be set up as:



$$F(\eta) = \sum_{i = 0}^{\infty} c_i (\eta - 1)^i$$



since in this way



$$\frac{d^n F}{d \eta^n} \Bigg|_{\eta = 1} = n! \ c_n$$



which makes the initial values very easy to evaluate. It is necessary to

rewrite the equation slightly in light of the definition of the power

series,



$$F F'' + 2 F'^2 + (\eta - 1) F' + F' = 0 \quad ; \quad F(1) = 0 \ , \ F'(1) = -\frac{1}{2}$$



so that the third term contains the same form $\eta - 1$ that shows in

the power series.



The last consideration is what to do with the products; substituting the

power series in would result in products of power series when it\'s

necessary that each term be its own power series. This is where the

\href{Cauchy_product ""wikilink""}{Cauchy product}



$$\left(\sum_{i = 0}^{\infty} a_i x^i\right) \left(\sum_{i = 0}^{\infty} b_i x^i\right) = 

\sum_{i = 0}^{\infty} x^i \sum_{j = 0}^i a_{i - j} b_j$$



is useful; substituting the power series into the differential equation

and applying this identity leads to an equation where every term is a

power series. After much rearrangement, the recurrence



$$\sum_{j = 0}^i \left((j + 1) (j + 2) c_{i - j} c_{j + 2} + 2 (i - j + 1) (j + 1) c_{i - j + 1} c_{j + 1}\right) + i c_i + (i + 1) c_{i + 1} = 0$$



is obtained, specifying exact values of the series coefficients. From

the initial values, $c_0 = 0$ and $c_1 = -1/2$, thereafter the above

recurrence is used. For example, the next few coefficients:



$$c_2 = -\frac{1}{6} \quad ; \quad c_3 = -\frac{1}{108} \quad ; \quad c_4 = \frac{7}{3240} \quad ; \quad c_5 = -\frac{19}{48600} \ \dots$$



A limitation of the power series solution shows itself in this example.

A numeric solution of the problem shows that the function is smooth and

always decreasing to the left of $\eta = 1$, and zero to the right. At

$\eta = 1$, a slope discontinuity exists, a feature which the power

series is incapable of rendering, for this reason the series solution

continues decreasing to the right of $\eta = 1$ instead of suddenly

becoming zero.



\paragraph{Licensing}



Content obtained and/or adapted from:



\begin{itemize}
    \item \href{https://en.wikipedia.org/wiki/Power_series_solution_of_differential_equations}{Power series solution of differential equations,
\end{itemize}
    Wikipedia}

    under a CC BY-SA license"

%%===============================================================
\section{Laplace Transform}
\label{sec:laplace-transform}
%%===============================================================

"In mathematics, the Laplace transform, named after its inventor

Pierre-Simon Laplace, is an integral transform that converts a function

of a real variable t (often time) to a function of a complex variable s

(complex frequency). The transform has many applications in science and

engineering because it is a tool for solving differential equations. In

particular, it transforms linear differential equations into algebraic

equations and convolution into multiplication.



For suitable functions f, the Laplace transform is the integral



$\mathcal{L}\{f\}(s) = \int_0^\infty f(t)e^{-st} \, dt.$



  $f(t)$              $F(s) = \mathcal{L}\{f\}(s)$

  ------------------- -----------------------------------------------------

  $1$                 $\frac{1}{s}$           $s > 0$

  $t$                 $\frac{1}{s^2}$           $s > 0$

  $t^{n}$             $\frac{n!}{s^{n+1}}$           $s > 0$

  $t^{n}e^{at}$       $\frac{n!}{s^{n+1}}$           $s > a$

  $e^{at}$            $\frac{1}{s-a}$           $s > a$

  $\sin{at}$          $\frac{a}{s^2+a^2}$           $s > 0$

  $\cos{at}$          $\frac{s}{s^2+a^2}$           $s > 0$

  $\sinh{at}$         $\frac{a}{s^2-a^2}$           $s > |a|$

  $\cosh{at}$         $\frac{s}{s^2-a^2}$           $s > |a|$

  $e^{at}\sin{bt}$    $\frac{b}{(s-a)^2+b^2}$           $s > a$

  $e^{at}\cos{bt}$    $\frac{s-a}{(s-a)^2+b^2}$           $s > a$

  $e^{at}\sinh{bt}$   $\frac{b}{(s-a)^2-b^2}$           $(s - a) > |b|$

  $e^{at}\cosh{bt}$   $\frac{s-a}{(s-a)^2-b^2}$           $(s - a) > |b|$



  : List of Common Laplace Transforms



\paragraph{Resources}



\begin{itemize}
    \item \href{https://tutorial.math.lamar.edu/classes/de/LaplaceIntro.aspx}{Laplace
\end{itemize}
    Transforms},

    Paul\'s Online Notes



\paragraph{Licensing}



Content obtained and/or adapted from:



\begin{itemize}
    \item \href{https://en.wikipedia.org/wiki/Laplace_transform}{Laplace transform,
\end{itemize}
    Wikipedia} under a

    CC BY-SA license"

%%===============================================================
\section{Inverse Laplace Transform}
\label{sec:inverse-laplace-transform}
%%===============================================================

"In mathematics, the \textbf{inverse Laplace transform} of a function \textit{F}(\textit{s})

is the piecewise-continuous and exponentially-restricted real function

\textit{f}(\textit{t}) which has the property:



$$\mathcal{L}\{f\}(s) = \mathcal{L}\{f(t)\}(s) = F(s),$$ where

$\mathcal{L}$ denotes the Laplace transform.



It can be proven that, if a function \textit{F}(\textit{s}) has the inverse Laplace

transform \textit{f}(\textit{t}), then \textit{f}(\textit{t}) is uniquely determined (considering

functions which differ from each other only on a point set having

Lebesgue measure zero as the same). This result was first proven by

Mathias Lerch in 1903 and is known as Lerch\'s theorem.



The Laplace transform and the inverse Laplace transform together have a

number of properties that make them useful for analysing linear

dynamical systems.



\paragraph{Mellin\'s inverse formula}



An integral formula for the inverse Laplace transform, called the

\textit{Mellin\'s inverse formula}, the \textit{Bromwich integral}, or the

\textit{Fourier-Mellin integral}, is given by the line integral:



$$f(t) = \mathcal{L}^{-1} \{F(s)\}(t) =  \frac{1}{2\pi i}\lim_{T\to\infty}\int_{\gamma-iT}^{\gamma+iT}e^{st}F(s)\,ds$$

where the integration is done along the vertical line Re(\textit{s}) = \textit{?} in

the complex plane such that \textit{?} is greater than the real part of all

singularities of \textit{F}(\textit{s}) and \textit{F}(\textit{s}) is bounded on the line, for

example if contour path is in the region of convergence. If all

singularities are in the left half-plane, or \textit{F}(\textit{s}) is an entire

function , then \textit{?} can be set to zero and the above inverse integral

formula becomes identical to the inverse Fourier transform.



In practice, computing the complex integral can be done by using the

Cauchy residue theorem.



\paragraph{Post\'s inversion formula}



\textbf{Post\'s inversion formula} for Laplace transforms, named after Emil

Post, is a simple-looking but usually impractical formula for evaluating

an inverse Laplace transform.



The statement of the formula is as follows: Let \textit{f}(\textit{t}) be a continuous

function on the interval \\href{http://reference.wolfram.com/mathematica/ref/InverseLaplaceTransform.html}{0, 8) of exponential order, i.e.



:   $\sup_{t>0} \frac{f(t)}{e^{bt}} < \infty$



for some real number \textit{b}. Then for all \textit{s} \> \textit{b}, the Laplace transform

for \textit{f}(\textit{t}) exists and is infinitely differentiable with respect to

\textit{s}. Furthermore, if \textit{F}(\textit{s}) is the Laplace transform of \textit{f}(\textit{t}), then

the inverse Laplace transform of \textit{F}(\textit{s}) is given by



:   $f(t) = \mathcal{L}^{-1} \{F(s)\}(t)

     = \lim_{k \to \infty} \frac{(-1)^k}{k!} \left( \frac{k}{t} \right) ^{k+1} F^{(k)} \left( \frac{k}{t} \right)$



for \textit{t} \> 0, where \textit{F}^(\textit{k})^ is the \textit{k}-th derivative of \textit{F} with

respect to \textit{s}.



As can be seen from the formula, the need to evaluate derivatives of

arbitrarily high orders renders this formula impractical for most

purposes.



With the advent of powerful personal computers, the main efforts to use

this formula have come from dealing with approximations or asymptotic

analysis of the Inverse Laplace transform, using the Grunwald--Letnikov

differintegral to evaluate the derivatives.



Post\'s inversion has attracted interest due to the improvement in

computational science and the fact that it is not necessary to know

where the poles of \textit{F}(\textit{s}) lie, which make it possible to calculate the

asymptotic behaviour for big \textit{x} using inverse Mellin transforms for

several arithmetical functions related to the Riemann hypothesis.



\paragraph{Software tools}



\begin{itemize}
    \item [InverseLaplaceTransform}
\end{itemize}
    performs symbolic inverse transforms in Mathematica

\begin{itemize}
    \item \href{http://library.wolfram.com/infocenter/MathSource/5026/}{Numerical Inversion of Laplace Transform with Multiple Precision
\end{itemize}
    Using the Complex

    Domain} in

    Mathematica gives numerical solutions

\begin{itemize}
    \item \href{http://www.mathworks.co.uk/help/symbolic/ilaplace.html}{ilaplace}
\end{itemize}
    performs symbolic inverse transforms in MATLAB

\begin{itemize}
    \item \href{http://www.mathworks.co.uk/matlabcentral/fileexchange/32824-numerical-inversion-of-laplace-transforms-in-matlab}{Numerical Inversion of Laplace Transforms in
\end{itemize}
    Matlab}

\begin{itemize}
    \item \href{https://www.mathworks.com/matlabcentral/fileexchange/71511-a-cme-based-numerical-inverse-laplace-transformation-method}{Numerical Inversion of Laplace Transforms based on concentrated
\end{itemize}
    matrix-exponential

    functions}

    in Matlab



\paragraph{Licensing}



Content obtained and/or adapted from:



\begin{itemize}
    \item \href{https://en.wikipedia.org/wiki/Inverse_Laplace_transform}{Inverse Laplace transform,
\end{itemize}
    Wikipedia}

    under a CC BY-SA license"

%%===============================================================
\section{L-Transform to ODEs}
\label{sec:l-transform-to-odes}
%%===============================================================

"The Laplace transform is a powerful integral transform used to switch a

function from the time domain to the s-domain. The Laplace transform can

be used in some cases to solve linear differential equations with given

initial conditions.



First consider the following property of the Laplace transform:



$$\mathcal{L}\{f'\}=s\mathcal{L}\{f\}-f(0)$$



$$\mathcal{L}\{f''\}=s^2\mathcal{L}\{f\}-sf(0)-f'(0)$$



One can prove by induction that



$$\mathcal{L}\{f^{(n)}\}=s^n\mathcal{L}\{f\}-\sum_{i=1}^{n}s^{n-i}f^{(i-1)}(0)$$



Now we consider the following differential equation:



$$\sum_{i=0}^{n}a_if^{(i)}(t)=\phi(t)$$



with given initial conditions



$$f^{(i)}(0)=c_i$$



Using the linearity of the Laplace transform it is equivalent to rewrite

the equation as



$$\sum_{i=0}^{n}a_i\mathcal{L}\{f^{(i)}(t)\}=\mathcal{L}\{\phi(t)\}$$



obtaining



$$\mathcal{L}\{f(t)\}\sum_{i=0}^{n}a_is^i-\sum_{i=1}^{n}\sum_{j=1}^{i}a_is^{i-j}f^{(j-1)}(0)=\mathcal{L}\{\phi(t)\}$$



Solving the equation for $\mathcal{L}\{f(t)\}$ and substituting

$f^{(i)}(0)$ with $c_i$ one obtains



$$\mathcal{L}\{f(t)\}=\frac{\mathcal{L}\{\phi(t)\}+\sum_{i=1}^{n}\sum_{j=1}^{i}a_is^{i-j}c_{j-1}}{\sum_{i=0}^{n}a_is^i}$$



The solution for \textit{f}(\textit{t}) is obtained by applying the inverse Laplace

transform to $\mathcal{L}\{f(t)\}.$



Note that if the initial conditions are all zero, i.e.



$$f^{(i)}(0)=c_i=0\quad\forall i\in\{0,1,2,...\ n\}$$



then the formula simplifies to



$$f(t)=\mathcal{L}^{-1}\left\{{\mathcal{L}\{\phi(t)\}\over\sum_{i=0}^{n}a_is^i}\right\}$$



\paragraph{An example}



We want to solve



$$f''(t)+4f(t)=\sin(2t)$$



with initial conditions \textit{f}(0) = 0 and \textit{f'}(0)=0.



We note that



$$\phi(t)=\sin(2t)$$



and we get



$$\mathcal{L}\{\phi(t)\}=\frac{2}{s^2+4}$$



The equation is then equivalent to



$$s^2\mathcal{L}\{f(t)\}-sf(0)-f'(0)+4\mathcal{L}\{f(t)\}=\mathcal{L}\{\phi(t)\}$$



We deduce



$$\mathcal{L}\{f(t)\}=\frac{2}{(s^2+4)^2}$$



Now we apply the Laplace inverse transform to get



$$f(t)=\frac{1}{8}\sin(2t)-\frac{t}{4}\cos(2t)$$



\paragraph{Laplace Transform to Systems of ODEs}



View an example from \href{https://tutorial.math.lamar.edu/classes/de/SystemsLaplace.aspx}{Laplace Transforms, Paul\'s Online

Notes}



\paragraph{Licensing}



Content obtained and/or adapted from:



\begin{itemize}
    \item \href{https://en.wikipedia.org/wiki/Laplace_transform_applied_to_differential_equations}{Laplace transform applied to ODEs,
\end{itemize}
    Wikipedia}

    under a CC BY-SA license"

%%===============================================================
\section{L-Transform to Systems of ODEs}
\label{sec:l-transform-to-systems-of-odes}
%%===============================================================

L-Transform to Systems of ODEs
